\chapter{Obstetrics and Gynecology}

\section{Obstetrics \& Maternal-Fetal Medicine}
\medterm{Abdominal Circumference (AC)}-- The circumference of the fetal abdomen measured at the level of the stomach, umbilical vein (portal sinus), and spine (ribs not visible). The best single parameter for estimating fetal weight. AC is also the most sensitive measurement for detecting IUGR (the liver is the first organ affected by placental insufficiency). AC growth: ~1-1.5 cm per week. A declining AC percentile (or growth arrest) strongly suggests placental insufficiency.

\medterm{Abdominal Pregnancy}-- A very rare ectopic pregnancy (1\%) implanted in the peritoneal cavity outside the fallopian tubes, ovaries, and uterus. Risk factors: prior ectopic, tubal damage, ART. Diagnosis: ultrasound (empty uterus, gestational sac outside the uterus, fetal parts adjacent to maternal organs). High maternal mortality (hemorrhage from placental separation). Treatment: surgical removal. The placenta may be left in situ (risks of infection, hemorrhage) if adherence to vital organs prevents safe removal.

\medterm{Abortion}-- The termination of a pregnancy by medical or surgical means before the fetus reaches viability (typically before 24 weeks gestation in the UK). Medical abortion uses mifepristone (antiprogestogen) followed by misoprostol (prostaglandin analog) to induce uterine contractions and expel pregnancy tissue. Surgical abortion involves vacuum aspiration (suction evacuation) or dilatation and evacuation (D\&E). In the UK, abortion is regulated under the Abortion Act 1967, requiring two doctors to certify that the grounds for abortion are met. Complications are rare: infection, hemorrhage, uterine perforation, and incomplete evacuation. Psychological outcomes are generally positive with appropriate support.

\medterm{Abruptio Placentae}-- Premature separation of the placenta from the uterine wall before delivery, occurring in approximately 0.5-1\% of pregnancies. Classified as revealed (vaginal bleeding present, blood tracks downward) or concealed (bleeding remains behind the placenta, no visible bleeding, higher risk of coagulopathy). Risk factors include hypertension, preeclampsia, maternal smoking, cocaine use, abdominal trauma, thrombophilia, prior abruption, and advanced maternal age. Clinical features include vaginal bleeding (may be absent in concealed abruption), abdominal pain, uterine tenderness and hypertonicity (board-like abdomen), and fetal distress or demise. Complications include maternal hemorrhage, disseminated intravascular coagulation (DIC) from release of thromboplastin into the maternal circulation, fetal hypoxia, preterm delivery, and perinatal death. Diagnosis is primarily clinical; ultrasound has limited sensitivity for detecting retroplacental clot. Management depends on gestational age and severity: immediate delivery (usually by caesarean) is indicated for maternal or fetal compromise irrespective of gestational age, with aggressive blood product replacement and coagulopathy management.

\medterm{Active Management of Third Stage}-- A bundle of interventions to reduce the risk of postpartum hemorrhage: (1) administration of oxytocin (10 IU IM/IV) immediately after delivery of the baby, (2) controlled cord traction (Brandt-Andrews maneuver) to deliver the placenta, and (3) uterine massage after placental delivery. Reduces the incidence of PPH by 50-60\%, reduces mean blood loss, and shortens the third stage. The WHO recommends active management for all vaginal deliveries. Misoprostol (600-800 mcg sublingual/rectal) is an alternative in settings where oxytocin is unavailable.

\medterm{Active Phase of Labor}-- The part of the first stage of labor from 6 cm to full dilatation (10 cm). Characterized by more rapid cervical change (minimum: 1 cm/hour in nulliparas, 1.2 cm/hour in multiparas). The cervix dilates more rapidly as labor progresses. Arrest of active phase: no cervical change for \(\ge\)4 hours with adequate contractions (\(\ge\)200 Montevideo units for 2 hours) or \(\ge\)6 hours with inadequate contractions. Management: oxytocin augmentation for inadequate contractions, cesarean delivery for failed progress despite adequate contractions.

\medterm{Acute Fatty Liver of Pregnancy (AFLP)}-- A rare (1:7000-1:15,000) but life-threatening pregnancy complication characterized by microvesicular fatty infiltration of the liver leading to hepatic failure. Occurs in the third trimester (most commonly 35-37 weeks). Risk factors: primiparity, multiple gestation, male fetus, and disorders of fatty acid oxidation (LCHAD deficiency in the fetus, transmitted to the mother). Presents: nausea, vomiting, abdominal pain, jaundice, and encephalopathy. Lab findings: markedly elevated LFTs (though often mild relative to disease severity), elevated bilirubin, hypoglycemia, elevated ammonia, leukocytosis, DIC (prolonged PT/aPTT, low fibrinogen), and renal impairment. The Swansea criteria are used for diagnosis: \(\ge\)6 features in the setting of no other cause. Treatment: immediate delivery (regardless of gestational age), aggressive supportive care (glucose, coagulation factors, dialysis, and ICU), and liver transplantation if not improving within days. Maternal mortality: 10-20\%. Fetal mortality: 10-20\%.

\medterm{Adhesions}-- Bands of fibrous scar tissue that form between organs or tissues, most commonly following abdominal or pelvic surgery. In gynecology, adhesions can develop after cesarean section, myomectomy, ovarian cystectomy, or treatment of endometriosis. Adhesions can cause pelvic pain, bowel obstruction, and female infertility by distorting tubo-ovarian anatomy and impairing oocyte pickup. Diagnosis is made at laparoscopy (the gold standard). Prevention: gentle tissue handling, meticulous hemostasis, minimizing operative time, and use of adhesion barriers (e.g., Interceed, Seprafilm). Treatment: laparoscopic adhesiolysis for symptomatic adhesions, though adhesions can reform.

\medterm{Advanced Maternal Age}-- Maternal age \(\ge\)35 years at the time of delivery. Associated with increased risks: aneuploidy (trisomy 21: 1:350 at age 35, 1:100 at age 40, 1:25 at age 45), miscarriage, gestational diabetes, preeclampsia, preterm birth, IUGR, placenta previa, and stillbirth. Management: offers of genetic counseling and prenatal diagnosis (NIPT, CVS, amniocentesis). Older mothers also have increased medical comorbidities that require management.

\medterm{Amenorrhea}-- Absence of menstruation. Primary amenorrhea: no menarche by age 15 with normal secondary sexual development, or by age 13 without. Secondary amenorrhea: absence of menses for \(\ge\)3 cycles or \(\ge\)6 months in women with prior menstruation. Causes: pregnancy (most common), hypothalamic suppression (stress, exercise, weight loss), PCOS, hyperprolactinemia, premature ovarian insufficiency, thyroid disorders, outflow obstruction. Evaluation: pregnancy test, FSH, LH, prolactin, TSH, and imaging.

\medterm{Amniocentesis}-- See medical_tests.

\medterm{Amnion}-- The innermost fetal membrane, lining the amniotic cavity and containing the amniotic fluid. A thin, avascular membrane composed of a single layer of cuboidal epithelium overlying a basement membrane and connective tissue. The amnion produces and resorbs amniotic fluid. It fuses with the chorion by 12-14 weeks. Amniotic membranes rupture spontaneously during labor (spontaneous ROM) or are artificially ruptured (AROM). Amniotic bands are rare complications of amnion rupture.

\medterm{Amniotic Fluid}-- The fluid surrounding the fetus in the amniotic cavity, providing cushioning, temperature regulation, and allowing fetal movement, breathing, and swallowing. Composed of water, electrolytes, and fetal components (urine, cells); volume increases with gestation (peak ~800-1000 mL at 34-36 weeks). Abnormalities: oligohydramnios (low—PROM, fetal renal anomaly, IUGR, postterm) and polyhydramnios (high—GI obstruction, neural tube defects, maternal diabetes, twin transfusion). Assessed by ultrasound (amniotic fluid index, deepest vertical pocket). Amniocentesis samples fluid for karyotype/genetic testing and lung maturity.

\medterm{Amniotic Fluid Embolism}-- A rare, catastrophic obstetric emergency where amniotic fluid enters the maternal circulation, causing acute cardiovascular collapse. Characterized by sudden respiratory distress, hypotension, coagulopathy, and seizure. Mortality rate 20-60\% among survivors; 50\% of affected mothers die. Treatment is supportive: cardiopulmonary resuscitation, mechanical ventilation, and management of disseminated intravascular coagulation. High association with mortality and severe maternal morbidity.

\medterm{Amniotic Fluid Index}-- A semi-quantitative assessment of amniotic fluid volume using ultrasound. The uterus divided into four quadrants, and the single deepest vertical pocket measured in each quadrant, then summed. Normal range is 5-24 cm at 16-40 weeks. Less than 5 cm indicates oligohydramnios; greater than 24 cm indicates polyhydramnios. Used in conjunction with fetal growth assessment and Doppler studies to evaluate fetal well-being.

\medterm{Anatomy Survey (Level II Ultrasound)}-- A comprehensive fetal ultrasound performed at 18-22 weeks to evaluate fetal anatomy and screen for structural anomalies. Standard views: fetal head (BPD, HC, lateral ventricles, choroid plexus, cerebellum, cisterna magna), face (lips, orbits, nasal bone), chest (heart rate, rhythm, four-chamber view, outflow tracts, lungs), abdomen (stomach, AC, kidneys, bladder, bowel, anterior abdominal wall), spine (longitudinal and transverse), extremities (femur, humerus, hands, feet), and placenta (location, grade, cord insertion) and amniotic fluid (AFI or single deepest pocket). Sensitivity for major anomalies: 60-80\% depending on BMI, gestational age, and operator experience.

\medterm{Antenatal}-- The period before birth, also called prenatal. The medical supervision and screening provided to pregnant women to monitor the health of the mother and fetus, detect complications early, and provide health education. Antenatal appointments typically follow a schedule: booking appointment (8-12 weeks), anomaly scan (18-21 weeks), and regular follow-ups. Key components: blood pressure monitoring, urine testing, fundal height measurement, fetal heart rate auscultation, and screening for gestational diabetes, pre-eclampsia, and fetal anomalies.

\medterm{Antepartum Fetal Surveillance}-- The monitoring of fetal well-being in the third trimester for high-risk pregnancies. Methods: fetal kick counts (daily maternal counting of fetal movements), non-stress test (NST), biophysical profile (BPP, AFI, fetal tone, movement, breathing), and Doppler (umbilical artery, MCA). Indications: postterm pregnancy, IUGR, preeclampsia, diabetes (pregestational), decreased fetal movement, maternal medical conditions (SLE, antiphospholipid syndrome), and prior stillbirth. Timing: weekly or twice-weekly starting at 32-34 weeks (earlier for severe conditions).

\medterm{Antihypertensives in Pregnancy}-- Medications used to treat hypertension in pregnancy. First-line agents: labetalol (combined alpha/beta blocker, 100-1200 mg BID-TID PO, or IV push for acute hypertension), nifedipine (calcium channel blocker, immediate-release 10-20 mg PO q30min for urgent, or XL 30-60 mg daily for maintenance), and methyldopa (central alpha agonist, 250-1000 mg BID, historically safest, first-line for chronic hypertension). Second-line: hydralazine (direct vasodilator, 5-10 mg IV q20min for acute hypertension). Contraindicated: ACE inhibitors and ARBs (fetal renal agenesis, oligohydramnios, fetal death), atenolol (IUGR), and nitroprusside (cyanide toxicity).

\medterm{APGAR Score}-- A scoring system developed by Virginia Apgar in 1952 to assess the newborn's immediate condition at 1 and 5 minutes after birth. Five components (each scored 0, 1, or 2): Appearance (skin color), Pulse (heart rate), Grimace (reflex irritability), Activity (muscle tone), and Respiration (breathing effort). Score 7-10: normal; 4-6: moderately depressed (may require stimulation or oxygen); 0-3: severely depressed (requires full resuscitation). The 5-minute score is more predictive of neonatal outcome. A low APGAR score does not equate to birth asphyxia.

\medterm{Artificial Rupture of Membranes (AROM)}-- Also called amniotomy. The intentional rupture of the amniotic sac with a small hook (Amnihook) during a cervical exam. Indications: induction or augmentation of labor, placement of internal monitors (fetal scalp electrode, intrauterine pressure catheter), and assessment of amniotic fluid (meconium staining). Risks: cord prolapse (especially with high presenting part), chorioamnionitis (with prolonged ROM), fetal injury, umbilical cord compression, and variable decelerations. AROM alone may augment labor by releasing prostaglandins.

\medterm{Assisted Breech Delivery}-- Vaginal breech delivery with operator assistance after spontaneous delivery to the umbilicus. Maneuvers: delivery of the legs (pinch them at the knees and deliver one at a time), delivery of the arms (Løvset maneuver: rotate the body to deliver the posterior arm, then the anterior arm), and delivery of the aftercoming head (Mauriceau-Smellie-Veit maneuver: the operator's hand supports the head flexed while traction is applied to the shoulders; or Piper forceps). The head must be delivered flexed to enter the pelvis with its smallest diameter. Episiotomy is generally performed.

\medterm{Assisted Reproductive Technology}-- ART: fertility treatments involving manipulation of oocytes/sperm/embryos: in vitro fertilization (IVF), intracytoplasmic sperm injection (ICSI), intrauterine insemination (IUI, some definitions), embryo cryopreservation, gamete donation, and surrogacy. Indications: tubal disease, male factor, unexplained infertility, endometriosis, genetic screening (PGT), and same-sex/single parenthood. IVF: ovarian stimulation (gonadotropins + GnRH analogues), oocyte retrieval, fertilization, embryo culture, transfer, luteal support. Complications: multiple pregnancy, ovarian hyperstimulation syndrome (OHSS), and procedural risks. Success depends on age, cause, and center; significant cost; ethical/legal considerations (embryos, gametes).

\medterm{Asymptomatic Bacteriuria in Pregnancy}-- The presence of >100,000 CFU/mL of bacteria on a clean-catch urine culture without urinary symptoms. Occurs in 2-10\% of pregnancies. Risk factors: prior UTI, diabetes, sickle cell trait, and multiparity. If untreated, 20-40\% develop pyelonephritis. Screening: urine culture at the first prenatal visit. Treatment: nitrofurantoin (100 mg BID x 5-7 days, avoid near term), cephalexin (500 mg QID x 5-7 days), or amoxicillin (500 mg TID x 5-7 days). Recurrence is common; monthly surveillance cultures and suppression may be needed.

\medterm{Augmentation of Labor}-- The stimulation of spontaneous contractions that are judged to be inadequate to achieve progressive cervical dilatation. Most commonly uses oxytocin (Pitocin) infusion. Also includes artificial rupture of membranes (AROM, amniotomy). Indications: arrest of active phase (no cervical change for \(\ge\)4 hours with adequate MVUs), or prolonged latent phase. Pitocin is started at 1-2 mU/min and titrated every 15-30 minutes (up to 40 mU/min) until achieving adequate uterine activity (200-250 Montevideo units). Hyperstimulation (>5 contractions in 10 minutes or contractions lasting >2 minutes) is managed by reducing or stopping oxytocin.

\medterm{Bacterial Vaginosis}-- BV: replacement of normal vaginal lactobacilli by an overgrowth of mixed anaerobes (Gardnerella vaginalis, Mobiluncus, etc.), the most common cause of vaginal discharge in reproductive-age women. Risk: new/multiple partners, douching, smoking. Features: thin, gray-white, fishy-smelling (amine) discharge, pH >4.5, no significant inflammation; many asymptomatic. Diagnosis: clinical, clue cells on wet mount, positive whiff (KOH) test, and Amsel criteria or Nugent score. Complications: increased risk of STI acquisition, PID, and preterm birth. Management: metronidazole (oral or intravaginal) or clindamycin; recurrence is common (maintenance, probiotics); screen/treat in pregnancy. Not sexually transmitted per se but partner treatment not required.

\medterm{Bartholin Cyst}-- A cyst of the Bartholin gland (paired mucus-secreting glands at the posterolateral introitus) due to obstruction of the duct; if infected, a Bartholin abscess. Features: painless (cyst) or painful, tender (abscess) swelling at the posterior labium minus; may cause dyspareunia or labial discomfort; abscess: fluctuant, erythematous, painful. Diagnosis: clinical. Management: small asymptomatic cysts—observation; symptomatic/abscess—incision and drainage with marsupialization or Word catheter insertion (abscess), antibiotics for cellulitis; recurrence may require excision; biopsy of cyst in older women (exclude malignancy). Refer to gynecology.

\medterm{Betamethasone}-- A corticosteroid administered to pregnant women at risk of preterm delivery to accelerate fetal lung maturation. Indications: all pregnancies between 24-36 weeks 6 days at risk for preterm delivery within 7 days. Also for late preterm (34-36 weeks) if not previously given and delivery anticipated in \(\le\)7 days. Dose: two doses of 12 mg IM, 24 hours apart. Mechanism: stimulates surfactant production in type II pneumocytes, induces structural maturation of the lungs (thinner alveolar walls), and reduces risk of RDS (by 40-50\%), IVH, and NEC. Optimal benefit: 24 hours to 7 days after the first dose. Can be repeated (rescue course) if delivery is still imminent >14 days after the initial course.

\medterm{Bile Acids}-- Substances produced in the liver that aid digestion of fats and absorption of fat-soluble vitamins. In pregnancy, elevated bile acids (\(\ge\)10 \(\mu\)mol/L) indicate intrahepatic cholestasis of pregnancy (obstetric cholestasis), a liver disorder characterized by pruritus (itching, typically palms and soles) and elevated serum bile acids. ICP is associated with increased risks of preterm birth, meconium-stained amniotic fluid, and stillbirth. Treatment: ursodeoxycholic acid reduces pruritus and improves bile acid levels. Delivery is typically planned at 37-39 weeks to reduce stillbirth risk.

\medterm{Biophysical Profile}-- A non-stress test combined with ultrasound assessment of five fetal parameters, each scored 0 or 2 for a total of 0-10. Parameters include non-stress test (fetal heart rate reactivity), fetal breathing movements, fetal tone, fetal movement, and amniotic fluid volume. Score of 8-10 is normal; 6 is equivocal; 4 or less is abnormal and may warrant delivery. Used to assess fetal well-being in high-risk pregnancies.

\medterm{Biparietal Diameter (BPD)}-- The transverse diameter of the fetal skull measured at the level of the thalami and cavum septum pellucidum. Used for dating (14-26 weeks, ±7-10 days) and assessing growth. Normal growth: ~2-4 mm per week. BPD is less reliable after 26 weeks and in abnormal fetal head shapes (dolichocephaly, brachycephaly). The cephalic index (BPD/OFD ratio, normal 0.74-0.84) helps identify head shape abnormalities.

\medterm{Birth Control}-- Contraception; see Contraception. Methods to prevent unintended pregnancy: hormonal contraceptives (pill, patch, ring, injection, implant), intrauterine devices (copper, levonorgestrel), barrier methods (condoms, diaphragm), emergency contraception, fertility awareness, and sterilization. Effectiveness varies; LARC (IUDs, implant) is most effective (>99 percent). Considerations: efficacy, side effects, non-contraceptive benefits (cycle regulation, dysmenorrhea, acne, endometriosis), and medical eligibility (WHO MEC—thrombosis risk with estrogen, etc.). Counseling, STI prevention, and follow-up are essential components.

\medterm{Bishop Score}-- A pre-induction cervical scoring system that predicts the success of labor induction. Five components: cervical dilatation (0-3), effacement (0-3), station (-3 to +2, scored 0-3), consistency (firm, medium, soft; scored 0-2), and position (posterior, mid, anterior; scored 0-2). Total score: 0-13. Favorable (\(\ge\)6-8): high likelihood of successful induction. Unfavorable (\(\le\)5): cervical ripening needed (prostaglandins or mechanical dilation with a Foley balloon). The Bishop score is used to choose the induction method and predict the risk of failed induction.

\medterm{B-Lynch Suture}-- A surgical technique for controlling postpartum hemorrhage due to uterine atony unresponsive to medical management. A compressive suture placed around the uterus to mechanically compress the uterine walls and reduce blood flow. The suture creates a brace-like effect. Can be combined with other compression sutures (Cho, Hayman). Often successful in avoiding hysterectomy. Requires laparotomy and is performed under general anesthesia.

\medterm{Braxton Hicks Contractions}-- Irregular, painless uterine contractions occurring throughout pregnancy, becoming more frequent in the third trimester. Usually last 30-60 seconds and are felt as a tightening of the abdomen. Distinguished from true labour contractions by their irregularity, lack of cervical change, and resolution with rest or hydration. Named after John Braxton Hicks, the English physician who first described them in 1872. They do not cause progressive cervical dilatation or effacement. Often confused with prodromal labour or preterm labour, especially in primiparous women. Increased frequency at the end of the day or with physical activity. No treatment is necessary; reassurance and hydration are sufficient.

\medterm{BRCA}-- BReast CAncer gene: BRCA1 and BRCA2 are tumor suppressor genes involved in DNA double-strand break repair via homologous recombination. Mutations in these genes confer a significantly increased lifetime risk of breast cancer (up to 70-80\% for BRCA1, 40-70\% for BRCA2) and ovarian cancer (up to 40-50\% for BRCA1, 10-20\% for BRCA2). In obstetrics and gynecology, BRCA testing is offered to women with a strong family history of breast or ovarian cancer. Management options for mutation carriers: enhanced breast surveillance (MRI and mammography), risk-reducing salpingo-oophorectomy (RRSO, typically after childbearing is complete, around age 35-40 for BRCA1, 40-45 for BRCA2), and risk-reducing mastectomy.

\medterm{Breast Engorgement}-- Overfilling of the breasts with milk, causing bilateral pain, firmness, warmth, and nipple flattening. Occurs 3-5 days postpartum. Pathophysiology: vascular congestion and edema before milk secretion is established. Management: frequent feeding or pumping, cold compresses between feedings, warm compresses before feeding, and NSAIDs for pain relief. Differentiate from mastitis (usually unilateral, with fever and erythema). Engorgement is a common cause of early breastfeeding cessation if not managed properly.

\medterm{Breastfeeding}-- The feeding of an infant with breast milk, recommended exclusively for the first 6 months and continued with complementary foods to 2 years and beyond. Benefits: infant (nutrition, immune protection—IgA, reduced infection/allergy/obesity, bonding) and maternal (uterine involution, contraception, reduced breast/ovarian cancer). Physiology: prolactin (milk production) and oxytocin (milk ejection); demand feeding. Challenges: latch difficulties, sore nipples, engorgement, mastitis, insufficient milk, and work/social barriers. Contraindications: maternal HIV (in high-resource settings), active TB, chemotherapy/radiation, and infant galactosemia; some medications contraindicated. Support: lactation consultants, positioning, frequent feeding.

\medterm{Breech Delivery}-- Vaginal delivery of a breech presentation. Types: assisted breech (the baby delivers spontaneously to the umbilicus, then the operator assists delivery of the shoulders and head) and total breech extraction (the operator delivers the entire baby by grasping the feet and pulling). Total breech extraction is reserved for the second twin or extreme prematurity. Prerequisites for planned breech delivery: frank or complete breech, estimated fetal weight 2000-3800 g, adequate pelvis, head flexed, and no fetal anomalies. The Term Breech Trial (2000) showed increased perinatal mortality with planned vaginal breech delivery, leading to a sharp decline in vaginal breech births.

\medterm{Breech Presentation}-- A fetal presentation where the buttocks or feet are positioned to deliver first. Types include frank breech (hips flexed, knees extended, most common), complete breech (hips and knees flexed), and footling breech (one or both feet presenting). Associated with increased risk of cord prolapse, head entrapment, and birth trauma. External cephalic version attempted at 36-38 weeks. Cesarean delivery recommended for persistent breech at term.

\medterm{Brow Presentation}-- A fetal presentation where the fetal head is in a military position between full flexion (vertex) and hyperextension (face). The presenting diameter is the largest (occipitofrontal, ~13.5 cm). Incidence: 1:1000-1:2000. Vaginal delivery is impossible with a persistent brow in a term fetus unless it converts to vertex or face. Diagnosis: palpation of the anterior fontanelle and orbital ridges side by side. Management: conversion to vertex by applying fundal pressure and flexing the head or conversion to face by encouraging hyperextension. If conversion fails, cesarean delivery.

\medterm{CA125}-- Cancer Antigen 125, a protein tumor marker used primarily in the evaluation of ovarian cancer. CA125 is elevated in approximately 80\% of epithelial ovarian cancers but can also be elevated in benign conditions (endometriosis, pelvic inflammatory disease, uterine fibroids, pregnancy, menstruation) and other cancers (pancreatic, breast, lung). In premenopausal women, the positive predictive value of CA125 for ovarian cancer is low due to frequent false positives. The Risk of Malignancy Index combines CA125, ultrasound findings, and menopausal status to stratify ovarian cancer risk. CA125 is also used to monitor response to ovarian cancer treatment and detect recurrence.

\medterm{Caesarean Birth}-- An operative procedure in which a baby is delivered through an incision in the abdominal wall and uterus. Indications: fetal distress, failure to progress in labor, breech presentation, placenta praevia, multiple pregnancy, previous cesarean, and maternal request. Types: planned (elective) cesarean, typically performed at 39 weeks; emergency cesarean, performed when urgent delivery is needed for maternal or fetal compromise. Procedure: Joel-Cohen or Pfannenstiel skin incision, lower uterine segment transverse incision (isthmic), delivery of the baby, delivery of the placenta, and closure in layers. Regional anesthesia (spinal or epidural) is preferred. Recovery: 2-4 day hospital stay, 6-week activity restriction.

\medterm{Cardiotocography (CTG)}-- A continuous electronic recording of the fetal heart rate and uterine contractions used to assess fetal well-being during pregnancy and labor. The CTG machine uses two transducers placed on the maternal abdomen: an ultrasound transducer records the fetal heart rate, and a tocodynamometer records uterine contractions. CTG interpretation uses the four features of the fetal heart rate: baseline rate (normal 110-160 bpm), baseline variability (normal 5-25 bpm), accelerations (increase of \(\ge\)15 bpm for \(\ge\)15 seconds, reassuring), and decelerations (early, late, or variable). CTG is categorized as normal (reassuring), suspicious (one non-reassuring feature), or pathological (two or more non-reassuring or abnormal features). Fetal blood sampling may be used to confirm acidosis in suspicious/pathological CTG.

\medterm{Cardiovascular Changes in Pregnancy}-- Plasma volume increases 40-50\% (total volume ~5 L), red cell mass increases 20-30\% (physiologic anemia of pregnancy due to dilution). Cardiac output increases 30-50\% by 32 weeks (stroke volume increases initially, then heart rate). Systemic vascular resistance decreases (progesterone-mediated vasodilation) causing a slight decrease in blood pressure despite increased cardiac output. Supine hypotensive syndrome: the gravid uterus compresses the IVC, reducing venous return and cardiac output. Auscultation: S3 gallop, systolic flow murmur. ECG: left axis deviation.

\medterm{Category I Fetal Heart Tracing}-- A normal fetal heart rate tracing by NICHD 2008 criteria. All three characteristics present: baseline 110-160 bpm, moderate variability (6-25 bpm), and no late or variable decelerations. May have accelerations. Category I tracings are associated with normal fetal acid-base status. Routine management: continue labor management; no intervention needed. Predictive of normal fetal oxygenation.

\medterm{Category II Fetal Heart Tracing}-- An indeterminate fetal heart rate tracing that falls between Category I and III. Includes a wide range of patterns: baseline tachycardia or bradycardia (100-110 bpm), minimal or absent variability without decelerations, absent variability with recurrent decelerations, marked variability, variable decelerations with slow return to baseline or with other concerning features, and prolonged decelerations. Category II tracings are the most common in labor (70-80\%). Management: intrauterine resuscitation (position change, maternal oxygen, IV fluids, reduce or stop oxytocin), further evaluation (fetal scalp stimulation, scalp pH), and continued close observation. Progression to Category III warrants intervention.

\medterm{Category III Fetal Heart Tracing}-- An abnormal fetal heart rate tracing indicating potential fetal hypoxia or acidosis. Two specific patterns: (1) absent baseline variability with recurrent late decelerations, (2) absent baseline variability with recurrent variable decelerations, and (3) sinusoidal pattern. Category III tracings require prompt evaluation and intervention: intrauterine resuscitation measures, and if the pattern does not resolve, immediate delivery (operative vaginal or cesarean) to prevent fetal acidemia and neurologic injury.

\medterm{Cell-Free Fetal DNA (cfDNA) Screening}-- Also known as noninvasive prenatal testing (NIPT). Analyzes cell-free fetal DNA (cffDNA) circulating in maternal blood, derived from placental trophoblast apoptosis. Performed from 10 weeks onward. Highly sensitive and specific for trisomy 21 (>99\%, FPR <0.1\%), trisomy 18 (>97\%), and trisomy 13 (>90\%). Also screens for sex chromosome aneuploidies and microdeletions (22q11.2, 1p36, Prader-Willi/Angelman). False positives can occur from placental mosaicism, vanishing twin, or maternal chromosomal abnormalities. NIPT is a screening test, not diagnostic; positive results require confirmatory diagnostic testing (CVS or amniocentesis).

\medterm{Cephalhaematoma}-- A collection of blood between the periosteum and the skull bone, caused by rupture of blood vessels during childbirth, most commonly associated with ventouse (vacuum) delivery. Unlike caput succedaneum, a cephalhaematoma does not cross suture lines and is limited by the periosteal attachments. It presents as a firm, unilateral swelling on the newborn's scalp, appearing within hours to days after birth. It resolves spontaneously over weeks to months as the blood is reabsorbed. Complications: jaundice (from breakdown of red blood cells), anemia, and rarely infection or calcification. No treatment is required unless complications develop.

\medterm{Cerclage}-- A surgical placement of a suture around the cervix to prevent preterm delivery in patients with cervical insufficiency. Types include transvaginal (McDonald or Shirodkar) and transabdominal. Usually placed at 12-14 weeks gestation and removed at 36-37 weeks. Indicated for history of painless second trimester loss, cervical dilatation on exam, or short cervix on ultrasound. Emergency cerclage performed when dilatation occurs in the second trimester.

\medterm{Cerebroplacental Ratio (CPR)}-- The ratio of the middle cerebral artery (MCA) pulsatility index to the umbilical artery (UA) pulsatility index. CPR <5th percentile (or <1.0-1.1) suggests fetal brain-sparing effect (vasodilation of cerebral vessels in response to hypoxemia) due to placental insufficiency. A low CPR in IUGR fetuses predicts adverse perinatal outcomes (stillbirth, acidosis, NICU admission) better than UA Doppler alone. CPR may also be abnormal in small-for-gestational-age fetuses before UA Doppler becomes abnormal.

\medterm{Cervical Cerclage}-- Surgical placement of a non-absorbable suture around the cervix to provide mechanical support and prevent preterm delivery in patients with cervical insufficiency. Most commonly performed using the McDonald technique (purse-string suture at the cervicovaginal junction without bladder dissection) or the Shirodkar technique (with bladder and rectal dissection to place the suture at the internal os). Elective cerclage is placed at 12-14 weeks gestation, with removal at 36-37 weeks or at the onset of labour. Emergency or rescue cerclage may be performed in the second trimester if cervical dilation is detected before 24 weeks. Contraindications include chorioamnionitis, active vaginal bleeding, preterm labour, preterm prelabour rupture of membranes, and fetal demise. Risks include membrane rupture, chorioamnionitis, cervical trauma, and preterm labour.

\medterm{Cervical Dilatation}-- Progressive opening of the cervical os during labor, measured in centimeters from 0 (closed) to 10 cm (fully dilated). The cervix must also efface (thin out) before significant dilatation occurs. Cervical status assessed by digital vaginal examination during labor.

\medterm{Cervical Dilation}-- The opening of the cervical canal during labor, measured in centimeters from closed (0 cm) to fully dilated (10 cm); required for vaginal delivery. Occurs with effacement (thinning/shortening of the cervix) during the first stage of labor, driven by effective uterine contractions. Assessed by vaginal examination (digital—dilation, effacement, station, position). Progress: active phase typically dilates ~1-1.2 cm/hour (nullipara) and 1.5 cm/hour (multipara); abnormal dilation (protracted/arrested) indicates labor dystocia, requiring augmentation, reassessment for cephalopelvic disproportion, or cesarean.

\medterm{Cervical Effacement}-- Thinning and shortening of the cervix during labor, measured as a percentage from 0\% (thick, normal length) to 100\% (paper thin). Effacement precedes and accompanies cervical dilatation. Nulliparous women often efface before dilating; multiparous women may efface and dilate simultaneously. Effacement occurs as the internal os and external os are pulled upward into the lower uterine segment. Assessed by digital vaginal examination. In nulliparous women, effacement typically occurs before significant dilatation; in multiparous women, effacement and dilatation often occur concurrently. Contributes to the Bishop score for predicting labor induction success.

\medterm{Cervical Exam in Labor}-- Serial digital vaginal examinations to assess cervical dilatation (cm), effacement (\%), station (relationship to ischial spines, -5 to +5), consistency (firm, medium, soft), and position (posterior, mid, anterior). Performed every 2-4 hours in active labor. In the second stage, assessing fetal head position and asynclitism. Risks: infection with frequent exams (especially with ruptured membranes). Current guidelines: minimize exams after ROM, limit to q4h in active labor unless clinically indicated.

\medterm{Cervical Incompetence}-- The inability of the cervix to maintain pregnancy in the second trimester due to painless cervical dilatation. Results in recurrent second trimester pregnancy loss or preterm delivery. Risk factors include prior cervical trauma, cone biopsy, and congenital anomalies. Management includes cerclage (McDonald or Shirodkar) typically placed at 12-14 weeks and removed at 36-37 weeks. Cervical length screening by transvaginal ultrasound in high-risk patients.

\medterm{Cervical Length Screening}-- Transvaginal ultrasound measurement of the cervix between 18-24 weeks to assess risk of spontaneous preterm birth. Normal cervix measures greater than 25 mm. Short cervix (less than 25 mm) increases preterm delivery risk. Cervical funneling and dilatation are additional concerning findings. Interventions for short cervix include progesterone supplementation, cerclage, and cervical pessary. Most beneficial in asymptomatic women with prior preterm birth.

\medterm{Cervical Mucus}-- Secretions from the cervical glands that change throughout the menstrual cycle under hormonal influence. Estrogen produces thin, clear, stretchy mucus favorable for sperm transport around ovulation (spinnbarkeit). Progesterone produces thick, sticky, opaque mucus that forms a plug in the cervical canal during pregnancy. The mucus plug (operculum) protects against ascending infection throughout gestation.

\medterm{Cervical Pregnancy}-- A rare ectopic pregnancy (0.1-1\%) implanted in the cervical canal. Presents with painless vaginal bleeding in the first trimester. Diagnosis: ultrasound (hourglass-shaped uterus, gestational sac in the cervix with absent sliding sign, empty uterine cavity). High risk of severe hemorrhage. Treatment: methotrexate (systemic or local injection), uterine artery embolization, dilation and curettage with Foley balloon tamponade, or hysterectomy for uncontrolled bleeding.

\medterm{Cervical Ripening Agents}-- Medications or mechanical devices used to soften and dilate the cervix before induction of labor when the Bishop score is unfavorable (\(\le\)5). Pharmacologic: prostaglandin E1 (misoprostol/Cytotec, 25 mcg PV q4h), prostaglandin E2 (dinoprostone/Cervidil vaginal insert 10 mg over 12 hours, or Prepidil gel 0.5 mg intracervical q6h). Mechanical: transcervical Foley catheter (16-30 Fr, 30-60 mL balloon), osmotic dilators (laminaria, Dilapan-S). Membrane stripping (sweeping) is a low-risk office procedure. Combination of pharmacologic and mechanical methods is common. Contraindication: prior classical cesarean (risk of uterine rupture with prostaglandins). Misoprostol has the highest risk of tachysystole.

\medterm{Cervicitis}-- Inflammation of the uterine cervix, most commonly caused by sexually transmitted infections (Chlamydia trachomatis and Neisseria gonorrhoeae) but also by other organisms (Mycoplasma genitalium, Trichomonas vaginalis, herpes simplex virus), or non-infectious causes (chemical irritation, foreign body, allergy). Presents with mucopurulent vaginal discharge, cervical friability (bleeding on touch), intermenstrual or postcoital bleeding, and dyspareunia. Many cases are asymptomatic. Diagnosis: microscopic examination of cervical secretions showing increased polymorphonuclear leukocytes, nucleic acid amplification tests (NAATs) for Chlamydia and Neisseria gonorrhoeae, and cervical cultures. Treatment includes appropriate antibiotics (azithromycin or doxycycline for Chlamydia, ceftriaxone for gonorrhea). Sexual partners require treatment and screening. Untreated cervicitis can ascend to cause pelvic inflammatory disease, leading to tubal factor infertility, ectopic pregnancy, and chronic pelvic pain. Barrier contraception reduces transmission risk.

\medterm{Cervix}-- The lower, narrow part of the uterus that opens into the vagina. The cervix is approximately 3-4 cm long in the non-pregnant state and consists of two parts: the ectocervix (visible portion protruding into the vagina, covered by stratified squamous epithelium) and the endocervical canal (lined by columnar epithelium). The transformation zone (squamocolumnar junction) is where cervical cell changes and cancer typically arise. Functions: produces cervical mucus that changes in consistency during the menstrual cycle to facilitate or impede sperm transport; dilates during labor to allow passage of the baby; and provides mechanical support to maintain pregnancy.

\medterm{Cesarean Delivery}-- Surgical delivery of the fetus through abdominal and uterine incisions. Performed when vaginal delivery is contraindicated or fails. Indications include malpresentation, placenta previa, previous cesarean, fetal distress, and failed induction. Most common uterine incision is low transverse (classical for very preterm or transverse lie). Regional anesthesia preferred. Complications include infection, hemorrhage, thromboembolism, and injury to adjacent organs.

\medterm{Cesarean Section}-- C-section: surgical delivery of the infant through incisions in the abdominal wall and uterus. Indications: fetal distress, failure to progress, malpresentation (breech, transverse), placenta previa, abruption, previous cesarean (some), multiple pregnancy, preeclampsia, cord prolapse, and maternal conditions. Types: lower segment (transverse, most common—less blood loss) vs. classical (vertical). Procedure: regional (spinal/epidural) or general anesthesia, Pfannenstiel incision, bladder reflection, uterine incision, delivery, placental removal, and layered closure. Complications: hemorrhage, infection (wound, endometritis), thrombosis, adhesions, placenta accreta in future pregnancies, and neonatal respiratory issues. Recovery longer than vaginal delivery.

\medterm{Chignon}-- A temporary swelling (edema) on the baby's scalp caused by the vacuum cap used during ventouse (vacuum) assisted delivery. The chignon forms as the vacuum creates negative pressure pulling the scalp tissue into the cup. It settles within 24-48 hours after birth without any treatment. It should be distinguished from cephalhaematoma (subperiosteal bleed) and caput succedaneum (generalized scalp edema). No specific management is needed; parents should be reassured.

\medterm{Chorioamnionitis}-- Infection of the chorion and amnion, typically due to ascending bacterial infection from the genital tract. Risk factors include prolonged rupture of membranes, multiple vaginal exams, and preterm labor. Presents with maternal fever, uterine tenderness, and fetal tachycardia. Diagnosis is clinical; amniocentesis for confirmation. Treatment includes intravenous antibiotics (ampicillin and gentamicin) and delivery. Neonatal sepsis is a major concern.

\medterm{Chorion}-- The outermost fetal membrane, derived from the trophoblast. The chorion frondosum (villous chorion) develops into the fetal portion of the placenta. The chorion laeve (smooth chorion) is the non-villous portion. Chorionic villi contain fetal blood vessels and are bathed in maternal blood. Chorionic villus sampling (CVS) biopsies these villi for prenatal genetic diagnosis. The chorion is separated from the amnion by the extraembryonic coelom until the amnion fuses with it at ~12-14 weeks.

\medterm{Chorionic Villus Sampling (CVS)}-- A first-trimester invasive prenatal diagnostic procedure (10-13 weeks) that samples chorionic villi for genetic analysis. Transcervical (catheter through the cervix into the placenta under ultrasound guidance) or transabdominal (needle through the abdominal wall). Advantages over amniocentesis: earlier results (24-48 hours for rapid FISH, 7-14 days for karyotype). Indications: advanced maternal age, abnormal NIPT, prior aneuploidy, family history of genetic disorders, and abnormal ultrasound findings. Risks: miscarriage (0.2-0.5\% above background), and limb reduction defects (rare, associated with CVS before 10 weeks). Confined placental mosaicism (CPM, 1-2\%): karyotype differs between placenta and fetus, requiring amniocentesis for confirmation.

\medterm{Choroid Plexus Cyst (CPC)}-- A sonolucent fluid-filled cyst within the choroid plexus of the lateral cerebral ventricle, seen in 1-3\% of second-trimester fetuses. In isolation, it is a normal variant with no clinical significance. CPC is associated with trisomy 18 (30-50\% of trisomy 18 fetuses have CPC, but most fetuses with CPC are normal). If isolated, the risk of trisomy 18 is very low, particularly if NIPT is normal. No further workup is needed for isolated CPC after 22 weeks gestation.

\medterm{Chromosomal Abnormality}-- See genetics.

\medterm{Chronic Hypertension in Pregnancy}-- Hypertension (BP \(\ge\)140/90 mmHg) that predates pregnancy or is diagnosed before 20 weeks gestation. Affects 1-5\% of pregnancies. Increased risk of superimposed preeclampsia (20-50\%), IUGR, preterm birth, and placental abruption. Management: continue or adjust antihypertensives (avoid ACE inhibitors and ARBs during pregnancy; labetalol, nifedipine, and methyldopa are safe). Low-dose aspirin (81 mg daily) from 12-36 weeks reduces the risk of preeclampsia. Delivery timing depends on blood pressure control and fetal status.

\medterm{Clavicle Fracture (Intentional)}-- Intentional fracture of the fetal clavicle as a last resort to reduce the bisacromial diameter in shoulder dystocia. The operator applies upward pressure on the midportion of the anterior fetal clavicle with a finger, fracturing it. The clavicle heals without long-term sequelae. Usually unnecessary because the posterior arm delivery and Gaskin maneuver are more effective and less traumatic. Accidental clavicle fracture occurs in 0.5-2\% of normal deliveries.

\medterm{Colostrum}-- The thick, yellowish fluid produced by the mammary glands in late pregnancy and the first few days postpartum. Rich in immunoglobulins (especially IgA), proteins, leukocytes, and growth factors.

\medterm{Combined Spinal-Epidural (CSE) Analgesia}-- A neuraxial analgesia technique for labor that combines a spinal injection (short-acting opioid or opioid + low-dose local anesthetic) with an epidural catheter. Advantages: rapid onset of pain relief (3-5 minutes from the spinal component), minimal motor block, and the ability to dose through the epidural catheter for prolonged labor. The spinal block provides immediate relief for women in active labor; the epidural provides sustained relief. The CSE may be associated with more maternal pruritus (from intrathecal opioids) and a higher rate of fetal bradycardia than epidural alone.

\medterm{Compound Presentation}-- A fetal presentation where an extremity (hand, arm, foot) presents alongside the vertex or breech. Most common: hand alongside the head. Incidence: 1:300-1:1000. Fetal prognosis is generally good; the prolapsed extremity often retracts as labor progresses. Management: expectant management in most cases. The prolapsed hand may be gently pushed back after descent. The cord should be assessed for prolapse. Cesarean if labor fails to progress or the extremity causes obstruction.

\medterm{Conception (Fertilization)}-- The fusion of a sperm and an ovum in the ampulla of the fallopian tube, forming a zygote. The sperm penetrates the corona radiata and zona pellucida (acrosome reaction). After fertilization, the zygote undergoes cleavage (mitotic division) as it travels toward the uterus. Implantation occurs approximately 6-10 days after ovulation. Conception timing is critical for dating pregnancy and determining the window of fertility.

\medterm{Conjoined Twins}-- Monozygotic, monochorionic, monoamniotic twins who are physically connected at birth, resulting from incomplete fission of the fertilized ovum after day 13 of fertilization. Incidence: 1 in 50,000-100,000 births. Classification by the site of union: thoracopagus (chest, most common, 40\%, often with shared heart — poor prognosis), omphalopagus (abdomen, 33\%), pygopagus (sacrum/buttocks, 19\%), ischiopagus (pelvis, 6\%), craniopagus (head, 2\%), and parapagus (lateral fusion). Determination of shared vital organs (especially the heart) is critical for prognosis and surgical planning. Diagnosis by prenatal ultrasound and MRI. Delivery is typically by planned cesarean at 34-36 weeks. Surgical separation after birth is complex and requires a multidisciplinary team. Prognosis depends on the degree of organ sharing.

\medterm{Continuous Electronic Fetal Monitoring (EFM)}-- The continuous recording of fetal heart rate (FHR) and uterine contractions using external transducers (Doppler ultrasound and tocodynamometer) or internal devices (fetal scalp electrode and intrauterine pressure catheter). Used intrapartum to detect fetal hypoxia. The NICHD 2008 classification categorizes tracings as Category I (normal), II (indeterminate), or III (abnormal). EFM is associated with an increased rate of cesarean and operative vaginal delivery without clear reduction in cerebral palsy (false positives are common). Recommended for high-risk labor and women who choose augmentation. Intermittent auscultation is appropriate for low-risk labor.

\medterm{Contraception}-- Methods to prevent pregnancy. Hormonal: combined oral contraceptive pill (COC), progestin-only pill, depot medroxyprogesterone, implant (etonogestrel), vaginal ring, patch; IUDs (copper, levonorgestrel); barrier (male/female condom, diaphragm, cervical cap, spermicide); emergency contraception (levonorgestrel, ulipristal, copper IUD); fertility awareness; and permanent (sterilization—tubal ligation, vasectomy). Selection: efficacy, reversibility, side effects, and medical eligibility (WHO MEC). Long-acting reversible contraception (LARC—IUDs, implant) has the highest effectiveness. Counseling includes STI protection (condoms), risk assessment (thrombosis—COC, migraine with aura), and patient preference.

\medterm{Contraception and Birth Control}-- Contraception encompasses methods used to prevent pregnancy, including hormonal options like oral contraceptives, patches, rings, injections, and implants, as well as non-hormonal methods like copper IUDs, barrier methods, diaphragms, and permanent sterilization. Emergency contraception is available for use shortly after unprotected intercourse. Method selection considers effectiveness, side effects, reversibility, and individual health factors.

\medterm{Contractions}-- Rhythmic uterine muscular contractions of labor that produce cervical dilation and effacement, driving descent of the fetus. Characterized by frequency, duration, and intensity; progress is monitored by partogram and tocography (CTG). Braxton Hicks contractions: irregular, non-progressive, painless prelabor contractions. Abnormalities: inadequate/hypotonic (dystocia—augmentation with oxytocin), hypertonic/tetanic (fetal distress, placental abruption risk), and preterm contractions (threatened preterm labor). Management: assess cervical change to confirm labor, monitor fetal well-being during contractions (CTG), analgesia, and address abnormalities (augmentation, tocolysis for preterm labor if indicated).

\medterm{Contraction Stress Test}-- An antepartum fetal surveillance technique assessing fetal heart rate response to induced uterine contractions. Oxytocin or nipple stimulation produces contractions. A negative test shows no late decelerations. A positive CST demonstrates late decelerations with more than 50\% of contractions. Equivocal results show intermittent late decelerations. Contraindications include preterm labour, placenta previa, previous classical caesarean, multiple gestation, and preterm prelabour rupture of membranes. The CST has largely been replaced by the biophysical profile and modified BPP as first-line antepartum surveillance due to its higher false-positive rate, greater time requirement, and need for intravenous access. However, it remains useful when BPP is equivocal or when fetal status is uncertain.

\medterm{Contraction Stress Test (CST)}-- An antepartum test that evaluates the fetal heart rate response to uterine contractions. Contractions are induced with oxytocin or nipple stimulation (until 3 contractions in 10 minutes). Negative (normal): no late or significant variable decelerations. Positive (abnormal): late decelerations with \(\ge\)50\% of contractions. Equivocal: intermittent late decelerations. The CST is less commonly used due to the higher success and lower risk of the BPP. Indications: when BPP is equivocal or maternal perception of decreased movement.

\medterm{Controlled Cord Traction (CCT)}-- A technique for active management of third-stage labor: gentle, steady traction on the umbilical cord with one hand while the other hand applies suprapubic counter-pressure (to prevent uterine inversion). Performed after signs of placental separation. The Brandt-Andrews maneuver and the Finocchietti technique are CCT methods. CCT reduces the duration of the third stage and the risk of postpartum hemorrhage. Contraindication: signs of uterine inversion.

\medterm{Cord Presentation}-- The presence of the umbilical cord between the presenting fetal part and the cervix, with intact membranes. Risk factors: malpresentation, unstable lie, low-lying placenta, and long cord. If the membranes rupture, the cord may prolapse (umbilical cord prolapse). Diagnosis: digital exam (palpation of a pulsating cord through intact membranes) or ultrasound. Management: admit to the hospital, strict bed rest, avoid vaginal exams, and cesarean delivery before membrane rupture.

\medterm{COVID-19 in Pregnancy}-- Pregnant women with COVID-19 are at increased risk of severe illness (ICU admission, mechanical ventilation, death) compared to non-pregnant women. Risks to the fetus: preterm birth, cesarean delivery, and possibly stillbirth. Vertical transmission is rare but possible. Management: vaccination is recommended (no increased risk of miscarriage or adverse outcomes). Treatment: remdesivir, dexamethasone (for oxygen requirement), and anticoagulation as indicated. Pregnant women should be offered vaccination at any stage of pregnancy.

\medterm{Crowning}-- The stage of the second stage of labor when the fetal head (widest diameter) distends the perineum and becomes visible at the introitus and is no longer receding between contractions. At crowning, controlled delivery of the head (guard the perineum, support, or episiotomy if needed) minimizes perineal trauma; the mother is guided not to push forcefully and to pant to allow controlled expulsion. Following the head, shoulders and body deliver. Complete breech and other presentations differ. Crowning precedes imminent delivery; neonatal and maternal readiness is required.

\medterm{Crown-Rump Length (CRL)}-- The measurement from the top of the fetal head (crown) to the bottom of the buttocks (rump), the most accurate method for establishing gestational age in the first trimester. Measured at 7-13 weeks. Accuracy: ±3-5 days. The EDD is set using CRL if the difference from LMP dating exceeds 7 days. CRL growth: approximately 1 mm per day. CRL <7 mm without cardiac activity: uncertain viability, repeat in 1 week. Absent cardiac activity at CRL \(\ge\)7 mm is diagnostic of pregnancy failure.

\medterm{Cystocele}-- Anterior vaginal wall prolapse: herniation of the bladder into the anterior vaginal wall due to weakening of the pubocervical fascia, most commonly after childbirth and menopause. Features: vaginal bulge/pressure, difficulty voiding, incomplete emptying, urinary urgency/frequency, and stress urinary incontinence; worse with standing/straining. Grades by POP-Q. Diagnosis: clinical. Management: pelvic floor physiotherapy, pessary, and surgical repair (anterior colporrhaphy) for symptomatic prolapse; treat urinary symptoms and address contributing factors (obesity, constipation). Complications of surgery: recurrence, dyspareunia, and urinary symptoms.

\medterm{Decidua}-- The transformed endometrium of pregnancy, shed after delivery. Three regions: decidua basalis (under the placenta, site of trophoblast invasion), decidua capsularis (over the gestational sac, covering the embryo), and decidua parietalis (the remainder of the uterine lining). The decidua provides nutrition to the early embryo and produces hormones (prolactin, relaxin). Decidual cells and natural killer cells play key roles in immune tolerance of the fetus.

\medterm{Delayed Cord Clamping}-- The practice of delaying umbilical cord clamping for \(\ge\)30-60 seconds after delivery, allowing placental transfusion of blood to the newborn. Benefits: increases neonatal blood volume by 30\%, improves iron stores (ferritin, hemoglobin) for the first 6 months, reduces anemia, and improves transitional circulation. For preterm infants (<34 weeks), delayed clamping for 30-60 seconds reduces intraventricular hemorrhage, necrotizing enterocolitis, and need for blood transfusion. Contraindications: need for immediate resuscitation, cord avulsion, placental abruption with hemorrhage, and maternal instability.

\medterm{Delivery}-- The process of childbirth: vaginal delivery (spontaneous vertex most common; assisted—vacuum, forceps) or cesarean delivery (abdominal). Stages of labor culminate in delivery of the infant and placenta. Considerations: fetal presentation (vertex, breech, face), position, and station; maternal pushing and expulsive forces; perineal trauma (tears/episiotomy); and immediate neonatal care (Apgar, cord clamping). Complications: shoulder dystocia, postpartum hemorrhage, retained placenta, cord prolapse. Management: preparation (delivery pack, neonatal resuscitation), controlled delivery of head/shoulders, third-stage management (oxytocin to reduce hemorrhage), and documentation. Facility-based delivery with skilled attendance reduces maternal/neonatal mortality.

\medterm{Delivery of Posterior Arm}-- A shoulder dystocia maneuver: the operator's hand enters the posterior vagina along the fetal chest, finds the fetal hand, grasps the forearm, sweeps it across the chest and face, and delivers the posterior arm and hand over the perineum. This reduces the bisacromial diameter by 20\% and is often highly effective. Flexing the elbow helps retrieve the hand. Risk: humerus fracture (treat with immobilization; heals well). Often the definitive maneuver after McRoberts and rotational maneuvers fail.

\medterm{Dexamethasone}-- An alternative corticosteroid for fetal lung maturation, similar to betamethasone. Dose: four doses of 6 mg IM, 12 hours apart. While betamethasone is the preferred agent in the US (more evidence for RDS reduction), dexamethasone is used in some settings and is effective. Both cross the placenta and induce fetal lung surfactant production. Dexamethasone has slightly more evidence for long-term neurodevelopmental outcomes in certain studies. WHO recommends either agent.

\medterm{Dichorionic Twins}-- Twins each having their own chorion and amnion, occurring in approximately 30\% of monozygotic twins and all dizygotic twins. Each twin has its own placenta, though placentas may be adjacent. Complications include growth restriction, preterm delivery, and preeclampsia, but not twin-to-twin transfusion syndrome. Require regular ultrasound monitoring for growth discordance.

\medterm{Dinoprostone (Cervidil, Prepidil)}-- A synthetic prostaglandin E2 (PGE2) used for cervical ripening and induction. Forms: Cervidil (10 mg vaginal insert, slow-release over 12 hours, removed if tachysystole occurs), Prepidil (0.5 mg intracervical gel, may be repeated q6h), and Prostin E2 (vaginal suppository). Indications: unfavorable cervix (Bishop \(\le\)5). Side effects: tachysystole (5-10\%), fever, nausea, vomiting, and diarrhea. Contraindications: prior uterine surgery (including classical cesarean), active vaginal bleeding, and fetal compromise. Remove promptly if tachysystole or nonreassuring FHR occurs.

\medterm{Donor Insemination}-- A fertility treatment in which sperm from a donor (anonymous or known) is placed into a woman's vagina, cervix, or uterus to achieve pregnancy. Indications: male factor infertility (azoospermia, severe oligospermia), genetic disorders in the male partner, single women, and same-sex female couples. In the UK, donor insemination is regulated by the Human fertilization and Embryology Authority, requiring donor screening for infectious diseases, genetic disorders, and medical conditions. Recipients may require ovulation induction or intrauterine insemination (IUI) to optimize timing.

\medterm{Douche}-- A stream of water or solution directed into a body cavity, most commonly the vagina, for cleansing, therapeutic, or contraceptive purposes. Medical use: therapeutic douches with antiseptic solutions (povidone-iodine, vinegar, chlorhexidine) for treating specific vaginal infections (trichomoniasis, bacterial vaginosis) as adjunctive therapy. Routine vaginal douching is not recommended and is associated with adverse outcomes: disruption of normal vaginal flora (increased risk of bacterial vaginosis, candida overgrowth), pelvic inflammatory disease, ectopic pregnancy, and increased HIV transmission risk.

\medterm{Doula Support}-- A trained non-medical birth companion who provides continuous emotional, physical, and informational support to the laboring mother. Services: massage, positioning guidance, breathing techniques, advocacy, and reassurance. Doula support has Level 1 evidence for improved outcomes: reduced duration of labor (by ~40 minutes), reduced need for cesarean (by 25\%), reduced epidural requests, reduced Pitocin use, and higher satisfaction. Doulas do not perform clinical tasks (vaginal exams, FHR monitoring). They complement but do not replace midwives or physicians.

\medterm{Early Decelerations}-- A fetal heart rate pattern with a gradual decrease (onset to nadir >30 seconds) that mirrors the contraction and returns to baseline as the contraction ends. Caused by fetal head compression during contractions (vagal response). A benign, physiologic finding associated with normal fetal acid-base status. No intervention required. Early decelerations are characteristic of active labor and descent of the fetal head.

\medterm{Echogenic Bowel}-- Increased echogenicity of the fetal bowel relative to adjacent bone (iliac wing). Identified in 0.1-0.5\% of second-trimester ultrasounds. Associations: aneuploidy (trisomy 21, 18, 13), cystic fibrosis (echogenic bowel may be the first sign of meconium ileus), congenital infection (CMV, toxoplasmosis, parvovirus B19), IUGR, and swallowing of intra-amniotic blood. Workup: NIPT, cystic fibrosis carrier screening, infection serologies (TORCH), and serial growth scans.

\medterm{Echogenic Intracardiac Focus (EIF)}-- A small bright spot (calcification) in the fetal heart, typically in the left ventricle (papillary muscle). Seen in 3-5\% of normal pregnancies. In isolation, it is a normal variant with minimal aneuploidy risk. However, EIF is more common in trisomy 21 fetuses (10-30\%), particularly when combined with other soft markers. If isolated and NIPT/low-risk serum screen, no further action is needed.

\medterm{Eclampsia}-- The occurrence of generalized tonic-clonic seizures in a patient with preeclampsia, representing a medical emergency. Complications include aspiration, placental abruption, disseminated intravascular coagulation, and maternal/fetal death. Prevention and treatment involves intravenous magnesium sulfate (4-6 g IV bolus followed by 1-2 g/hour maintenance), which reduces the risk of seizures by 50\%. Management includes airway protection, seizure control (magnesium sulfate first-line; benzodiazepines if magnesium fails), blood pressure control with antihypertensives (labetalol or hydralazine for BP \(\ge\)160/110 mmHg), and delivery once maternal stability is achieved. Eclampsia can occur antepartum (50\%), intrapartum (25\%), or postpartum (25\%, most within 48 hours). Recurrence risk in future pregnancies is approximately 2\%. Maternal mortality is approximately 2-5\% in high-resource settings.

\medterm{Ectopic Pregnancy}-- Implantation of a fertilized ovum outside the uterine cavity, most commonly in the fallopian tube (95\%). Risk factors: prior ectopic, PID, tubal surgery, IUD, smoking, and assisted reproductive technology. Presents with lower abdominal pain, amenorrhea, and vaginal bleeding. Rupture causes hemodynamic instability. Diagnosis: transvaginal ultrasound (no intrauterine pregnancy with hCG >discriminatory zone), serial hCG monitoring. Treatment: methotrexate for early, unruptured ectopic; salpingectomy or salpingostomy for surgical management.

\medterm{Effacement}-- Thinning (shortening) of the cervix during labor, expressed as a percentage (0-100 percent), as the cervix is taken up into the lower uterine segment, enabling dilation. Occurs with cervical dilation in the first stage of labor, driven by contractions and the presenting part. Assessed digitally on vaginal examination along with dilation and station. Effacement precedes or accompanies dilation, particularly in nulliparas (effacement first), whereas multiparas may dilate with less effacement. Abnormal (lack of progress) indicates dystocia.

\medterm{Embryo}-- The developing human organism from fertilization (conception) through the eighth week of gestation (10th week by last menstrual period). After fertilization, the zygote undergoes cleavage (cell division) as it travels through the fallopian tube, reaching the morula stage (~16 cells) and then the blastocyst stage (~5--6 days), which implants in the endometrium. During the embryonic period, all major organ systems develop (organogenesis): gastrulation forms the three germ layers (ectoderm, mesoderm, endoderm), neurulation forms the neural tube, and the heart begins beating at ~5--6 weeks. The embryonic period is the most vulnerable to teratogens (medications, infections, radiation, alcohol), which can cause major structural malformations. After 8 weeks, the developing human is called a fetus.

\medterm{Endometritis}-- Infection of the endometrium, the most common cause of puerperal fever after cesarean delivery. Risk factors include prolonged labor, multiple vaginal examinations, premature rupture of membranes, and cesarean section. Presents with fever, uterine tenderness, foul-smelling lochia, and leukocytosis. Treatment with broad-spectrum IV antibiotics (ampicillin and gentamicin plus clindamycin) for polymicrobial infection including anaerobes.

\medterm{Engagement}-- The descent of the fetal presenting part into the pelvis, typically occurring 2-3 weeks before labor in nulliparas or at the onset of labor in multiparas.

\medterm{Epidural}-- Regional analgesia (or anesthesia) delivering local anesthetic/opioid via a catheter in the epidural space (lumbar), used for pain relief during labor and for cesarean (higher dose) anesthesia. Provides excellent analgesia while allowing maternal mobility and pushing; reduces catecholamine stress. Complications: hypotension (sympathetic block), postdural puncture headache, motor block, back pain, urinary retention, infection/hematoma (rare), and fetal effects (transient). Contraindications: coagulopathy, infection at site, patient refusal, and severe hypovolemia. Alternatives: spinal (single dose), combined spinal-epidural, nitrous oxide, and systemic opioids.

\medterm{Epidural Analgesia for Labor}-- A neuraxial technique for labor analgesia: a catheter is placed in the epidural space (L3-L4 or L2-L3 interspace) and local anesthetic (bupivacaine 0.0625-0.125\%) with opioid (fentanyl 2 µg/mL) is infused continuously or in patient-controlled epidural analgesia (PCEA) boluses. Produces segmental analgesia (T10-L1 for the first stage, T10-S4 for the second stage). Benefits: excellent pain relief, mother can rest, and can be converted to surgical anesthesia for cesarean. Side effects: hypotension, motor block (reduced with low-dose solutions), fever (can increase the diagnosis of chorioamnionitis), pruritus, and urinary retention. Contraindications: maternal coagulopathy, infection at the insertion site, and untreated maternal hypovolemia.

\medterm{Episiotomy}-- A surgical incision of the perineum to enlarge the vaginal outlet during delivery. Two types: mediolateral (most common worldwide) and median/midline. Previously performed routinely; current evidence supports selective use. Indications include operative vaginal delivery, fetal distress requiring expedited delivery, and shoulder dystocia. Classified with perineal lacerations as first degree (vaginal mucosa only), second degree (mucosa and perineal muscles), third degree (involving the anal sphincter complex), or fourth degree (extending through the rectal mucosa). Healing occurs over 2-4 weeks. Complications include pain, infection, dyspareunia, and wound dehiscence. The mediolateral episiotomy is preferred over midline because it is associated with a lower risk of third- and fourth-degree tears involving the anal sphincter. Modern obstetrics recommends restrictive rather than routine use of episiotomy.

\medterm{Episiotomy Repair}-- Surgical repair of an episiotomy incision after delivery. Technique: continuous suturing of the vaginal mucosa (2-0 chromic), interrupted sutures for the perineal muscles, and subcuticular closure of the perineal skin. The rectum should be examined after repair (suture inadvertently through the rectal mucosa). Mediolateral episiotomy is repaired in the same fashion. Pain management: ice packs, topical anesthetics, sitz baths, NSAIDs. Avoid prolonged sitting and constipation.

\medterm{Estimated Due Date (EDD)}-- The expected date of delivery, calculated as 40 weeks (280 days) from the first day of the last menstrual period. Confirmed or adjusted by first-trimester ultrasound crown-rump length measurement (most accurate ±5-7 days). Second-trimester ultrasound is less accurate (±10-14 days). Third-trimester is least accurate (±21-30 days). Only ~5\% of women deliver on their exact EDD; most deliver within 2 weeks before or after. Used to determine gestational age for prenatal screening, timing of interventions, and postterm management.

\medterm{Estimated Fetal Weight (EFW)}-- An ultrasound-derived estimate of fetal weight, calculated using formulas that combine BPD, HC, AC, and FL. The Hadlock formula (using HC, AC, FL) is the most commonly used. Accuracy: ±10-15\%. Indications: assess fetal growth, determine size for gestational age, and plan delivery. EFW <10th percentile: SGA/IUGR. EFW >90th percentile: LGA/macrosomia. EFW >4500 g may influence delivery decisions (consider cesarean).

\medterm{External Cephalic Version (ECV)}-- A procedure to manually rotate a breech or transverse fetus to a cephalic (head-down) presentation from outside the abdomen. Performed after 36-37 weeks (term, with adequate amniotic fluid). Success rate: 30-80\% (higher with prior vaginal delivery, transverse lie, adequate fluid, and non-anterior placenta). Contraindications: non-reassuring fetal status, oligohydramnios, uterine anomaly, multiple gestation, and contraindications to vaginal delivery. Fetal monitoring before and after. RhoGAM for Rh-negative patients. Tocolysis (terbutaline) may increase success.

\medterm{Extrauterine Pregnancy}-- A pregnancy in which the fertilized ovum implants outside the uterine cavity, an alternative term for ectopic pregnancy. By far the most common site is the fallopian tube (ampulla 70\%, isthmus 12\%, fimbria 11\%, cornual/interstitial 2--4\%). Other extrauterine sites: ovary, cervix, cesarean scar, abdominal cavity, and the rudimentary horn of a bicornuate uterus. Risk factors: prior ectopic pregnancy, tubal surgery, pelvic inflammatory disease, endometriosis, smoking, assisted reproductive technology, and IUD in situ. Presentation: abdominal pain, amenorrhea, vaginal bleeding--classic triad; rupture causes hemoperitoneum, hypovolemic shock. Diagnosis: quantitative beta-hCG and transvaginal ultrasound. Treatment: medical (methotrexate for unruptured, low hCG, no fetal cardiac activity) or surgical (salpingectomy or salpingostomy by laparoscopy).

\medterm{Face Presentation}-- A fetal presentation where the neck is hyperextended and the face is the presenting part. Incidence: 1:500-1:1000. Types: mentum anterior (chin toward pubis, delivers vaginally), mentum posterior (chin toward the sacrum, cannot deliver vaginally unless rotates to anterior). Diagnosis: ultrasound or palpation (feel nose, eyes, mouth, and chin during vaginal exam). Management: if mentum anterior, labor may proceed; if mentum posterior and persistent, cesarean. Preterm and anencephalic fetuses have a higher rate of face presentation.

\medterm{Failed Induction of Labor}-- The inability to achieve active labor and progressive cervical change after an adequate trial of induction. Definition: no active labor (cervical change) after \(\ge\)12-18 hours of oxytocin with ruptured membranes and adequate contractions. Management: the latent phase may be prolonged (up to 24 hours or more). If the cervix does not change after adequate contractions with a favorable Bishop score and ROM, cesarean delivery is indicated. Failed induction is more common in nulliparas and with unfavorable cervical status at the beginning.

\medterm{Fecundity}-- The biological ability to conceive and produce a live birth. Fecundity is influenced by age (declines significantly after 35), ovarian reserve, tubal patency, uterine factors, and male fertility (sperm count, motility, morphology). Fecundability: the probability of achieving pregnancy in a single menstrual cycle. Normal fecundability is approximately 20-25\% per cycle in healthy couples. Infertility is defined as failure to conceive after 12 months of regular unprotected intercourse.

\medterm{Female Genital Mutilation}-- FGM: the partial or total removal of female external genitalia for non-medical reasons, a human rights violation practiced in parts of Africa, the Middle East, and Asia (types I-IV). Immediate complications: severe pain, hemorrhage, infection, urinary retention. Long-term: chronic infection, scarring, painful menstruation, sexual dysfunction, and obstetric complications (prolonged labor, perineal tears, fistula, increased maternal/neonatal mortality). Legal/ethical: illegal in most countries; mandatory reporting and safeguarding. Management: multidisciplinary care (gynecology, urology, psychosexual), deinfibulation for obstetric/functional indications, counseling, and community education/prevention. Support survivors with sensitive, non-judgmental care.

\medterm{Femur Length (FL)}-- The length of the fetal femur (diaphysis), measured from the greater trochanter to the distal femoral metaphysis. Used for dating (14-26 weeks) and growth assessment. Normal growth: ~2 mm per week. Short femur: can be constitutional (especially with short parents), associated with aneuploidy (trisomy 21, skeletal dysplasias), or IUGR. The FL/AC ratio helps differentiate IUGR (increased ratio) from constitutional smallness.

\medterm{Fertility}-- The biological capacity to conceive and produce offspring. Fertility depends on the proper functioning of the reproductive system in both males and females, including ovulation, sperm production, and hormonal regulation. Infertility affects a significant proportion of couples and may require medical evaluation and assisted reproductive technologies such as in vitro fertilization (IVF).

\medterm{Fetal Acid-Base Status}-- The assessment of fetal oxygenation and acid-base balance during labor, reflecting fetal metabolic condition. Normal: umbilical artery pH 7.25-7.35, pCO2 40-50 mmHg, base excess 0 to -5. Fetal acidosis: pH <7.20 (mild), pH <7.00 (severe, associated with neonatal encephalopathy). Respiratory acidosis: elevated pCO2 from cord compression (rapidly reversible). Metabolic acidosis: decreased base excess from prolonged hypoxia (associated with neurologic injury). Fetal scalp pH: direct measurement from fetal blood; pH >7.25 = normal, <7.20 = acidosis.

\medterm{Fetal Alcohol Spectrum Disorder (FASD)}-- A broad spectrum of physical, behavioral, and cognitive impairments caused by prenatal alcohol exposure. The most severe form is fetal alcohol syndrome (FAS): characteristic facial features (smooth philtrum, thin vermillion border, small palpebral fissures), growth deficiency (prenatal or postnatal), and central nervous system abnormalities (microcephaly, developmental delay, intellectual disability). No amount of alcohol is considered safe in pregnancy. Diagnosis: based on confirmed maternal alcohol use, dysmorphology exam, and neurodevelopmental assessment. Prevention: screening and counseling in pregnancy.

\medterm{Fetal Alcohol Syndrome}-- A collection of birth defects resulting from exposure of the fetus to alcohol during pregnancy, representing the most severe form of fetal alcohol spectrum disorder (FASD). Characteristic features include distinctive facial dysmorphology (smooth philtrum, thin vermillion border, small palpebral fissures), prenatal and postnatal growth deficiency, and central nervous system abnormalities (microcephaly, developmental delay, intellectual disability, behavioral problems). No amount of alcohol is considered safe during pregnancy. Diagnosis requires confirmed prenatal alcohol exposure, characteristic facial features, growth deficits, and neurodevelopmental impairment. Management is supportive with early intervention services.

\medterm{Fetal Attitude}-- The relationship of fetal body parts to one another, describing the degree of flexion or extension of the fetal head and extremities. Normal attitude is full flexion (vertex presentation with chin tucked to chest). attitudes include deflexion (less flexion), military (neutral position), brow, face, or complete extension. Deflexed attitudes result in larger presenting diameters, potentially complicating vaginal delivery.

\medterm{Fetal Bradycardia}-- A baseline fetal heart rate <110 bpm for \(\ge\)10 minutes. Moderate (100-110 bpm): often benign, especially in post-term fetuses. Severe (<100 bpm): concerning for fetal hypoxia, acidosis, congenital heart block (SLE, anti-Ro/anti-La antibodies), cord accident, placental abruption, or maternal seizures. A prolonged deceleration lasting >2 minutes is managed as an emergency. Congenital heart block: fixed bradycardia 50-80 bpm, unrelated to contractions, diagnosed prenatally by fetal echocardiography.

\medterm{Fetal Heart Rate}-- Normal fetal heart rate ranges from 110-160 beats per minute. Monitored via auscultation with a fetoscope, Doppler ultrasound, or continuous electronic fetal monitoring. Baseline FHR is determined over a 10-minute window and is modulated by the autonomic nervous system. Variability (fluctuations in baseline) is classified as absent, minimal (\(\le\)5 bpm), moderate (6-25 bpm), or marked (>25 bpm) and reflects fetal acid-base status. Accelerations (transient increases of \(\ge\)15 bpm for \(\ge\)15 seconds in fetuses beyond 32 weeks) are reassuring signs of fetal well-being. Decelerations are classified as early (head compression, benign), variable (cord compression, most common), or late (uteroplacental insufficiency, concerning). Fetal tachycardia (>160 bpm) may indicate maternal fever, infection, or hypoxia; bradycardia (<110 bpm) may suggest fetal distress, congenital heart block, or maternal hypotension.

\medterm{Fetal Kick Counts (Fetal Movement Counting)}-- A method of antepartum fetal surveillance where the mother counts fetal movements (kicks, rolls, stretches) daily. Cardiff count (10 kicks in 2 hours): lie on the left side, record the time until 10 movements are felt. Normal: 10 movements in \(\le\)2 hours. Decreased fetal movement is associated with increased risk of stillbirth and IUGR. A single episode of decreased fetal movement warrants NST/BPP. Maternal perception of fetal movement is generally reliable and may be the earliest sign of fetal compromise.

\medterm{Fetal Lie}-- The relationship of the long axis of the fetus to the long axis of the mother. Longitudinal lie (parallel, 99\%) or transverse lie (perpendicular, 1\%). Most common is left occiput transverse or left occiput anterior. Transverse lie at term is abnormal and requires cesarean delivery if not corrected. Lie determined by palpation and ultrasound, particularly important for planning delivery route and management of labor.

\medterm{Fetal Presentation}-- The relationship of the fetal presenting part to the maternal pelvis. Includes cephalic (head first, most common at 97\%), breech (buttocks or feet first), shoulder/oblique, and transverse lie. Cephalic presentations include vertex, face, brow, and compound. Breech presentations include complete, frank, and footling. Presentation assessed by Leopold maneuvers and confirmed with ultrasound. External cephalic version may be attempted for breech at 36-38 weeks.

\medterm{Fetal Scalp Electrode (FSE)}-- A spiral wire electrode placed on the fetal scalp through the cervix to directly record the fetal heart rate. Provides a more accurate FHR signal than external Doppler, particularly in obese patients or when external monitoring is inadequate. Indications: poor-quality external tracing, need for internal monitoring. Contraindications: HIV and hepatitis C (risk of vertical transmission), placenta previa, and scalp infection at the electrode site. The electrode is placed after membrane rupture, with cervical dilatation \(\ge\)2 cm.

\medterm{Fetal Scalp Stimulation}-- A test of fetal well-being during labor by digitally rubbing or pinching the fetal scalp during a vaginal exam. A reactive acceleration (increase of \(\ge\)15 bpm for \(\ge\)15 seconds) in the FHR in response to stimulation is a reassuring sign of normal fetal acid-base status (negative predictive value 95\% for fetal acidosis). If no acceleration, fetal scalp pH may be indicated. The test is noninvasive and can be performed repeatedly.

\medterm{Fetal Tachycardia}-- A baseline fetal heart rate >160 bpm for \(\ge\)10 minutes. Mild (161-180 bpm) or severe (>180 bpm). Causes: maternal fever (most common, often from chorioamnionitis), medications (terbutaline, atropine, scopolamine), maternal hyperthyroidism, fetal infection (myocarditis), fetal hypoxia (compensatory response), and fetal tachyarrhythmia (SVT, atrial flutter, 200-300 bpm). Management: identify and treat the cause (maternal cooling, antibiotics for chorioamnionitis). Persistent tachycardia with minimal variability is concerning.

\medterm{Fetal Viability}-- The gestational age at which a fetus has the potential to survive outside the uterus with medical support, generally considered to be 23-24 weeks of gestation. Survival rates increase significantly with each additional week: approximately 30-50\% at 23 weeks, 50-70\% at 24 weeks, and over 90\% by 27-28 weeks. Factors influencing viability include birth weight, sex, multiple gestation, access to neonatal intensive care, and presence of congenital anomalies. The limit of viability is defined by each country's legal and ethical framework. In the UK, the upper limit for termination of pregnancy is 24 weeks for most grounds. Prognosis for intact survival without severe neurodevelopmental impairment remains guarded at the earliest gestational ages, with approximately 40\% of survivors at 23 weeks having moderate to severe disability.

\medterm{Fetomaternal Hemorrhage}-- Transplacental passage of fetal red blood cells into the maternal circulation. Occurs in small amounts during normal pregnancy but can be massive during trauma, placental abruption, or cesarean delivery.

\medterm{Fetus}-- An unborn offspring from the 9th week of gestation (after the embryonic period) until birth. The fetal period is characterized by rapid growth, maturation of organ systems, and increasing complexity of function. Key fetal milestones: 12 weeks — sex distinguishable, kidneys begin producing urine; 20 weeks — fetal movements felt by mother (quickening); 24 weeks — viability (potential to survive outside the womb with intensive neonatal support); 28-32 weeks — lung surfactant production increases; 37-42 weeks — term.

\medterm{First Stage of Labor}-- From the onset of regular contractions to full cervical dilatation. Two phases: latent (0-6 cm, slow progress, 6-12 hours in nulliparas) and active (6-10 cm, faster progress, 1-2 cm/hour in nulliparas). Active phase begins at 6 cm (the Friedman curve was revised: the old 4 cm threshold was updated to 6 cm by ACOG/NICHD in 2014). Arrest of first stage: no cervical change for \(\ge\)4 hours with adequate contractions or \(\ge\)6 hours with inadequate contractions.

\medterm{First Trimester}-- The period from conception to 12 weeks gestational age. Characterized by rapid embryonic organogenesis, rising hCG levels, and establishment of the placenta. Common symptoms include nausea, vomiting, fatigue, and breast tenderness. Dating ultrasound performed to confirm gestational age. Risk of miscarriage highest during this trimester. Prenatal vitamins with folic acid recommended.

\medterm{Folic Acid Supplementation}-- Daily folic acid (400-800 mcg) taken before conception and during the first trimester to prevent neural tube defects (NTDs: spina bifida, anencephaly, encephalocele). Folic acid is a B vitamin (folate) essential for DNA synthesis and cell division. Higher doses (4-5 mg daily) are recommended for women with a prior NTD-affected pregnancy, diabetes, obesity, or those taking antiepileptic medications (valproate, carbamazepine). NTDs occur at 21-28 days post-conception, often before pregnancy is recognized, making preconceptional supplementation critical.

\medterm{Forceps-Assisted Delivery}-- Operative vaginal delivery using forceps (two curved blades applied to the fetal head). Types: outlet (forceps applied at the introitus, fetal scalp visible), low (station \(\ge\)+2, not at outlet), midpelvic (station \(\ge\)+2 but head above ischial spines), and rotational (forceps used to rotate from OP to OA). Most common types: Simpson (for molded heads), Elliot (for unmolded heads), and Piper (for aftercoming head in breech delivery). Prerequisites: full dilatation, ROM, engaged head, known position, adequate pelvis, and consent. Neonatal risks: facial nerve palsy, brachial plexus injury, bruising, and cephalohematoma. Maternal risks: third/fourth-degree lacerations (higher than vacuum).

\medterm{Forceps Delivery}-- Assisted vaginal delivery using forceps (curved blades applied to the fetal head) to expedite delivery for maternal or fetal indications (prolonged second stage, fetal distress, maternal exhaustion, medical conditions limiting pushing). Requires: full dilation, engaged vertex (station and position known), ruptured membranes, empty bladder, and adequate analgesia; contraindications: fetal malpresentation, suspected cephalopelvic disproportion. Complications: maternal (vaginal/perineal lacerations, episiotomy, incontinence) and fetal (facial bruising, cephalohematoma, nerve injury—rare serious injury). Compared with vacuum, forceps have higher maternal trauma but lower fetal scalp trauma. Increasingly replaced by cesarean for difficult cases.

\medterm{Fourth-Degree Tear}-- The most severe type of perineal tear during childbirth, extending through the anal sphincter (both internal and external) and into the rectal mucosa. Fourth-degree tears occur in approximately 0.5-1\% of vaginal deliveries and are associated with forceps delivery, shoulder dystocia, large baby, and episiotomy. Repair requires expert surgical reconstruction in the operating room: layered closure of rectal mucosa, internal anal sphincter, external anal sphincter, and perineal muscles. Complications: anal incontinence (flatus and fecal), rectovaginal fistula, and dyspareunia.

\medterm{Friedman Curve}-- The classic labor curve described by Emanuel Friedman in the 1950s, plotting cervical dilatation against time. Defined the latent and active phases, with active phase beginning at 4 cm. The Friedman curve was the standard for labor management for decades. Modern revisions (Consortium on Safe Labor, Zhang 2010) showed that labor progresses more slowly than Friedman described, particularly in nulliparas and with epidural anesthesia. The active phase now starts at 6 cm. The new paradigm recommends more patience before diagnosing arrest of labor.

\medterm{Fundal Height}-- The distance from the symphysis pubis to the top of the uterus measured in centimeters. Used to estimate gestational age and fetal growth. Generally correlates with weeks of gestation between 20-36 weeks (e.g., 24 cm at 24 weeks). Discrepancies may suggest fetal growth restriction, macrosomia, polyhydramnios, or multiple gestation. Measured with patient supine using a tape measure.

\medterm{Gestation}-- The period of development of a fetus from fertilization to birth. Human gestation averages 280 days (40 weeks) from the first day of the last menstrual period (LMP) or 266 days (38 weeks) from conception. Trimesters: first (weeks 1-12: organogenesis, highest vulnerability to teratogens), second (weeks 13-27: fetal growth and maturation, quickening at 18-22 weeks), third (weeks 28-40: rapid weight gain, lung maturation, acquisition of immune competence). Gestational age is assessed by: LMP, early ultrasound (crown-rump length at 11-13 weeks most accurate), symphysis-fundal height measurement after 20 weeks, and clinical assessment. Gestational age categories: preterm (<37 weeks), term (37-41 weeks), post-term (>41 weeks).

\medterm{Gestational Age}-- The age of pregnancy measured in weeks from the first day of the last menstrual period (LMP). Also called menstrual age. Typically 2 weeks longer than fertilization age. Calculated using LMP or firsttrimester ultrasound. Used to estimate due date (40 weeks from LMP). Preterm is less than 37 weeks; post-term is beyond 42 weeks. Key metric for fetal development milestones and management decisions.

\medterm{Gestational Diabetes}-- Glucose intolerance first diagnosed during pregnancy, affecting 6-9\% of pregnancies. Risk factors include obesity, advanced maternal age, family history, and prior gestational diabetes. Complications include macrosomia, birth trauma, neonatal hypoglycemia, and preeclampsia. Screening by 50g glucose challenge test at 24-28 weeks, with confirmatory 100g oral glucose tolerance test. Management includes dietary modification, glucose monitoring, and insulin therapy when diet fails.

\medterm{Gestational Hypertension}-- New-onset hypertension (BP \(\ge\)140/90 mmHg) after 20 weeks of gestation in a woman with previously normal blood pressure, without proteinuria or other features of preeclampsia. Affects 6-8\% of pregnancies. Management: monitoring blood pressure, screening for preeclampsia progression (proteinuria, labs, symptoms). Delivery at term (37-40 weeks) is typical. Some patients progress to preeclampsia. Treatment with antihypertensives (labetalol, nifedipine, methyldopa) is recommended for BP \(\ge\)160/110 mmHg.

\medterm{Gestational Sac}-- The first sonographic sign of an intrauterine pregnancy, visible on transvaginal ultrasound at approximately 4.5-5 weeks gestational age. A fluid-filled structure within the decidua, surrounded by the chorionic ring. Mean sac diameter (MSD) correlates with gestational age. The yolk sac is visible within the gestational sac at 5-6 weeks, and the embryo with cardiac activity at 6-7 weeks. Abnormal findings: empty gestational sac without yolk sac (pseudogestational sac in ectopic, anembryonic pregnancy).

\medterm{Gestational Trophoblastic Disease}-- A spectrum of conditions arising from abnormal trophoblast proliferation, including complete and partial hydatidiform moles and gestational choriocarcinoma. Complete mole has 46,XX karyotype (all paternal origin); partial mole is triploid (69 chromosomes). Presents with elevated hCG, uterus large for dates, and vaginal bleeding. Diagnosis by ultrasound (snowstorm pattern) and histology. Managed by uterine evacuation and hCG monitoring.

\medterm{Group B Streptococcus}-- A gram-positive bacterium (Streptococcus agalactiae) colonizing the genitourinary and gastrointestinal tracts. Asymptomatic vaginal colonization occurs in 20-30\% of pregnant women. Screening is performed at 36-37 weeks gestation via rectovaginal swab culture. Intrapartum antibiotic prophylaxis (IV penicillin G or ampicillin) is indicated for women with positive GBS screening, prior infant with GBS disease, or GBS bacteriuria during the current pregnancy. Penicillin-allergic women receive cefazolin (if low risk) or clindamycin/vancomycin (if high risk). Prophylaxis reduces early-onset neonatal GBS disease (sepsis, pneumonia, meningitis) by 80\%. Early-onset GBS disease occurs in the first 7 days of life and has a mortality rate of 5-10\% in term infants.

\medterm{Growth Scan}-- An ultrasound performed in the third trimester (usually 28-32 weeks and/or 36-38 weeks) to assess fetal growth. Measurements: BPD, HC, AC, FL, and estimated fetal weight (EFW). Growth percentiles plotted on a growth curve. Abnormal: <10th percentile (SGA/IUGR) or >90th percentile (LGA/macrosomia. Amniotic fluid volume and Doppler studies (umbilical artery, MCA) often included. Indications: suspected growth abnormality, multiple gestation, maternal diabetes, hypertension, prior IUGR, and postterm surveillance.

\medterm{Gynaecologist}-- A medical doctor who specializes in the health of the female reproductive system (vagina, cervix, uterus, fallopian tubes, ovaries) and breasts. Gynecologists diagnose and treat conditions including menstrual disorders, pelvic pain, endometriosis, fibroids, ovarian cysts, pelvic organ prolapse, incontinence, infections, and gynecologic cancers. They perform procedures including colposcopy, hysteroscopy, laparoscopy, and surgery. Many gynecologists also practice obstetrics (obstetrician-gynecologists, OB/GYNs). Subspecialties: gynecologic oncology, reproductive endocrinology and infertility, urogynecology, and pediatric/adolescent gynecology.

\medterm{Hand Prolapse}-- Prolapse of a fetal hand or arm alongside or below the presenting part. May occur spontaneously or after manual rotation. The prolapsed arm may be mistaken for a cord or foot. Management: if the hand is above the presenting part (compound presentation), it may be gently pushed back during contractions. If it presents beside the head after full dilatation, spontaneous delivery may occur, but rotation may be impeded. If vaginal delivery is obstructed, cesarean is performed.

\medterm{Head Circumference (HC)}-- The circumference of the fetal skull, measured at the same level as BPD (thalami, cavum septum pellucidum). Used for dating (second trimester) and for assessing head growth. HC is more reliable than BPD for dating when head shape is abnormal. Microcephaly: HC <3rd percentile (workup for infection, aneuploidy, genetic syndromes). Macrocephaly: HC >97th percentile.

\medterm{HELLP Syndrome}-- A severe variant of preeclampsia characterized by Hemolysis, Elevated Liver enzymes, and Low Platelets. Occurs in 10-20\% of patients with severe preeclampsia. Diagnostic criteria include haemolysis (LDH >600 U/L, schistocytes on peripheral smear), elevated liver enzymes (AST \(\ge\)70 U/L), and low platelets (<100,000/µL). HELLP can develop before or after delivery and may present without classic preeclampsia features (15-20\% of cases). Symptoms include right upper quadrant or epigastric pain, nausea, vomiting, headache, and malaise. Complications include DIC, placental abruption, hepatic rupture or infarction, acute kidney injury, pulmonary edema, and maternal death. Management: delivery is indicated regardless of gestational age. Corticosteroids (dexamethasone) may improve platelet count but are not recommended as routine treatment. Plasma exchange may be considered for persistent thrombocytopenia postpartum.

\medterm{HELPERR Mnemonic}-- A mnemonic for the sequential management of shoulder dystocia: H — Help (call for additional personnel, anesthesia, and pediatrics). E — Evaluate for episiotomy (generous episiotomy to provide more room). L — Legs (McRoberts maneuver: hyperflex maternal legs onto the abdomen). P — Suprapubic pressure (downward pressure on the maternal suprapubic area to dislodge the anterior shoulder). E — Enter internal rotation (Rubin II or Woods screw maneuver: rotate the fetal shoulders 180 degrees). R — Remove the posterior arm (sweep the posterior arm across the chest and deliver it). R — Roll the patient (Gaskin maneuver: onto all fours). If all maneuvers fail: Zavanelli maneuver (replace the fetal head and proceed with emergency cesarean).

\medterm{Hematologic Changes in Pregnancy}-- Plasma volume increases 40-50\% while red cell mass increases 20-30\%, resulting in physiologic anemia (Hb 10-11 g/dL). Leukocytosis (WBC up to 15,000-16,000/µL, predominantly neutrophils) is normal. The coagulation cascade shifts to a prothrombotic state: increased fibrinogen (factor I, 300-600 mg/dL), factors VII, VIII, IX, X, XII, and von Willebrand factor. Decreased protein S activity (free S decreases). The overall effect is an increased risk of venous thromboembolism (5-10 fold) and the utility of DVT prophylaxis in certain populations.

\medterm{hemorrhage}-- Excessive or severe bleeding. In obstetrics, hemorrhage is classified by timing: antepartum hemorrhage (bleeding after 24 weeks of pregnancy before birth, causes include placenta praevia, placental abruption, vasa praevia) and postpartum hemorrhage (bleeding >500 mL after vaginal delivery or >1000 mL after cesarean, causes include uterine atony, retained placenta, genital tract trauma, coagulopathy). Postpartum hemorrhage is a leading cause of maternal mortality worldwide. Management: uterine massage, oxytocin, ergometrine, carboprost, tranexamic acid, intrauterine balloon tamponade, uterine compression sutures (B-Lynch), and emergency hysterectomy.

\medterm{Herpetic Whitlow}-- A painful viral infection of the finger caused by the herpes simplex virus (typically HSV-1), characterized by swelling, redness, and painful vesicular lesions on the finger. In obstetrics, healthcare workers (midwives, obstetricians) are at risk when attending deliveries of women with genital herpes, as the virus can enter through breaks in the skin of the gloved hand. Treatment: oral antiviral medication (acyclovir, valacyclovir). Prevention: wearing double gloves during vaginal examinations and deliveries of women with active genital herpes lesions.

\medterm{Heterotopic Pregnancy}-- Simultaneous intrauterine and ectopic pregnancy. Rare in spontaneous conception (1:30,000) but more common with ART (1:100-1:500). Risk factors: ovulation induction, IVF with multiple embryo transfer. The presence of an intrauterine pregnancy does not exclude an ectopic pregnancy. Diagnosis: ultrasound showing intrauterine and extrauterine gestational sacs. Treatment: surgical removal of the ectopic pregnancy (salpingectomy) with preservation of the intrauterine pregnancy (which can proceed normally).

\medterm{High-Dependency Unit}-- A ward or area within a hospital that provides a level of care between a general ward and an intensive care unit, for patients who require intensive observation or treatment but not mechanical ventilation. In obstetrics, women may be admitted to HDU for: severe pre-eclampsia, postpartum hemorrhage, sepsis, or medical complications in pregnancy. HDU provides continuous monitoring (pulse oximetry, ECG, invasive blood pressure), one-to-one or one-to-two nursing, and the ability to administer vasoactive drugs or blood products.

\medterm{Hormone Replacement Therapy}-- HRT (menopause hormone therapy): estrogen ± progestin to relieve menopausal symptoms (vasomotor, vaginal atrophy, sleep, mood) and prevent osteoporosis. Regimens: estrogen-only (post-hysterectomy) or combined (with progestin for endometrial protection in uterus-intact women); routes: oral, transdermal, vaginal. Benefits: symptom relief, bone protection, and improved quality of life. Risks: VTE, stroke, breast cancer (estrogen+progestin; duration-dependent), and gallbladder disease. Selection: lowest effective dose, shortest duration, personalized (transdermal for VTE risk), timing (within 10 years of menopause for benefit). Contraindications: breast/endometrial cancer, VTE, liver disease, unexplained bleeding. Alternatives for vaginal/osteoporosis indications.

\medterm{Hot Flashes}-- Vasomotor symptoms: sudden sensations of intense heat, sweating, flushing, and palpitations, most commonly from estrogen withdrawal around menopause; also in hypogonadism, chemotherapy-induced ovarian failure, and some medications (SERMs, GnRH agonists). Mechanism: hypothalamic thermoregulatory dysfunction. Features: episodic (30 s-several min), often at night with disrupted sleep. Diagnosis: clinical; exclude secondary causes. Management: lifestyle (avoid triggers, cooling), non-hormonal (SSRIs/SNRIs—paroxetine, venlafaxine, gabapentin, clonidine, exercise), hormonal (HRT—estrogen for symptomatic menopause in appropriate candidates) as most effective; weigh risks (breast cancer, thrombosis); alternative/complementary options.

\medterm{Human Chorionic Gonadotropin}-- A glycoprotein hormone produced by trophoblast cells after implantation. Detectable in serum 8-11 days post-conception and in urine 12-14 days. Levels double approximately every 48-72 hours in early pregnancy, peaking at 8-11 weeks then plateauing. Used in pregnancy testing, ectopic pregnancy diagnosis (abnormal rise), and molar pregnancy monitoring. Beta-hCG subunit measured quantitatively for clinical decisions.

\medterm{Hydatidiform Mole}-- A type of gestational trophoblastic disease characterized by abnormal proliferation of trophoblastic tissue. Complete mole: 46,XX (all paternal), no fetal tissue, diffuse vesicular swelling; risk of malignant transformation (15-20\%). Partial mole: triploid (69,XXX), some fetal tissue, focal swelling; lower risk. Presents with vaginal bleeding in first trimester, excessive uterine size, theca lutein cysts, hyperemesis, and very high hCG. Diagnosis: ultrasound (snowstorm appearance, no fetal parts), markedly elevated hCG. Treatment: suction curettage with careful hCG monitoring until undetectable. Chemotherapy (methotrexate, actinomycin D) for persistent or malignant disease.

\medterm{Hydrops Fetalis}-- Abnormal accumulation of fluid in two or more fetal compartments, including pleural effusion, pericardial effusion, ascites, and subcutaneous edema. Causes include Rh isoimmunization, chromosomal abnormalities, cardiac defects, infections (parvovirus B19), and twin-to-twin transfusion syndrome. Often associated with severe fetal anemia or cardiac failure. Mortality rate high; management depends on etiology. May require intrauterine transfusion.

\medterm{Hymen}-- A thin fold of mucous membrane that partially covers the external opening of the vagina. The hymen is a remnant of fetal development, with no known physiological function in postnatal life. Shape and configuration vary widely: annular (most common), crescentic, cribriform (multiple small openings), septate, redundant, imperforate (complete closure, requiring surgical incision). The hymen typically has a central or crescentic opening of variable size that allows passage of menstrual blood. Rupture (tearing) of the hymen was traditionally associated with first vaginal intercourse (defloration), but the hymen can be stretched or torn by: tampon use, vigorous exercise, horseback/bicycle riding, gynaecological examination, or trauma. Its appearance does not reliably indicate virginity or sexual activity. Clinical significance: (1) Imperforate hymen: presents in adolescents with primary amenorrhoea and cyclic pelvic pain due to haematocolpos (accumulated menstrual blood in the vagina), treated by hymenotomy (cruciate incision). (2) Hymenal tags or polyps. (3) Hymenal tears in sexual abuse cases (forensic examination). (4) Hymenal stenosis or scarring. The concept of the hymen as a marker of virginity is a cultural and social construct, not a reliable anatomical indicator.

\medterm{Hyperemesis Gravidarum}-- Severe, persistent nausea and vomiting during pregnancy resulting in weight loss greater than 5\%, dehydration, electrolyte imbalances, and ketosis. Occurs in 0.3-2\% of pregnancies, typically presenting in the first trimester. Differentiated from normal morning sickness by severity and metabolic derangement. Management includes IV fluids, antiemetics (ondansetron, promethazine), thiamine supplementation to prevent Wernicke encephalopathy, and dietary modification. May require hospitalisation for IV hydration and electrolyte correction. Risk factors include female fetus, multiple gestation, molar pregnancy, and history of hyperemesis in a prior pregnancy. Complications include Wernicke encephalopathy (from thiamine deficiency), hypokalaemia, hyponatraemia, and vitamin deficiencies. Prognosis is generally good with symptoms improving by 16-20 weeks, though some women experience symptoms throughout pregnancy.

\medterm{Hyperprolactinaemia}-- Elevated levels of the hormone prolactin in the blood, above the normal range (typically >500 mIU/L in women). Causes: physiological (pregnancy, breastfeeding, stress, exercise), pharmacological (antipsychotics, metoclopramide, SSRIs, verapamil), and pathological (prolactinoma — pituitary adenoma, hypothyroidism, polycystic ovary syndrome, chronic renal failure). In women, hyperprolactinaemia causes galactorrhoea (milk production), oligomenorrhea or amenorrhea, anovulation, and infertility. Diagnosis: serum prolactin measurement (fasting, morning, avoid venipuncture stress), thyroid function tests, and pituitary MRI for macroprolactinoma. Treatment: dopamine agonists (cabergoline first-line due to better tolerability, bromocriptine second-line).

\medterm{Hysterectomy}-- The surgical removal of the uterus, the second most common major surgical procedure in women after cesarean section in the United States. The types of hysterectomy are classified as total (complete removal of the uterus and cervix, the most common), subtotal or supracervical (removal of the uterine body, leaving the cervix in place, preserving pelvic support and potentially reducing operative complications but requiring continued cervical cancer screening), radical (removal of the uterus, cervix, upper vagina, parametrial tissue, and pelvic lymph nodes, reserved for cervical and uterine malignancies), and total hysterectomy with bilateral salpingooophorectomy (removal of the uterus, cervix, fallopian tubes, and ovaries, inducing surgical menopause).

\medterm{Implantation}-- The embedding of the blastocyst into the endometrium, occurring approximately 6-10 days after ovulation. The blastocyst hatches from the zona pellucida, and the trophoblast invades the decidua. Implantation triggers hCG production, which maintains the corpus luteum. Failure of implantation is a common cause of early pregnancy loss. Ectopic implantation occurs outside the uterus, most commonly in the fallopian tube.

\medterm{Induction of Labor}-- The artificial initiation of uterine contractions before the onset of spontaneous labor. Indications: postterm pregnancy (>41-42 weeks), hypertensive disorders (preeclampsia, gestational hypertension), maternal medical conditions (diabetes, renal disease), fetal compromise (IUGR, oligohydramnios), PROM, chorioamnionitis, and elective reasons (shared decision-making after 39 weeks). Methods: cervical ripening agents (prostaglandins, mechanical dilation with Foley balloon), followed by oxytocin if needed. Contraindications: placenta previa, vasa previa, transverse lie, prior classical cesarean, active genital herpes, and absolute cephalopelvic disproportion.

\medterm{Influenza Vaccination in Pregnancy}-- Annual inactivated influenza vaccine recommended for all pregnant women. Pregnancy increases the risk of severe influenza (pneumonia, ICU admission, death) due to changes in immune function and cardiopulmonary physiology. Maternal vaccination also protects the newborn during the first 6 months of life (transplacental antibody transfer) and reduces the risk of influenza-like illness in the infant by 50-70\%. The vaccine is safe at any trimester.

\medterm{Intensive Care Unit}-- A specialist hospital unit providing continuous monitoring and life support for critically ill patients. In obstetrics, women may require ICU admission for: severe pre-eclampsia or eclampsia, massive postpartum hemorrhage, amniotic fluid embolism, sepsis, cardiac disease in pregnancy, or complications of anesthesia. ICU provides mechanical ventilation, continuous hemodynamic monitoring (invasive arterial pressure, central venous pressure), renal replacement therapy, and vasoactive drug support. Multidisciplinary care involves intensivists, obstetricians, and midwives.

\medterm{Intermittent Auscultation (IA)}-- The periodic auscultation of the fetal heart rate with a fetoscope or Doppler ultrasound during labor, without continuous electronic monitoring. Protocol: auscultate after a contraction for 60 seconds, during the active phase every 15-30 minutes in the first stage and every 5-15 minutes in the second stage. Used for low-risk pregnancies (no meconium, normal labor progress, no maternal risk factors). Advantages: less restricted mobility, lower intervention rate. Disadvantages: requires one-to-one nursing, may miss subtle changes. ACOG supports IA for low-risk patients.

\medterm{Internal Podalic Version}-- A maneuver where the operator inserts a hand into the uterus, grasps one or both fetal feet, and converts a transverse lie or cephalic presentation to a breech, then delivers by total breech extraction. Indication: delivery of the second twin (transverse lie) or severe fetal distress in the second stage. Rare in modern obstetrics. Contraindications: contracted pelvis, uterine scar, ruptured uterus, and fetal distress from other causes.

\medterm{Interstitial (Cornual) Pregnancy}-- A rare ectopic pregnancy (2-4\%) implanted in the intramural portion of the fallopian tube within the uterine wall. Higher morbidity than tubal ectopic because the myometrial layer expands, allowing later rupture (12-16 weeks) with massive hemorrhage. Risk factors: prior ipsilateral salpingectomy, IVF. Diagnosis: ultrasound (eccentric gestational sac surrounded by <5 mm myometrium), interstitial line sign. Treatment: cornual resection (wedge resection of the uterine horn) or hysterectomy for rupture. Methotrexate may be considered for early diagnosis with low hCG.

\medterm{Intimate Partner Violence (IPV) in Pregnancy}-- Physical, sexual, or psychological abuse by a current or former partner during pregnancy. Prevalence: 3-9\% of pregnancies. Risk factors: unplanned pregnancy, late prenatal care, substance use, and young age. IPV increases the risk of miscarriage, preterm birth, low birth weight, placental abruption, and maternal death. Screening: universal screening with validated tools (HITS, AAS, PRISMA). All positive screens require safety assessment, referral to domestic violence resources (hotline, shelter), and documentation. Mandatory reporting varies by state.

\medterm{Intrahepatic Cholestasis of Pregnancy (ICP)}-- A pregnancy-specific liver disorder characterized by pruritus (palms and soles, often nocturnal, with no rash) and elevated serum bile acids (>10-40 µmol/L). Incidence: 0.5-2\% of pregnancies. Risk factors: prior ICP, multiple gestation, hepatitis C, and high estrogen states. Onset: third trimester (after 28 weeks), resolves within 48 hours of delivery. Diagnosis: elevated total bile acids (fasting), elevated liver enzymes (AST, ALT, GGT). Fetal risks: preterm birth (spontaneous and iatrogenic), meconium-stained amniotic fluid, and stillbirth (with bile acids >40-100 µmol/L). Management: ursodeoxycholic acid (UDCA, 10-15 mg/kg/day) for pruritus, vitamin K supplementation (coagulopathy risk from fat malabsorption), and planned delivery at 36-39 weeks depending on bile acid level. Fetal surveillance advised.

\medterm{Intrapartum}-- The period during labor and childbirth, from the onset of regular uterine contractions leading to cervical dilation until delivery of the baby and placenta. Intrapartum care includes: monitoring maternal vital signs and fetal wellbeing (CTG), pain relief (epidural, gas and air, opioids), management of labor progress (partogram), assisted delivery if needed (forceps, ventouse), and active management of the third stage (oxytocin for placental delivery). Intrapartum complications include fetal distress, failure to progress, shoulder dystocia, and postpartum hemorrhage.

\medterm{Intrauterine Device}-- IUD: a long-acting reversible contraceptive (LARC) device placed in the uterus. Types: copper IUD (ParaGard—hormone-free, up to 10-12 years, also emergency contraception) and levonorgestrel-releasing IUS (Mirena/Skyla—progestin, up to 3-8 years, reduces bleeding). Mechanisms: copper—spermicidal/inflammatory; LNG—thickens cervical mucus, endometrial changes. Insertion by trained clinician; expulsion and perforation are rare. Benefits: highly effective, reversible, no systemic hormones (copper), reduced menstrual bleeding (LNG). Side effects: cramping, spotting, irregular bleeding; risk of PID is low. Contraindications: current PID, pregnancy, uterine anomaly, some. Used for contraception and for menorrhagia/endometriosis.

\medterm{Intrauterine Fetal Demise (IUFD) / Stillbirth}-- Fetal death before complete expulsion or extraction from the mother, at \(\ge\)20 weeks gestation. Affects 1:160 pregnancies in the US (~25,000 annually). Causes: placental abruption, umbilical cord accidents, infection, genetic abnormalities, maternal medical conditions (diabetes, hypertension, SLE), and unexplained (>25\%). Evaluation: fetal autopsy, placental pathology, karyotype/microarray, and maternal testing (Kleihauer-Betke, lupus anticoagulant, anticardiolipin, thrombophilia, HbA1c, infection screen, thyroid). Management: induction of labor or D\&E. Grief support and counseling for subsequent pregnancy.

\medterm{Intrauterine Growth Restriction}-- IUGR: a fetus with estimated weight below the 10th percentile (or growth velocity decline), usually with placental insufficiency or maternal/fetal factors. Causes: maternal (preeclampsia, hypertension, chronic disease, malnutrition, smoking, drugs), placental (insufficiency, abruption), and fetal (chromosomal, congenital infection, multiple pregnancy). Types: symmetric (early onset, affects head and body—anomaly/infection) vs. asymmetric (head-sparing, later onset—placental insufficiency). Complications: perinatal mortality, hypoxia, meconium aspiration, hypoglycemia, and long-term metabolic effects. Diagnosis: serial ultrasound (fundal height, estimated fetal weight, Dopplers—UA/MCA). Management: surveillance, Doppler monitoring, treat cause, and timed delivery for fetal compromise.

\medterm{Intrauterine Pressure Catheter (IUPC)}-- A fluid-filled catheter inserted into the uterine cavity alongside the fetus to measure the intensity and frequency of contractions in mm Hg (Montevideo units, MVU). Indications: inadequate labor progress despite oxytocin, need to assess contraction adequacy, determining if tachysystole is present, and when external tocodynamometer cannot provide adequate data. MVU: sum of peak contraction pressures (minus resting tone) over 10 minutes. Adequate labor: \(\ge\)200-250 MVU. Placed through the cervix after ROM, typically during active labor.

\medterm{Intrauterine Transfusion}-- Ultrasonically guided intravascular or intraperitoneal transfusion of fetal red blood cells to treat severe fetal anemia, most commonly from Rh isoimmunization. Performed through the umbilical vein using a 22-25 gauge needle. Packed red blood cells, CMV-negative and crossmatched, are transfused slowly. Risks include bradycardia, infection, preterm labor, and fetal death. May need repeated transfusions until delivery.

\medterm{In Vitro Fertilisation (IVF)}-- An assisted reproductive technology in which eggs are retrieved from the ovaries, fertilized with sperm in a laboratory, and resulting embryos are transferred to the uterus. Stages: ovarian stimulation (FSH injections for 10-14 days), oocyte retrieval (transvaginal ultrasound-guided aspiration under sedation), insemination (conventional or ICSI), embryo culture (3-6 days), and embryo transfer (1-2 embryos placed in the uterine cavity, typically at blastocyst stage). Success rates: approximately 30-40\% live birth per cycle for women under 35, declining with age.

\medterm{Involution}-- The process by which the uterus returns to its prepregnancy size after delivery. The uterus weighs approximately 1000 grams immediately postpartum and returns to approximately 60 grams by 6 weeks. Fundal height decreases approximately 1 cm per day.

\medterm{IVF}-- In vitro fertilization: the assisted reproductive technology in which oocytes are retrieved, fertilized by sperm outside the body (in vitro), and the resulting embryos transferred to the uterus. Indications: tubal factor, male factor (with ICSI), endometriosis, unexplained, ovulation disorders, genetic testing, and same-sex parenthood. Steps: controlled ovarian stimulation, ovulation trigger, oocyte retrieval (transvaginal ultrasound-guided), fertilization, embryo culture (days 3-6, blastocyst), embryo transfer (fresh or frozen), and luteal phase support. Success: live birth per cycle depends on maternal age (40-50 percent <35, declining after). Risks: OHSS, multiple pregnancy, ectopic, and cost. PGT screens embryos.

\medterm{Kangaroo Care (Skin-to-Skin Contact)}-- Holding a newborn against the mother's bare chest, typically for at least 60 minutes, promoting physiologic stability, thermoregulation, bonding, and breastfeeding initiation. Especially beneficial for preterm and low-birth-weight infants: reduces mortality (by 40\%), hypothermia, and hospital-acquired infections. Improves maternal confidence and reduces postpartum depression. Now standard practice for all newborns, including term infants in the immediate postpartum period.

\medterm{Karyotyping}-- A laboratory technique that produces an image (karyotype) of an individual's complete set of chromosomes, arranged in pairs by size, centromere position, and banding pattern. Karyotyping is used to detect chromosomal abnormalities: aneuploidy (trisomy 21, 18, 13), sex chromosome abnormalities (Turner, Klinefelter), structural abnormalities (deletions, duplications, translocations), and marker chromosomes. In obstetrics, karyotyping is performed on samples obtained by chorionic villus sampling or amniocentesis, or on products of conception after miscarriage.

\medterm{Ketones}-- Acidic byproducts produced when the body breaks down fat for energy instead of glucose (ketogenesis). Ketones can be detected in blood and urine. In pregnancy, ketones may be elevated due to: hyperemesis gravidarum (prolonged vomiting and starvation), gestational diabetes or poorly controlled pre-existing diabetes (diabetic ketoacidosis), and prolonged labor with inadequate caloric intake. Testing for ketones (urine dipstick or blood ketone meter) is part of the assessment of women with vomiting, diabetes, or reduced oral intake in pregnancy. Significant ketonuria may indicate dehydration or metabolic decompensation requiring IV fluids and glucose.

\medterm{Labor}-- The process of uterine contractions leading to progressive cervical dilation and effacement and delivery of the fetus. Stages: first (latent—cervical dilation to 3-4 cm, then active—3-10 cm), second (full dilation to delivery of the infant), third (delivery of the placenta), fourth (first hour postpartum). Features of normal labor: regular contractions, cervical change, rupture of membranes. Onset: may begin spontaneously (with prelabor: engagement, bloody show, Braxton Hicks). Abnormal labor: prolonged/protracted (dystocia), requiring augmentation, operative delivery, or cesarean. Management: maternal/fetal monitoring (CTG), analgesia (nitrous, epidural), progress evaluation (partogram), and support of the second stage; immediate newborn and third-stage care.

\medterm{Labour}-- The process of childbirth, consisting of regular uterine contractions that cause progressive cervical effacement and dilation, leading to delivery of the baby and placenta. Three stages: First stage — from onset of labor to full cervical dilation (10 cm), divided into latent (slow dilation to 4-6 cm) and active (rapid dilation from 4-6 cm to 10 cm) phases. Second stage — from full dilation to delivery of the baby, with passive descent and active pushing. Third stage — delivery of the placenta and membranes, active management (oxytocin + controlled cord traction) recommended. The partogram records labor progress and alerts to abnormalities.

\medterm{Lactation (Breastfeeding)}-- The production of milk by the mammary glands stimulated by the hormone prolactin. Colostrum is produced for the first 2-5 days (rich in IgA, proteins, and growth factors). Mature milk appears at day 5-7. The milk ejection reflex (let-down) is mediated by oxytocin. Benefits to the baby: optimal nutrition, passive immunity (IgA, lactoferrin, lysozyme), reduced risk of infections (otitis, gastroenteritis, pneumonia), SIDS, and long-term protection against obesity and diabetes. Benefits to mother: reduced risk of breast and ovarian cancer, postpartum weight loss, and uterine involution from oxytocin release.

\medterm{Laparoscopy}-- A minimally invasive surgical technique using a laparoscope (a thin, lighted telescope) inserted through a small incision (typically at the umbilicus) to visualize the abdominal and pelvic organs. Carbon dioxide gas is insufflated to create a working space. Additional small incisions (ports) allow insertion of instruments. In gynecology, laparoscopy is used for diagnosis and treatment of: endometriosis, ovarian cysts, ectopic pregnancy, pelvic adhesions, and tubal infertility. Also used for sterilization (tubal occlusion) and hysterectomy.

\medterm{Laparotomy}-- A surgical incision into the abdominal cavity, typically a midline vertical or transverse (Pfannenstiel) incision, providing wide access to the abdominal and pelvic organs. In obstetrics and gynecology, laparotomy is performed for: emergency cesarean delivery, ectopic pregnancy (ruptured), ovarian tumor (suspected malignant), pelvic abscess, hysterectomy for large fibroids or cancer, and repair of uterine rupture. Laparotomy carries greater morbidity than laparoscopy: longer recovery, more pain, higher infection risk, and increased adhesion formation.

\medterm{Last Menstrual Period (LMP)}-- The first day of the last normal menstrual period, used to estimate gestational age and the due date. Requires regular menstrual cycles and reliable recall. Limitations: irregular cycles, recent oral contraceptive use, early pregnancy bleeding mistaken for menses, and poor recall. Discriminatory zone: if ultrasound dating differs from LMP dating by more than 7 days in the first trimester, 14 days in the second trimester, or 21 days in the third trimester, ultrasound dating is preferred.

\medterm{Late Decelerations}-- A fetal heart rate pattern with a gradual decrease (onset to nadir >30 seconds) that starts at or after the peak of the contraction and returns to baseline after the contraction ends. Associated with uteroplacental insufficiency (reduced oxygen transfer across the placenta). Recurrent late decelerations (>50\% of contractions) with moderate variability require intrauterine resuscitation. Late decelerations with absent variability are Category III and indicate potential fetal acidosis. Management: maternal oxygen, position change (left lateral), IV fluids, reduce oxytocin.

\medterm{Latent Phase of Labor}-- The early part of the first stage of labor, from the onset of regular contractions to 6 cm cervical dilatation. Characterized by slow, progressive cervical change (usually <1 cm/hour). Duration is highly variable: up to 20 hours in nulliparas and 14 hours in multiparas. Active-phase labor arrest cannot be diagnosed during the latent phase. Prolonged latent phase (>20 hours nullipara, >14 hours multipara) is managed with rest, hydration, and observation or oxytocin augmentation.

\medterm{Late Preterm (34-36 weeks)}-- Preterm births between 34 weeks 0 days and 36 weeks 6 days, accounting for ~70\% of all preterm births. Late preterm infants are at increased risk compared to term infants: respiratory distress (TTN, RDS), temperature instability, hypoglycemia, jaundice, feeding difficulties, and hospital readmission. The risk of adverse outcomes is inversely related to gestational age. Antenatal corticosteroids (betamethasone) are recommended for all preterm births <37 weeks within 7 days of expected delivery.

\medterm{Laxatives}-- Medications used to treat constipation by stimulating bowel movements or softening stool. In obstetrics, constipation is common in pregnancy (progesterone slows bowel motility, iron supplements, and later mechanical compression) and postpartum (perineal pain, fear of damaging sutures, opioid pain relief). Types: bulk-forming (psyllium, Fybogel) — increase stool water content; osmotic (lactulose, polyethylene glycol/macrogol) — draw water into the bowel; stimulant (senna, bisacodyl) — increase peristalsis; and stool softeners (docusate). Safe in pregnancy and breastfeeding.

\medterm{Leiomyoma}-- Uterine fibroid: a benign smooth muscle tumor of the uterus (estrogen/progesterone-sensitive), the most common pelvic tumor in women. Locations: subserosal, intramural, submucosal (most symptomatic). Features: often asymptomatic; abnormal uterine bleeding (menorrhagia), pelvic pain/pressure, bulk symptoms (urinary frequency, constipation), infertility, pregnancy complications (pain, abruption, malpresentation). Diagnosis: pelvic exam, ultrasound (gold standard imaging), MRI. Management: asymptomatic—observation; symptomatic—medical (tranexamic acid, NSAIDs, hormonal—progestins, GnRH agonists, LNG-IUS) or surgical (myomectomy, hysterectomy), UAE (embolization), and MR-guided focused ultrasound; sarcoma risk is low.

\medterm{Leopold Maneuvers}-- A systematic abdominal palpation technique used to determine fetal position and presentation. Four maneuvers: (1) fundal grip identifies fetal lie (head vs breech), (2) umbilical grip locates the fetal back, (3) Pawlik grip identifies the presenting part and its position, (4) pelvic grip confirms presentation and attitude. Performed after 36 weeks with an empty bladder. Helps determine need for ultrasound confirmation.

\medterm{Libido}-- Sexual desire or drive, influenced by biological, psychological, and social factors. In women, libido is affected by hormonal changes (menstrual cycle, pregnancy, menopause), medical conditions (depression, thyroid disorders, PCOS), medications (SSRIs, hormonal contraceptives), and life circumstances (stress, relationship issues, fatigue). Low libido (hypoactive sexual desire disorder) is the most common female sexual complaint. Management: address underlying causes, relationship counseling, psychosexual therapy, and consider testosterone therapy in selected cases (postmenopausal women).

\medterm{Linea Nigra}-- A dark brown or black vertical line appearing on the abdomen during pregnancy, extending from the pubic symphysis to the umbilicus or beyond. Caused by hormonal stimulation of melanocytes, resulting in hyperpigmentation. More prominent in darker-skinned women. Usually appears in the second trimester and fades weeks to months after delivery. One of many normal dermatologic changes of pregnancy along with melasma and striae gravidarum.

\medterm{Listeriosis in Pregnancy}-- Infection with Listeria monocytogenes, a foodborne Gram-positive rod. Pregnant women are 10-20 times more susceptible than the general population. Ingestion of contaminated food (deli meats, soft cheeses, unpasteurized dairy, raw sprouts). Maternal symptoms: fever, myalgias, GI symptoms (nausea, diarrhea). Fetal/neonatal infection: chorioamnionitis, preterm labor, miscarriage, stillbirth, and neonatal sepsis/granulomatosis infantiseptica. Diagnosis: blood culture, placental culture. Treatment: IV ampicillin (high dose) + gentamicin. Prevention: avoid high-risk foods during pregnancy.

\medterm{Local anesthetic}-- A medication that causes reversible loss of sensation in a specific area of the body without affecting consciousness. Local anesthetics work by blocking voltage-gated sodium channels in nerve cell membranes, preventing the propagation of action potentials and transmission of pain signals. Common agents: lidocaine (1-2\%, most widely used), bupivacaine (longer-acting), and ropivacaine. Onset: 2-5 minutes (infiltration). Duration: 30 minutes to several hours depending on agent, dose, and use of vasoconstrictor (epinephrine prolongs action and reduces bleeding). In obstetrics: used for perineal infiltration before episiotomy repair, pudendal block, and as part of epidural/spinal anesthesia.

\medterm{Lochia}-- Vaginal discharge following delivery that undergoes characteristic changes over the puerperium. Lochia rubra (days 1-4): red, bloody discharge containing blood, decidual tissue, and fetal membranes. Lochia serosa (days 5-10): pinkish-brown discharge with serous fluid, leukocytes, and organisms. Lochia alba (days 11-42): yellowish-white discharge containing mucus, leukocytes, and bacteria. The sequence reflects progressive uterine involution and healing of the placental site. Foul-smelling lochia suggests endometritis or retained products of conception. Prolonged or heavy lochia warrants evaluation for retained tissue or subinvolution. Normal total volume is approximately 200-500 mL over the first 2 weeks.

\medterm{Macrosomia}-- Birth weight exceeding 4000 grams (8 pounds 13 ounces) regardless of gestational age. Risk factors include maternal diabetes, obesity, post-term pregnancy, and prior macrosomia. Complications include shoulder dystocia, brachial plexus injury, clavicular fracture, and neonatal hypoglycemia. Cesarean delivery may be considered for estimated fetal weight over 4500 grams in diabetic patients or over 5000 grams in non-diabetic patients.

\medterm{Magnesium Sulfate for Neuroprotection}-- Intravenous magnesium sulfate administered to women in imminent preterm labor (<32 weeks) to reduce the risk of cerebral palsy and major motor dysfunction in the child. Dose: 4-6 g IV bolus, followed by 1-2 g/hour maintenance for up to 24 hours or until delivery. Mechanism: thought to stabilize cerebral blood flow, prevent excitotoxic injury, and reduce inflammation. The BEAM trial showed a 30-45\% reduction in the risk of moderate-to-severe cerebral palsy. Side effects: flushing, warmth, nausea, dizziness, and respiratory depression (monitor deep tendon reflexes and respiratory rate). Contraindication: myasthenia gravis.

\medterm{Malpresentation}-- Any fetal presentation other than vertex (cephalic) at term. Includes breech (buttocks or feet first), face, brow, and shoulder presentations. Breech presentation occurs in 3-4\% of term pregnancies and increases the risk of cord prolapse and difficult delivery. External cephalic version may be attempted at 36-38 weeks. Persistent breech presentation typically managed with cesarean delivery.

\medterm{Manual Removal of Placenta}-- The manual extraction of a retained placenta from the uterine cavity. Indication: retained placenta beyond 30-60 minutes. Technique: under adequate anesthesia (regional or general), one hand inserts into the uterine cavity with fingers together (like a knife hand), and the placenta is gently separated from the uterine wall using a sweeping motion. The other hand is placed on the abdomen to stabilize the uterus. After removal, the uterus is explored to ensure complete evacuation. Oxytocin is given. Complications: hemorrhage, infection, perforation, and Asherman syndrome.

\medterm{Mastitis}-- Inflammation of the breast tissue, most commonly occurring during lactation (lactational mastitis, affecting 5-20 percent of breastfeeding women), typically within the first 6-12 weeks postpartum. Lactational mastitis results from milk stasis combined with bacterial infection, most frequently Staphylococcus aureus (including MRSA) entering through cracked or fissured nipples. Clinical presentation includes localized breast pain, erythema, warmth, swelling, induration, and systemic symptoms of fever, chills, and malaise. Diagnosis is primarily clinical. Treatment includes milk removal through continued breastfeeding or pumping (emptying the affected breast is essential), supportive measures (cold compresses between feedings, warm compresses before feedings, NSAIDs for pain), and antibiotics if symptoms are severe or persist beyond 12-24 hours of conservative management (dicloxacillin or cephalexin for 10-14 days; clindamycin for penicillin-allergic or MRSA-suspected cases). The mother should continue breastfeeding or pumping to maintain milk flow, as abrupt cessation worsens the condition and increases the risk of breast abscess. Preventive measures include proper latch technique, frequent feeding, and hand hygiene. Recurrent mastitis may require investigation for underlying ductal abnormalities or milk supply issues.

\medterm{Maternal Physiologic Changes in Pregnancy}-- Extensive adaptive changes across all organ systems to support the growing fetus. Key changes: plasma volume increases 40-50\% (hemodilution), cardiac output increases 30-50\%, heart rate increases 10-15 bpm, systemic vascular resistance decreases, tidal volume and minute ventilation increase (chronic respiratory alkalosis), renal blood flow increases 50\% (increased GFR, decreased creatinine and BUN), and insulin resistance increases (mediated by human placental lactogen, cortisol, and growth hormone). These changes impact medication dosing, interpretation of lab values, and response to illness.

\medterm{Maternal Pushing (Second Stage)}-- Maternal expulsive efforts during the second stage of labor to facilitate fetal descent. Directed (closed-glottis, Valsalva) pushing: sustained breath-hold for 10 seconds with three pushes per contraction. Spontaneous (open-glottis) pushing: shorter pushes with the mother pushing when she feels an urge. Evidence: spontaneous pushing results in less perineal trauma, less fatigue, and better fetal acid-base status. Directed pushing may slightly shorten the second stage but increases perineal injury and is associated with lower umbilical cord pH.

\medterm{MCA Peak Systolic Velocity (PSV)}-- Doppler measurement of the peak systolic velocity in the fetal middle cerebral artery. An indicator of fetal anemia: elevated MCA PSV (>1.5 MoM) correlates with moderate-to-severe anemia. The most important application is surveillance of Rh-sensitized pregnancies (detect anemia before hydrops develops). Also used to assess fetal anemia due to parvovirus B19 infection or fetomaternal hemorrhage. MCA PSV is also used as part of the cerebroplacental ratio (CPR = MCA PI / umbilical artery PI) to assess brain sparing in IUGR.

\medterm{McDonald Cerclage}-- A transvaginal cervical cerclage using a purse-string suture (Mersilene tape or Prolene) placed at the cervicovaginal junction, without bladder dissection. The most common cerclage technique. Indications: cervical insufficiency (history-indicated or ultrasound-indicated short cervix <25 mm before 24 weeks). Placed at 12-14 weeks. The suture is removed at 36-37 weeks or at the onset of labor. Suture around the cervix, not through it. No flap of the bladder or rectum is created. Success rate: 80-90\% for term delivery in selected patients.

\medterm{Meconium Aspiration}-- Meconium aspiration syndrome (MAS): aspiration of meconium-stained amniotic fluid into the fetal/neonatal airway causing airway obstruction, chemical pneumonitis, and surfactant dysfunction—respiratory distress in term/post-term infants. Risk: fetal distress/hypoxia (passage of meconium in utero). Features: meconium-stained skin/nails, respiratory distress at birth, tachypnea, cyanosis, hyperinflation on X-ray. Management: at birth, if depressed—suction of oropharynx (routine suction not recommended for vigorous infants); support ventilation, oxygen, surfactant, and possibly ECMO for severe; complications: pneumothorax, persistent pulmonary hypertension, death. Prevention: fetal monitoring and timely intervention for distress.

\medterm{Melasma}-- A common hyperpigmentation disorder characterized by brown or gray-brown patches on the face, particularly the cheeks, forehead, nose, and upper lip. Also called chloasma or the mask of pregnancy.

\medterm{Menstrual Cycle}-- The monthly hormonal and uterine changes preparing for pregnancy, typically 21-35 days. Phases: follicular (FSH stimulates follicle growth and estradiol; endometrium proliferates), ovulation (LH surge mid-cycle, ~day 14, releases the oocyte), luteal (progesterone from corpus luteum supports endometrium; if no pregnancy, corpus luteum regresses), and menstruation (shedding). Hormonal control: hypothalamus GnRH $\rightarrow$ pituitary FSH/LH $\rightarrow$ ovary estradiol/progesterone $\rightarrow$ feedback. Abnormalities: anovulation, cycle irregularity (PCOS, thyroid, perimenopause), and bleeding disorders. Tracking (calendar, basal body temperature, LH kits) assists fertility or contraception; hormonal methods modify the cycle for contraception.

\medterm{Menstruation}-- Menses: the cyclic shedding of the endometrial lining (menstrual period) as part of the menstrual cycle, occurring with the fall in estrogen/progesterone when implantation does not occur. Normal: interval 21-35 days, duration 2-7 days, blood loss 30-80 mL. First menses = menarche (~12-13 years); cessation = menopause (~51 years). Abnormalities: amenorrhea (absence), oligomenorrhea, menorrhagia (heavy), metrorrhagia (irregular), dysmenorrhea (painful), and polymenorrhea. Physiology: regulated by the HPO axis (GnRH, FSH/LH, ovarian steroids). Evaluation for abnormalities: pregnancy, hormonal, structural (fibroids, polyps), and systemic causes. Hygiene, pain management, and cycle tracking support.

\medterm{Midtrimester}-- The middle third of pregnancy, from 14 to 28 weeks gestation. The midtrimester is when many women feel their best: nausea and fatigue of the first trimester resolve, and the physical discomforts of the third trimester have not yet developed. Key events: fetal movements felt (quickening, around 18-20 weeks), anomaly scan (18-21 weeks), maternal serum screening, and the fetus grows rapidly. By 24 weeks, the fetus reaches viability. Midtrimester complications include miscarriage (late), cervical insufficiency, and preterm labor.

\medterm{Miscarriage}-- Also known as spontaneous abortion, the spontaneous loss of a pregnancy before 20 weeks of gestation, occurring in 15-20 percent of clinically recognized pregnancies. The majority (80 percent) occur in the first trimester, with risk decreasing significantly after visualization of fetal cardiac activity. Chromosomal abnormalities account for 50-60 percent of first-trimester miscarriages, most commonly autosomal trisomies (trisomy 16, trisomy 22, trisomy 21, trisomy 18, trisomy 13), monosomy X, and triploidy. Risk increases with advancing maternal age: 10 percent at age 20, 15-20 percent at age 35, 25-40 percent at age 40, and greater than 50 percent at age 45. Other causes include uterine anatomic abnormalities (septate uterus, fibroids, adhesions), maternal disorders (diabetes, thyroid disease, antiphospholipid syndrome, thrombophilias), infections (Listeria, Toxoplasma, rubella, parvovirus B19), cervical insufficiency, and environmental factors (smoking, alcohol, high caffeine). Diagnosis is classified as threatened (bleeding, closed cervix, viable pregnancy), inevitable (open cervix with bleeding), incomplete (partial tissue passage), complete (empty uterine cavity), or missed (fetal demise without expulsion). Management is expectant, medical (misoprostol), or surgical (suction evacuation) depending on clinical presentation and patient preference.

\medterm{Misoprostol (Cytotec)}-- A synthetic prostaglandin E1 analog used off-label for cervical ripening and labor induction. Dose: 25-50 mcg vaginally or 25-100 mcg orally q4-6h. Advantages: inexpensive, stable at room temperature, and multiple routes of administration. Risks: uterine tachysystole (with and without FHR changes), the highest rate among induction agents. Contraindicated in patients with prior cesarean or uterine scar (increased risk of uterine rupture). The oral route has less tachysystole than vaginal. The buccal/sublingual route is also used.

\medterm{Modified Biophysical Profile (Modified BPP)}-- A simplified antepartum test combining the non-stress test (NST) with the amniotic fluid index (AFI). The NST assesses acute fetal status (heart rate reactivity over 20-40 minutes) and the AFI assesses chronic uteroplacental function (oligohydramnios suggests chronic hypoxia or placental insufficiency). Scoring: NST (reactive: 2, nonreactive: 0), AFI (>5 cm: 2, \(\le\)5 cm: 0). The total modified BPP score is 4 (normal), 2 (equivocal, repeat in 24 hours), or 0 (abnormal, consider delivery). Equivalent to the full BPP for outcome prediction in most settings.

\medterm{Molar Pregnancy}-- See Hydatidiform Mole.

\medterm{Montevideo Units (MVU)}-- A quantitative measure of uterine contraction intensity, calculated by summing the peak contraction pressures (in mm Hg) minus the resting tone for each contraction in a 10-minute window. Adequate labor: \(\ge\)200 MVU (spontaneous) or \(\ge\)250 MVU (augmented). Measured with an intrauterine pressure catheter (IUPC). Low MVU may prompt oxytocin augmentation. The concept helps distinguish true arrest of labor from inadequate contractions.

\medterm{Montgomery Tubercles}-- Small, raised sebaceous glands visible on the areola of the breast during pregnancy and lactation. They appear as small bumps surrounding the nipple and secrete an oily substance that lubricates the nipple and areola, protecting against infection. Their prominence increases during pregnancy under hormonal influence. Often mistaken for pathological lesions but are normal anatomic structures. First described by William Fetherston Montgomery in 1837.

\medterm{Morning Sickness}-- Nausea and vomiting of pregnancy (NVP): nausea (and sometimes vomiting) commonly in the first trimester (6-12 weeks), typically resolving by 14-16 weeks; multifactorial (hCG, estrogen, progesterone, odor sensitivity). Management: dietary changes (small frequent meals, ginger, bland foods, avoid triggers), pyridoxine (vitamin B6) ± doxylamine (first-line pharmacotherapy), and reassurance. Severe/persistent vomiting with weight loss, dehydration, and electrolyte disturbance = hyperemesis gravidarum—requires IV fluids, antiemetics (ondansetron, metoclopramide), electrolyte replacement, thiamine, and monitoring; exclude other causes (multiple pregnancy, molar pregnancy, GI/hepatic disease, thyroid).

\medterm{Multiple Pregnancy}-- Gestation with more than one fetus: twins (most common), triplets, higher order. Types: dizygotic (two eggs, two placentas—fraternal, more common, familial and increased with IVF) vs. monozygotic (one egg split—identical, sharing determined by timing: dichorionic-diamniotic, monochorionic-diamniotic, monoamniotic, conjoined). Risks: preterm birth, low birth weight, preeclampsia, gestational diabetes, twin-twin transfusion syndrome (monochorionic), IUGR, congenital anomalies, and fetal loss. Diagnosis: ultrasound (chorionicity critical). Management: high-risk antenatal care, serial growth scans, prevention/treatment of preterm labor, planned delivery (vaginal for low-risk vertex twins; cesarean for others), and postpartum surveillance.

\medterm{Naegele Rule}-- A formula for estimating the due date (estimated date of delivery, EDD): add 7 days to the first day of the last menstrual period (LMP), then add 9 months (or subtract 3 months). Example: LMP January 1 yields EDD October 8. Assumes a 28-day cycle with ovulation on day 14. The rule is most accurate with regular cycles and confirmed LMP. First-trimester ultrasound dating is more accurate and preferred when LMP is uncertain or cycles are irregular.

\medterm{Neonatal Resuscitation}-- The systematic approach to stabilizing a newborn after delivery, guided by the Neonatal Resuscitation Program (NRP). Steps: (1) initial steps (warm, dry, stimulate, position, clear airway if needed), (2) assess breathing, heart rate, and color, (3) positive pressure ventilation (PPV) with a bag-mask if gasping or HR <100 bpm, (4) chest compressions (3:1 ratio with ventilations) if HR <60 bpm despite PPV, (5) endotracheal intubation if needed, (6) epinephrine (IV or ET) if HR remains <60 bpm, and (7) volume expansion if suspected hypovolemia. The transition from fetal to neonatal circulation requires successful lung aeration and closure of fetal shunts.

\medterm{Nitrous Oxide for Labor}-- A 50:50 mixture of nitrous oxide and oxygen (Entonox) inhaled by the mother during labor for analgesia. Self-administered via a mask or mouthpiece; the woman controls her own dosing, taking a breath 30 seconds before a contraction. Advantages: noninvasive, rapid onset and offset, no effect on the baby at low concentrations, and the mother can ambulate. Disadvantages: incomplete pain relief (does not eliminate pain, only modulates it), can cause nausea, dizziness, and drowsiness. Contraindications: pneumothorax, bowel obstruction, and vitamin B12 deficiency. Widely used in the UK and Canada; increasingly available in US hospitals.

\medterm{Noninvasive Prenatal Testing (NIPT)}-- See Cell-Free Fetal DNA (cfDNA) Screening.

\medterm{Non-Stress Test}-- A fetal heart rate monitoring technique assessing fetal heart rate reactivity over 20-40 minutes. Normal (reactive) tracing shows at least 2 fetal heart rate accelerations of 15 beats per minute above baseline lasting at least 15 seconds within a 20-minute period. Non-reactive tracings may indicate fetal hypoxia or sleep cycle. Used for antepartum surveillance in high-risk pregnancies, including post-dates pregnancy and maternal conditions affecting fetal well-being.

\medterm{Nuchal Translucency (NT) Screening}-- An ultrasound measurement of the fluid-filled space at the posterior fetal nuchal region, performed at 11-14 weeks (crown-rump length 45-84 mm). Increased NT (>95th percentile or >3.0-3.5 mm) is associated with trisomies 21, 18, 13, Turner syndrome, and cardiac defects. Combined with maternal serum biomarkers (PAPP-A, free $\beta$-hCG) in the first-trimester combined screen, it detects ~85-90\% of Down syndrome. NT measurement requires certified training and quality assurance (Fetal Medicine Foundation).

\medterm{Nutrition in Pregnancy}-- A balanced diet providing adequate calories, protein, vitamins, and minerals to support maternal and fetal health. Additional caloric needs: ~300 kcal/day in the second and third trimesters. Key nutrients: folic acid (400 mcg/day), iron (27 mg/day), calcium (1000 mg/day), vitamin D (600 IU/day), iodine (220 mcg/day), DHA (200-300 mg/day). Foods to avoid: raw/undercooked meat and fish (listeriosis, toxoplasmosis), high-mercury fish (shark, swordfish, king mackerel), unpasteurized dairy, and excess caffeine (<200 mg/day).

\medterm{Obstetric Anal Sphincter Injury (OASIS)}-- Third- and fourth-degree perineal lacerations involving the anal sphincter complex. Incidence: 0.5-5\% of vaginal deliveries. Risk factors: midline episiotomy (increases risk 2-3 fold), operative vaginal delivery, primiparity, macrosomia, and occiput posterior position. Repair: immediate repair in the delivery room under adequate anesthesia (regional or general). End-to-end or overlap repair of the EAS (3-0 PDS), repair of the IAS if involved. Postoperative: broad-spectrum antibiotics, stool softeners (lactulose, docusate), and avoid constipation. Follow-up: pelvic floor physiotherapy, ultrasound, or MRI for persistent symptoms.

\medterm{Obstetrician}-- A medical doctor who specializes in the care of women during pregnancy, labor, and the postpartum period (obstetrics). Obstetricians manage normal and high-risk pregnancies, perform cesarean deliveries and assisted vaginal deliveries, and treat pregnancy complications (pre-eclampsia, gestational diabetes, hemorrhage, preterm labor). In the UK, obstetrics and gynecology are combined as a single specialty (O\&G). Training: 2 years foundation training + 7 years specialty training to become a consultant obstetrician and gynecologist.

\medterm{Obstetrics}-- The branch of medicine and surgery concerned with the care of women during pregnancy, childbirth, and the postpartum period. Obstetricians manage normal and high-risk pregnancies, perform antenatal screening and diagnosis, conduct deliveries (spontaneous vaginal, assisted vaginal—forceps, ventouse, and caesarean section), and manage obstetric emergencies (pre-eclampsia/eclampsia, hemorrhage, sepsis, shoulder dystocia, cord prolapse, uterine rupture). Antenatal care in the UK follows a schedule of 10 appointments for nulliparous women and 7 for multiparous, including ultrasound dating (10--13 weeks), anomaly scan (18--21 weeks), screening for aneuploidy, gestational diabetes, and infections. Maternal medicine deals with pre-existing medical conditions in pregnancy (diabetes, hypertension, cardiac disease, epilepsy). The specialty has close links with midwifery (midwife-led care for low-risk pregnancies), neonatology, and fetal medicine.

\medterm{Occiput Posterior (OP) Position}-- The fetal head presenting with the occiput oriented toward the maternal sacrum (back-to-back position). Occurs in 15-20\% of labors, about half rotate spontaneously to OA by delivery. Persistent OP is associated with prolonged labor (especially the second stage), arrest of descent, severe back pain (back labor), and higher rate of operative vaginal delivery and cesarean. Management: maternal positioning (hands and knees, pelvic rocking), manual rotation (hands or forceps), and vacuum or forceps delivery if needed.

\medterm{Occiput Transverse (OT) Position}-- The fetal head presenting with the sagittal suture oriented transversely. Usually the normal descent position; most rotate to OA with contractions. When persistent, it can cause protracted descent and prolonged second stage. Management: maternal position changes, and if persistent, manual rotation to OA or direct OP extraction with forceps/vacuum. Deep transverse arrest: OT at the outlet, more common with an android pelvis and requires rotation before delivery.

\medterm{Oestrogen}-- The primary female sex hormone, produced mainly by the ovaries (granulosa cells of developing follicles) and in smaller amounts by adipose tissue, the adrenal glands, and the placenta during pregnancy. The three major forms: estradiol (E2, most potent, predominant in reproductive-age women), estrone (E1, predominant after menopause), and estriol (E3, produced by the placenta). Functions: development of secondary sexual characteristics, regulation of the menstrual cycle, endometrial proliferation, breast development, bone density maintenance, and cardiovascular health. Used in HRT, combined oral contraceptives, and treatment of anovulation.

\medterm{Oligohydramnios}-- Insufficient amniotic fluid volume, defined as amniotic fluid index less than 5 cm or single deepest pocket less than 2 cm. Causes include fetal renal anomalies, post-term pregnancy, premature rupture of membranes, maternal dehydration, and placental insufficiency. Complications include umbilical cord compression, fetal growth restriction, and pulmonary hypoplasia. Management includes amniotic fluid infusion, maternal hydration, and delivery when indicated.

\medterm{One-Hour Glucose Challenge Test (GCT)}-- A screening test for gestational diabetes mellitus (GDM) performed at 24-28 weeks. The patient drinks 50 g of glucose (regardless of fasting state), and plasma glucose is measured at 1 hour. Threshold: \(\ge\)130-140 mg/dL (7.2-7.8 mmol/L) is considered positive, warranting diagnostic 3-hour oral glucose tolerance test (OGTT). Sensitivity: 80-90\% depending on the threshold used. A negative test (glucose <130 mg/dL) excludes GDM. Earlier testing (first trimester) is performed in high-risk women (obesity, prior GDM, family history of diabetes).

\medterm{Operative Vaginal Delivery}-- Vaginal delivery assisted by vacuum or forceps. Indications: prolonged second stage, suspected fetal compromise, need to shorten the second stage for maternal benefit (exhaustion, cardiac disease, hypertensive crisis). Prerequisites: full dilatation, ruptured membranes, engaged head (station \(\ge\)+2), known fetal position, adequate pelvis, and consent. Vacuum: soft or rigid cup applied to the fetal occiput. Forceps: two blades applied to the fetal head (outlet, low, or mid-pelvic). Success rate: 80-95\%. Complications: maternal (perineal lacerations, anal sphincter injury, vaginal lacerations, hematoma) and neonatal (cephalohematoma, subgaleal hemorrhage with vacuum, facial nerve palsy with forceps).

\medterm{Oral Contraceptive}-- Birth control pill: combined oral contraceptive (COC—ethinylestradiol + progestin) or progestin-only pill (POP/minipill). Mechanisms: ovulation inhibition, cervical mucus thickening, endometrial change. Benefits: effective (~91-99 percent with typical/perfect use), cycle regulation, reduced dysmenorrhea/menorrhagia, acne, ovarian/endometrial cancer prevention, and PMDD. Risks: VTE (estrogen—higher with smoking, age >35, obesity, migraine with aura), hypertension, and rare hepatic effects. Contraindications: VTE history, thrombophilia, migraine with aura, estrogen-sensitive cancer, liver disease, and smokers >35. POP is safe during breastfeeding. Prescribed with counseling and follow-up; missed pills reduce efficacy (backup contraception).

\medterm{Oral Contraceptives}-- Hormonal medications taken orally to prevent pregnancy. Combined oral contraceptive pills (COCP) contain oestrogen (ethinylestradiol 20--35 \(\mu\)g) and a progestogen (levonorgestrel, norethisterone, drospirenone, desogestrel). Mechanisms: inhibit ovulation (suppression of LH surge), thicken cervical mucus (impedes sperm penetration), and alter endometrial receptivity. Perfect-use failure rate: 0.3\% per year; typical-use failure rate: 7--9\%. Benefits: cycle regulation, reduced dysmenorrhoea, decreased ovarian/endometrial cancer risk (protective effect persists >30 years after stopping), reduced acne, and improved bone density. Contraindications (WHO categories 3--4): age \(\ge\)35 and smoking \(\ge\)15 cigarettes/day, uncontrolled hypertension, migraine with aura, history of VTE, breast cancer, liver disease. Cardiovascular risks: increased VTE risk (3--9 per 10,000 woman-years vs 1--5 in non-users). Progestogen-only pill (POP, 'mini-pill') contains desogestrel 75 \(\mu\)g or norethisterone 350 \(\mu\)g; strict timing required (within 3-hour window); safe for use during breastfeeding and in women with contraindications to oestrogen.

\medterm{Ovarian Cyst}-- A fluid-filled sac within or on the ovary, most commonly functional (follicular or corpus luteum cysts) related to the menstrual cycle. Usually asymptomatic; can present with pelvic pain, pressure, and bloating. Complications: rupture, hemorrhage, torsion (sudden severe pain with nausea/vomiting). Diagnosis: transvaginal ultrasound (assess size, morphology, septations, solid components, and vascularity). Simple cysts <5 cm in reproductive-age women are usually benign. Treatment: observation for simple cysts; surgical removal for complex/malignant features or persistent large cysts.

\medterm{Ovarian Pregnancy}-- A rare ectopic pregnancy (0.5-3\%) implanted on the ovary. Diagnosis: Spiegelberg criteria (ipsilateral fallopian tube intact, gestational sac on the ovary, connected to the ovary by the uterovarian ligament, ovarian tissue in the sac wall). Treatment: ovarian wedge resection or oophorectomy.

\medterm{Ovarian Torsion}-- Twisting of the ovary (and often the fallopian tube) on its vascular pedicle, causing obstruction of blood flow—a gynecologic emergency. Risk: ovarian cysts/tumors (especially >5 cm), longer utero-ovarian ligament, younger women, pregnancy. Features: sudden severe unilateral lower abdominal/pelvic pain, nausea/vomiting, and a palpable adnexal mass; may be intermittent (partial torsion) or with fever/infection. Diagnosis: clinical suspicion; ultrasound with Doppler (absent/reduced arterial or venous flow, enlarged ovary), CT/MRI. Management: urgent surgery (laparoscopy/laparotomy) with detorsion (ovarian preservation) ± oophoropexy; necrosis requires oophorectomy. Delayed treatment risks ovarian loss and infertility.

\medterm{Oxytocics}-- A class of medications that stimulate uterine contractions, used for induction and augmentation of labor and prevention/treatment of postpartum hemorrhage. Oxytocin (Syntocinon) — synthetic analog of endogenous oxytocin, administered IV infusion for labor induction or IM/IV for PPH prevention. Ergometrine — causes sustained tetanic uterine contraction, used IM for PPH, contraindicated in hypertension and preeclampsia. Carboprost (Hemabate) — prostaglandin F2-alpha analog, IM for refractory PPH. Misoprostol (Cytotec) — PGE1 analog, oral/sublingual/rectal for PPH prevention and treatment in resource-limited settings.

\medterm{Oxytocin Infusion}-- Synthetic oxytocin (Pitocin) administered intravenously for induction or augmentation of labor. Standard dosing: start at 1-2 mU/min, increase by 1-2 mU/min every 15-30 minutes to achieve 200-250 Montevideo units (MVU, measured by IUPC). Maximum dose: 40 mU/min. After achieving adequate labor, the dose may be maintained or decreased. Complications: uterine tachysystole (>5 contractions/10 min, leading to fetal hypoxia), hyperstimulation (tachysystole with nonreassuring FHR), water intoxication (antidiuretic effect at high doses), and uterine rupture (rare, increased risk with prior cesarean). Stop or reduce oxytocin if hyperstimulation occurs.

\medterm{Pap Smear}-- Papanicolaou test (cervical cytology): a screening test collecting cervical cells (conventional or liquid-based) to detect precancerous cervical changes (CIN) and cancer, usually combined with HPV DNA testing (co-testing). Screening: start at 21-25 years (age varies), every 3-5 years depending on modality and age; cytology alone or cytology + HPV. Abnormal results: ASC-US, LSIL, HSIL, AGC—managed by HPV triage, colposcopy, or repeat testing per guidelines; CIN2-3 treated (LEEP, cryotherapy, conization). HPV vaccination prevents most cervical cancers. Reduced incidence/mortality from screening programs. Follow-up and surveillance after abnormal results are essential.

\medterm{Partograph (Partogram)}-- A graphic record of cervical dilatation, fetal descent, and uterine contractions over time, used to monitor labor progress and identify labor abnormalities. The WHO partograph has an alert line (1 cm/hour dilatation) and an action line (4 hours after the alert line, indicating need for intervention). Components: cervical dilatation (plotted against time), fetal descent (station), contractions (frequency and duration), FHR, and maternal vital signs. The partograph reduces prolonged labor, augmentation rates, and intrapartum stillbirths in low-resource settings. Modern use: most centers now use the Zhang labor curve (active phase at 6 cm).

\medterm{Parturition (Childbirth)}-- The process of giving birth, encompassing three stages: (1) First stage: from the onset of regular uterine contractions (cervical effacement and dilation) to full dilation (10 cm). Latent phase—slow dilation (0--4 cm) over hours, variable duration. Active phase—more rapid dilation (4--10 cm), typically 1--1.5 cm/hour in nulliparous women. (2) Second stage: from full dilation to delivery of the baby. Passive phase—descent of the presenting part without active pushing. Active phase—maternal pushing with contractions. Duration: up to 2 hours in nulliparous (3 hours with epidural), 1 hour in multiparous (2 hours with epidural). (3) Third stage: delivery of the placenta and membranes. Active management (oxytocin, controlled cord traction) reduces the risk of postpartum hemorrhage. Signs of labour: regular painful contractions, cervical dilation and effacement, show (blood-stained mucus), and spontaneous rupture of membranes (SROM). Induction of labour (IOL) is indicated for post-term pregnancy, pre-eclampsia, diabetes, intrauterine growth restriction, and premature rupture of membranes. Methods: prostaglandins (dinoprostone, misoprostol), balloon catheter, amniotomy, and oxytocin infusion. Assisted vaginal delivery: forceps or ventouse (vacuum) for prolonged second stage, fetal distress, or maternal exhaustion. Caesarean section: for failed IOL, fetal distress, malpresentation, placenta praevia, and previous caesarean section (VBAC may be attempted).

\medterm{Pelvic Floor Health and Disorders}-- The pelvic floor is a musculofascial sling (levator ani, coccygeus, and surrounding fascia) spanning the pelvic outlet, supporting the pelvic organs (bladder, uterus, rectum), maintaining continence, and facilitating urination, defecation, and sexual function. Pelvic floor disorders include pelvic organ prolapse (cystocele, rectocele, uterine prolapse, vaginal vault prolapse—affecting up to 50\% of parous women), stress urinary incontinence (leakage with cough, sneeze, exercise—treated with pelvic floor muscle training/Kegel exercises, pessaries, or midurethral sling), overactive bladder (urgency, frequency, urge incontinence—treated with behavioral modification, anticholinergics, beta-3 agonists), and pelvic pain syndromes (levator ani spasm, pudendal neuralgia). Risk factors: pregnancy/childbirth (vaginal delivery, especially with forceps or prolonged second stage), aging, obesity, chronic constipation, and heavy lifting. Diagnosis: Pelvic Organ Prolapse Quantification (POP-Q) system, urodynamic studies, and MRI defecography for complex prolapse. Treatment progresses from conservative (pelvic floor physical therapy, biofeedback, vaginal pessaries, lifestyle modification) to surgical (native tissue repair, sacrocolpopexy, colpocleisis for frail patients, sling procedures for SUI).

\medterm{Pelvic Inflammatory Disease}-- Infection and inflammation of the upper female genital tract (uterus, fallopian tubes, ovaries, and peritoneum), most commonly caused by sexually transmitted infections (Chlamydia trachomatis, Neisseria gonorrhoeae). Presents with lower abdominal pain, cervical motion tenderness, adnexal tenderness, and fever. Complications: tubo-ovarian abscess, infertility, ectopic pregnancy, and chronic pelvic pain. Diagnosis: clinical criteria (CDC minimal criteria: cervical motion, uterine, or adnexal tenderness). Treatment: empiric broad-spectrum antibiotics (ceftriaxone + doxycycline +/- metronidazole).

\medterm{Pelvic Inlet}-- The superior opening of the true pelvis, bounded by the sacral promontory posteriorly, the arcuate lines of the ilia laterally, and the symphysis pubis anteriorly. Its anteroposterior diameter (true conjugate) normally measures at least 11 cm. The transverse diameter is typically 13 cm. Contraction of the pelvic inlet (cephalopelvic disproportion) may prevent engagement of the fetal head, necessitating cesarean delivery.

\medterm{Pelvic Organ Prolapse}-- POP: descent of the pelvic organs through the vagina due to pelvic floor weakness. Types: cystocele (anterior—bladder), rectocele (posterior—rectum), uterine prolapse (apical), enterocele, and vault prolapse (post-hysterectomy). Risk factors: vaginal delivery (especially instrumental), multiparity, aging/menopause, obesity, chronic cough/constipation, and heavy lifting. Features: vaginal bulge/pressure, urinary incontinence/voiding difficulty, fecal symptoms, dyspareunia, and back pain; graded by POP-Q. Management: asymptomatic—observation; symptomatic—pelvic floor physiotherapy, pessary, and surgery (repair, sacrocolpopexy, hysterectomy); address contributing factors.

\medterm{Pelvic Outlet}-- The inferior opening of the true pelvis, bounded by the pubic symphysis, ischial tuberosities, sacrotuberous ligaments, and coccyx. Normal anteroposterior diameter is 11-13 cm (from the lower border of the pubic symphysis to the tip of the coccyx), and the transverse diameter (between the ischial tuberosities) is approximately 10-11 cm. A narrow pelvic outlet can cause shoulder dystocia or obstructed labour. Assessment of the pelvic outlet is part of clinical pelvimetry examination during prenatal care and is relevant when planning vaginal delivery, particularly in cases of suspected cephalopelvic disproportion.

\medterm{Pelvic Pain}-- Pain localized to the lower abdomen, pelvis, or perineum, lasting <3 months (acute) or >3--6 months (chronic), and can arise from gynecologic, urologic, gastrointestinal, musculoskeletal, or psychologic causes. Acute causes: ectopic pregnancy (life-threatening emergency), pelvic inflammatory disease (PID—sexually transmitted, polymicrobial; treated with antibiotics), ovarian torsion (sudden severe pain, adnexal mass; surgical emergency), acute cystitis, appendicitis, and diverticulitis. Chronic pelvic pain (CPP) affects ~15\% of women, with common etiologies: endometriosis (endometrial glands outside uterus, cyclic pain, dyspareunia, dysmenorrhea; laparoscopy for diagnosis; medical or surgical management), interstitial cystitis/bladder pain syndrome (bladder pressure/pain with filling, urgency, frequency), irritable bowel syndrome, pelvic floor myofascial pain, vulvodynia, and adenomyosis. Evaluation: bimanual exam, pregnancy test, imaging (transvaginal ultrasound, MRI, CT), laparoscopy for suspected endometriosis, and multidisciplinary assessment. Multimodal treatment: NSAIDs, hormones (OCPs, GnRH agonists for endometriosis), physical therapy, trigger point injections, CBT, and nerve blocks.

\medterm{Percutaneous Umbilical Blood Sampling (PUBS)}-- Also called cordocentesis. An invasive procedure for fetal blood sampling (cardocentesis) from the umbilical vein under ultrasound guidance, typically performed after 18-20 weeks. Indications: fetal anemia (Rh isoimmunization, parvovirus B19, hemoglobin Bart), fetal platelet disorders (immune thrombocytopenia), rapid karyotyping (late gestation), and fetal infection assessment (toxoplasmosis, CMV, parvovirus). Risks: fetal loss (1-2\%), bleeding from the puncture site, bradycardia, and cord hematoma.

\medterm{Perineal Laceration}-- Tearing of the perineum during vaginal delivery. First degree: skin only (perineal skin, fourchette). Second degree: perineal muscles (bulbocavernosus, transverse perineal) without anal sphincter involvement. Third degree: anal sphincter injury—3a (<50\% of EAS), 3b (>50\% of EAS), 3c (both EAS and IAS). Fourth degree: extends through the anal sphincter into the rectal mucosa. Risk factors: primiparity, operative vaginal delivery, midline episiotomy, large baby, and precipitous delivery. Repair: absorbable sutures, layers corresponding to the injury.

\medterm{Perineal Protection}-- Techniques used during the second stage to minimize perineal trauma (lacerations, episiotomy). Manual protection: the attendant's hand applies gentle pressure to the fetal head to control the speed of delivery (not too fast), supporting the perineum and flexing the fetal head. Warm compresses applied to the perineum during crowning may reduce third- and fourth-degree lacerations. Maternal pushing efforts should be gentle during crowning to allow the perineum to stretch gradually. Perineal massage in the second stage may reduce trauma.

\medterm{Perineal Tear}-- Laceration of the perineum during vaginal delivery, classified by degree: first degree (skin only), second degree (perineal muscles), third degree (sphincter complex involvement), and fourth degree (rectal mucosa). Most deliveries involve first or second degree tears. Third and fourth degree tears risk fecal incontinence. Risk factors include nulliparity, operative delivery, and midline episiotomy. Repair under local anesthesia with absorbable suture.

\medterm{Perineum}-- The diamond-shaped area between the pubic symphysis, coccyx, and ischial tuberosities, containing the urogenital and anal triangles. The urogenital triangle contains the external genitalia and urethral opening; the anal triangle contains the anus and coccygeus muscles. During delivery, the perineum stretches to allow passage of the fetus. Perineal tears classified as first through fourth degree.

\medterm{Peripartum Cardiomyopathy}-- See cardiology.

\medterm{Placental Abruption}-- Premature separation of the placenta from the uterine wall before delivery, occurring in 0.5-1\% of pregnancies. Risk factors: hypertension, preeclampsia, smoking, cocaine use, trauma, prior abruption, and thrombophilia. Presents with vaginal bleeding (concealed or revealed), abdominal pain, uterine tenderness, and uterine hypertonicity (board-like abdomen). Fetal distress may be present. Complications: maternal hemorrhage, DIC (due to release of thromboplastin), fetal hypoxia, and fetal death. Diagnosis: clinical; ultrasound (retroplacental clot, sensitivity limited). Management: immediate delivery regardless of gestational age if maternal or fetal compromise. Severe abruption: emergency cesarean, blood transfusion, and coagulopathy management.

\medterm{Placental Accreta Spectrum}-- Abnormal placental attachment to the myometrium, ranging from accreta (adherent to decidua basalis) to increta (invading myometrium) to percreta (invading through serosa). Major risk factor is prior caesarean section with anterior placenta previa. Diagnosis by prenatal ultrasound or MRI, with findings including loss of the clear zone between the placenta and myometrium, abnormal placental lacunae, turbulent flow on Doppler, and increased vascularity. Management requires a multidisciplinary team approach with planned caesarean hysterectomy, as attempted manual removal of the placenta leads to massive hemorrhage. Conservative management (leaving the placenta in situ) may be considered in carefully selected cases but carries risks of infection, delayed hemorrhage, and sepsis. Maternal mortality is approximately 5-10\%, primarily from hemorrhage. The rising incidence parallels the increasing caesarean delivery rate.

\medterm{Placental Separation Signs}-- Clinical signs indicating the placenta has separated from the uterine wall and is ready for delivery. Classic signs: lengthening of the umbilical cord (cord lengthens as the placenta descends), a sudden gush of blood, and uterine elevation and firming (the uterus rises and becomes globular as the placenta moves into the lower segment). The mother may feel an urge to push. Manual exam may show the placenta in the vagina. The Brandt-Andrews maneuver (controlled cord traction) is performed after signs of separation are observed.

\medterm{Placenta Praevia}-- Placental implantation over or near the internal cervical os, affecting approximately 0.5\% of pregnancies at term. Risk factors: previous caesarean section (risk increases with number of CS), placenta praevia in a prior pregnancy, multiple gestation, advanced maternal age (>35), smoking, cocaine use, and IVF. Classification per ultrasound: complete (placenta covers the os completely), partial (covers the os partially when dilated), marginal (placental edge reaches but does not cover the os), and low-lying (placental edge within 20 mm of the os). Presentation: painless, bright red vaginal bleeding in the second half of pregnancy (typically after 28 weeks—sentinel bleed). Diagnosis is by transvaginal ultrasound (accurate, safe—no risk of provoking bleeding). Management: (1) Antepartum: hospitalisation for significant bleeding, corticosteroids (betamethasone 12 mg IM \(\times\)2, 24 hours apart) for fetal lung maturity if preterm, anti-D immunoglobulin for Rh-negative women, avoid vaginal examination until placenta praevia is excluded. (2) Delivery: elective caesarean section at 36--37 weeks (increased bleeding risk), or earlier if bleeding is heavy. If the placenta is anterior and there is a prior CS scar, assess for placenta accreta spectrum (MRI or Doppler ultrasound). Complications: massive obstetric hemorrhage (blood transfusion often required), placenta accreta/increta/percreta, preterm delivery, and disseminated intravascular coagulation.

\medterm{Placenta Previa}-- A condition where the placenta implants in the lower uterine segment, partially or completely covering the internal cervical os. Classified as complete (total), partial, marginal, or low-lying. Major risk factor is prior cesarean section. Presents with painless, bright red vaginal bleeding in the third trimester. Diagnosis by transvaginal ultrasound. Management includes pelvic rest, corticosteroids for lung maturity, and cesarean delivery for complete previa.

\medterm{Polyhydramnios}-- Excess amniotic fluid volume, defined as amniotic fluid index greater than 24 cm or single deepest pocket greater than 8 cm. Occurs in 1-2\% of pregnancies. Causes include gestational diabetes, fetal anomalies (GI atresia, neural tube defects), chromosomal abnormalities, and idiopathic. Complications include preterm labor, preterm premature rupture of membranes, cord prolapse, and malpresentation. Mild cases managed expectantly; severe cases may require amnioreduction.

\medterm{Postpartum Blues}-- A transient, self-limited condition affecting 50-80\% of women, characterized by mood swings, tearfulness, irritability, and anxiety beginning 2-4 days after delivery and resolving within 2 weeks. No psychomotor disturbance or suicidal ideation. Treatment: reassurance, support, and rest. No pharmacotherapy needed. Differentiated from PPD by timing, duration, and lack of functional impairment. Women with postpartum blues are at increased risk of developing PPD and should be followed closely.

\medterm{Postpartum Cardiomyopathy (Peripartum Cardiomyopathy)}-- An idiopathic dilated cardiomyopathy presenting with heart failure in the last month of pregnancy or within 5 months of delivery. Incidence: 1:1000-1:4000 live births. Risk factors: advanced maternal age, African descent, multiple gestation, hypertension, and preeclampsia. Presents: dyspnea, orthopnea, edema, fatigue, and chest pain. Diagnosis: echocardiography (LVEF <45\%, LV dilatation). Management: standard heart failure therapy (beta-blockers, diuretics, ACE inhibitors/ARBs after delivery). Bromocriptine may improve outcomes (prolactin inhibition). Recovery: 50-60\% recover LVEF within 6 months. Prognosis: worse in those with LVEF <30\% at diagnosis.

\medterm{Postpartum Depression (PPD)}-- A major depressive episode occurring within 4 weeks of delivery, affecting 10-15\% of postpartum women. Symptoms persist >2 weeks and include depressed mood, anhedonia, sleep disturbance (insomnia despite exhaustion), appetite changes, fatigue, impaired concentration, guilt, and thoughts of harming self or the baby. Onset: most common within the first 3 months. Risk factors: prior depression, PPD in a previous pregnancy, stressful life events, lack of social support, and infant difficulties. Screening: Edinburgh Postnatal Depression Scale (EPDS, cutoff \(\ge\)10). Treatment: cognitive behavioral therapy (CBT), SSRIs (sertraline, citalopram, fluoxetine — compatible with breastfeeding), and social support.

\medterm{Postpartum Hemorrhage}-- Blood loss >500 mL after vaginal delivery or >1000 mL after cesarean section, the leading cause of maternal mortality worldwide. Causes: uterine atony (most common, 80\%), retained placenta, genital tract trauma (lacerations, episiotomy), and coagulopathy. Presents with heavy vaginal bleeding, hypotension, tachycardia, and pallor. Management: uterine massage, bimanual compression, uterotonic medications (oxytocin, misoprostol, methylergonovine, carboprost), balloon tamponade, uterine artery embolization, and surgical interventions (B-lynch suture, hysterectomy) for refractory cases.

\medterm{Postpartum Hemorrhage (Secondary)}-- Excessive bleeding occurring between 24 hours and 12 weeks postpartum. Causes: retained products of conception (most common), endometritis, subinvolution of the placental site, and coagulopathy. Presents with heavy or prolonged bleeding, cramping, and passing clots. Diagnosis: ultrasound (retained tissue, thickened endometrium). Management: uterine evacuation (D\&C), antibiotics, uterotonics, and rarely, uterine artery embolization. Retained products of conception can be missed if the ultrasound is performed too early (endometrial stripe not reliable before 2 weeks).

\medterm{Postpartum Infection}-- Infections occurring after delivery, a leading cause of maternal morbidity. Types: endometritis (most common, especially after cesarean), wound infection (cesarean incision, episiotomy, perineal laceration), urinary tract infection, mastitis, septic pelvic thrombophlebitis, and pneumonia. Risk factors: prolonged ROM, multiple vaginal exams, cesarean section, chorioamnionitis, and retained products of conception. Presentation: fever (\(\ge\)38°C after the first 24 hours postpartum), uterine tenderness, foul lochia, and wound erythema/drainage. Treatment: broad-spectrum antibiotics (gentamicin + clindamycin for endometritis, cephalexin for wound infection).

\medterm{Postpartum Psychosis}-- A rare but severe psychiatric emergency (1-2 per 1000 deliveries) characterized by the acute onset of delusions, hallucinations, disorganized behavior, confusion, and mood swings (mania or depression) within 1-4 weeks postpartum. Risk factors: bipolar disorder (50\% risk of recurrence in the postpartum period), prior postpartum psychosis, and a family history of bipolar disorder. High risk of infanticide and suicide requiring immediate psychiatric hospitalization. Treatment: mood stabilizers (lithium), antipsychotics (olanzapine, haloperidol), and ECT for severe cases. Breastfeeding considerations: lithium is relatively contraindicated; antipsychotics are preferred.

\medterm{Postpartum Thyroiditis}-- Autoimmune thyroid inflammation presenting within the first year postpartum (typically 2-6 months). Incidence: 5-10\% of women, especially those with positive TPO antibodies (50-70\% risk). Triphasic pattern: thyrotoxic phase (2-4 months, due to destructive release of stored thyroid hormone), followed by hypothyroid phase (4-8 months), followed by return to euthyroid (80\%). The thyrotoxic phase may be mistaken for postpartum depression (fatigue, palpitations, heat intolerance). Treatment: beta-blockers for thyrotoxic symptoms, levothyroxine for hypothyroidism. Up to 30\% develop permanent hypothyroidism. Recurrence in future pregnancies is common (70\%).

\medterm{Postterm Pregnancy}-- Pregnancy that extends beyond 42 weeks (294 days) of gestation, occurring in ~5-10\% of pregnancies. Risk factors: prior postterm pregnancy, inaccurate dating, male fetus, and obesity. Complications: increased perinatal mortality (2-3 fold), meconium aspiration syndrome, macrosomia with shoulder dystocia, oligohydramnios (leading to cord compression), and placental insufficiency. Management: accurate dating (first-trimester ultrasound), antepartum surveillance (BPP, NST, AFI) starting at 41 weeks. Induction of labor is recommended by 41-42 weeks to reduce perinatal complications.

\medterm{Preconception Counseling}-- Counseling provided to women of reproductive age before pregnancy to optimize outcomes. Components: folic acid supplementation (400 mcg, 4-5 mg for high risk), confirmation of immunity (rubella, varicella), management of chronic conditions (diabetes HbA1c <6.5\%, hypertension, epilepsy, thyroid disease), medication review (stop teratogens: isotretinoin, valproate, ACE inhibitors), vaccinations (MMR, varicella, hepatitis B, Tdap, influenza), healthy weight, smoking/alcohol cessation, genetic carrier screening (cystic fibrosis, spinal muscular atrophy, hemoglobinopathies per ethnicity), and reproductive life planning.

\medterm{Preeclampsia}-- A pregnancy-specific hypertensive disorder defined as new-onset hypertension (blood pressure 140/90 mmHg or higher) after 20 weeks gestation with proteinuria or other end-organ damage. Risk factors include nulliparity, advanced maternal age, obesity, chronic hypertension, diabetes, multiple gestation, preeclampsia in a prior pregnancy, and thrombophilia. Pathophysiology involves abnormal placentation, endothelial dysfunction, and systemic inflammatory response. Diagnosis requires BP \(\ge\)140/90 on two occasions at least 4 hours apart with proteinuria \(\ge\)300 mg/24 hours or protein/creatinine ratio \(\ge\)0.3. In the absence of proteinuria, preeclampsia is diagnosed if new-onset hypertension is accompanied by thrombocytopenia, renal insufficiency, impaired liver function, pulmonary edema, or cerebral/visual symptoms. Management includes antihypertensives for BP \(\ge\)160/110 mmHg (labetalol, nifedipine, hydralazine), magnesium sulfate for seizure prophylaxis, and delivery for gestational age \(\ge\)37 weeks or for severe features. Low-dose aspirin (81 mg daily) from 12-36 weeks reduces the risk in high-risk women.

\medterm{Preeclampsia with Severe Features}-- Preeclampsia (BP \(\ge\)140/90 + proteinuria/end-organ dysfunction) with one or more severe criteria: BP \(\ge\)160/110 mmHg, thrombocytopenia (<100,000/µL), impaired liver function (AST/ALT >2x normal, epigastric/RUQ pain), progressive renal insufficiency (creatinine >1.1 mg/dL), pulmonary edema, new-onset cerebral or visual symptoms (headache, scotomata, blurred vision). Management: antihypertensives for BP \(\ge\)160/110 (labetalol, hydralazine, nifedipine), magnesium sulfate for seizure prophylaxis, and delivery (regardless of gestational age if maternal/fetal instability, otherwise at 34-37 weeks with corticosteroids).

\medterm{Pregnancy}-- The physiological state in which a developing embryo or fetus grows within the female uterus, typically lasting approximately 40 weeks from the last menstrual period to childbirth. It involves complex hormonal and anatomical changes that support fetal development, including increased blood volume, altered metabolism, and growth of the placenta. Regular prenatal care is essential for monitoring maternal health and fetal growth, as well as identifying potential complications such as gestational diabetes, preeclampsia, or preterm labor.

\medterm{Pregnancy-Induced Hypertension}-- Hypertension developing after 20 weeks gestation without other features of preeclampsia. Blood pressure 140/90 mmHg or higher on two readings at least 4 hours apart. Affects 6-8\% of pregnancies. Managed with antihypertensive therapy (labetalol, nifedipine, methyldopa) when blood pressure exceeds 160/110 mmHg. Close monitoring for progression to preeclampsia. May resolve postpartum or persist as chronic hypertension.

\medterm{Premature Rupture of Membranes}-- PROM: spontaneous rupture of the amniotic membranes before the onset of labor; preterm PROM (PPROM) occurs before 37 weeks. Features: gush or leakage of fluid per vaginam; confirmed by sterile speculum (pooling, ferning, nitrazine pH test) or amniotic fluid biomarkers (IGFBP-1, PAMG-1). Risks: ascending infection (chorioamnionitis), cord prolapse, placental abruption, preterm birth. Management: confirm gestational age; for term—induction if not laboring (risk of infection); for PPROM—antenatal corticosteroids (lung maturity), latency antibiotics (erythromycin/azithromycin), magnesium sulfate (neuroprotection), and monitoring for chorioamnionitis; delivery for infection or fetal compromise.

\medterm{Prenatal Care}-- Regular medical care during pregnancy to monitor maternal and fetal health. Includes history and physical exam, weight and blood pressure monitoring, urine testing, blood work, fundal height measurements, fetal heart rate auscultation, and ultrasound. Schedules visits monthly in first trimester, biweekly in second, and weekly in third trimester. Screens for gestational diabetes, preeclampsia, anemia, infections, and fetal anomalies.

\medterm{Prenatal Vitamins}-- Multivitamin supplements formulated for pregnancy, containing folic acid (400-800 mcg), iron (27-60 mg), calcium (200-300 mg), vitamin D (400-600 IU), DHA (200 mg), and iodine (150-220 mcg). Started preconceptionally or at the first prenatal visit. Evidence supports NTD reduction from folic acid. Iron prevents maternal anemia. Additional screening detects deficiencies (vitamin D, B12, iron). Gummy vitamins may lack iron and calcium.

\medterm{Preterm Birth}-- Delivery before 37 weeks gestation, occurring in ~10\% of pregnancies worldwide. Subtypes: spontaneous preterm labor (50\%), preterm PROM (30\%), and medically indicated preterm birth (20\%, for maternal or fetal indications). Risk factors: prior preterm birth (strongest), short cervix, multiple gestation, smoking, infection, periodontal disease, and socioeconomic factors. Complications: respiratory distress syndrome (RDS), intraventricular hemorrhage (IVH), necrotizing enterocolitis (NEC), sepsis, and neurodevelopmental impairment. Prevention: progesterone supplementation, cervical cerclage, and screening for short cervix.

\medterm{Preterm Labor}-- Regular uterine contractions causing cervical change before 37 weeks gestation. Risk factors include prior preterm delivery, multiple gestation, infection, and cervical insufficiency. Diagnosis requires regular contractions with documented cervical change. Tocolytic therapy may delay delivery to allow corticosteroid administration for fetal lung maturity. Magnesium sulfate for neuroprotection if delivery anticipated before 32 weeks. Antenatal corticosteroids reduce respiratory distress syndrome.

\medterm{Preterm Premature Rupture of Membranes}-- Rupture of the amniotic membranes before the onset of labor at <37 weeks gestation. Causes: infection (subclinical chorioamnionitis), cervical insufficiency, and prior PPROM. Presents with sudden gush or leakage of fluid per vagina. Diagnosis: sterile speculum exam (pooling of amniotic fluid), ferning pattern on microscopy, and nitrazine test (turns blue, pH >6.5). Management: depends on gestational age. <24 weeks: counseling for expectant management or induction. 24-34 weeks: expectant management with corticosteroids (betamethasone for fetal lung maturity), antibiotics (ampicillin + erythromycin to prolong latency and reduce neonatal complications), and monitoring for chorioamnionitis. At term: induction of labor.

\medterm{Progesterone Supplementation}-- Administration of progesterone to prevent preterm birth in high-risk women. Natural progesterone (17-alpha hydroxyprogesterone caproate) shown to reduce preterm birth rate in women with prior spontaneous preterm delivery. Also beneficial for short cervix detected on screening. Typically started in the second trimester. Route options include weekly intramuscular injections or vaginal suppositories.

\medterm{Progestogen}-- A synthetic form of progesterone, the hormone that maintains the uterine lining in the secretory phase of the menstrual cycle and supports early pregnancy. Progestogens are used in: combined oral contraceptives (with estrogen), progestogen-only contraceptives (mini-pill, Depo-Provera injection, implant, Mirena IUS), HRT (with estrogen to protect the endometrium), treatment of menstrual disorders (heavy bleeding, endometriosis), and luteal phase support in IVF. Examples: norethisterone, levonorgestrel, medroxyprogesterone acetate, and micronized progesterone.

\medterm{Prolonged Deceleration}-- A fetal heart rate deceleration lasting \(\ge\)2 minutes but <10 minutes. If the deceleration persists \(\ge\)10 minutes, it is a baseline change. Causes: maternal hypotension (epidural-induced aortocaval compression), uterine hyperstimulation, umbilical cord prolapse, placental abruption, and maternal seizure. Management: immediate intrauterine resuscitation (position change, maternal oxygen, IV fluids, stop oxytocin, treat hypotension). Vaginal exam to exclude cord prolapse. If not resolved in minutes, emergency delivery (cesarean or operative vaginal).

\medterm{Prolonged Second Stage}-- Second stage of labor exceeding the time limits: nulliparas — 3 hours with regional anesthesia (2 hours without); multiparas — 2 hours with regional (1 hour without). Management: reassess position, contractions, and fetal status. For inadequate contractions: oxytocin augmentation. For malposition (OP, OT): manual rotation. For inadequate maternal effort: coaching or rest. If vaginal delivery is not imminent after a reasonable trial, operative vaginal delivery or cesarean is indicated. Prolonged second stage increases maternal (hemorrhage, chorioamnionitis, lacerations) and neonatal (acidosis, low APGAR, NICU admission) risks.

\medterm{Prostaglandins}-- Lipid compounds involved in labor initiation and cervical ripening. Endogenous prostaglandins (PGE2, PGF2alpha) stimulate uterine contractions and promote cervical softening. Synthetic analogs used clinically: misoprostol (PGE1) and dinoprostone (PGE2 gel or insert) for cervical ripening and labor induction. Also used for postpartum hemorrhage management. Contraindicated in patients with previous cesarean section due to uterine rupture risk.

\medterm{Pudendal Block}-- A local anesthetic block of the pudendal nerve (S2-S4) used for perineal anesthesia during the second stage of labor, outlet forceps delivery, or perineal repair. Technique: transvaginal approach — a needle guide is placed on the ischial spine, and the needle is advanced through the sacrospinous ligament. 1\% lidocaine (10 mL on each side) is injected. Provides perineal and lower vaginal analgesia. Does not provide uterine pain relief. Rarely used today because epidural analgesia provides superior pain relief. Complications: intravascular injection (seizures, cardiac toxicity), ischiorectal hematoma, and sciatic nerve block.

\medterm{Puerperal Fever}-- Fever higher than 38 degrees Celsius occurring during the first 10 days postpartum, not including the first 24 hours. Most common cause is endometritis, particularly after cesarean delivery. Other causes include urinary tract infection, wound infection, mastitis, and thrombophlebitis. Evaluation includes blood cultures, urinalysis, wound assessment, and pelvic exam. Treatment depends on suspected source; broad-spectrum antibiotics for endometritis.

\medterm{Puerperium}-- The period following delivery during which the reproductive organs return to their prepregnancy state, approximately 6 weeks in duration. Characterized by uterine involution, cervical closure, lochial discharge changes, and hormonal shifts. Physical recovery includes perineal healing, breast engorgement, and metabolic readjustment. Emotional changes common; screening for postpartum depression recommended. Return of ovulation variable, particularly with breastfeeding.

\medterm{Quickening}-- The first perception of fetal movement by the mother, typically felt between 16-20 weeks in multiparous women and 18-22 weeks in nulliparous women. Described initially as a subtle fluttering sensation that becomes stronger and more regular as pregnancy progresses. Factors affecting the timing include maternal body habitus, placental location (anterior placenta may delay perception), and gestational age. Quickening was historically used to estimate gestational age before ultrasound became available. The relationship between maternal perception and fetal activity patterns informs kick-count protocols for antepartum surveillance. Absence of perceived movement by 24 weeks warrants ultrasound evaluation for fetal wellbeing.

\medterm{Rectum}-- The final straight section of the large intestine, approximately 12-15 cm long, extending from the sigmoid colon to the anal canal. The rectum stores feces before defecation. In obstetrics, the rectum is relevant for: perineal tears (third- and fourth-degree involve the anal sphincter and rectal mucosa), rectocele (prolapse of the rectum into the posterior vaginal wall), and rectal examination (assessing descent of the fetal head or performing an enema).

\medterm{Renal Changes in Pregnancy}-- Renal blood flow increases 50-80\% by the second trimester, and GFR increases 50\% (creatinine decreases to 0.4-0.6 mg/dL, BUN to 8-10 mg/dL). Glycosuria and aminoaciduria are normal due to decreased tubular reabsorption. The ureters dilate (hydronephrosis of pregnancy, more prominent on the right) due to progesterone and mechanical compression by the gravid uterus. Proteinuria up to 300 mg/day is normal. The kidney increases in size (1-1.5 cm). Serum uric acid decreases due to increased clearance.

\medterm{Reproductive Organs}-- The organs involved in reproduction. Female: ovaries (produce eggs and hormones), fallopian tubes (transport eggs, site of fertilization), uterus (nurtures the developing fetus), cervix (lower part of uterus, dilates during labor), vagina (passage for sperm and baby), and vulva (external genitalia). Male: testes (produce sperm and testosterone), epididymis (sperm maturation), vas deferens (sperm transport), seminal vesicles and prostate (produce seminal fluid), and penis (delivers sperm). Reproductive years: the period from menarche to menopause during which a woman is capable of conceiving.

\medterm{Reproductive Years}-- The period of a woman's life during which she is capable of conceiving, typically from menarche (first menstrual period, average age 12-13) to menopause (final menstrual period, average age 51). During the reproductive years, women experience regular menstrual cycles, ovulation, and fertility that peaks in the mid-20s and declines after age 35. The most fertile window is approximately 5-6 days before ovulation. Contraception may be used to prevent pregnancy during the reproductive years.

\medterm{Respiratory Changes in Pregnancy}-- Functional residual capacity (FRC) decreases by 20\% due to elevated diaphragm from the gravid uterus. Tidal volume increases 30-40\%, minute ventilation increases 50\% (due to progesterone-mediated increase in respiratory drive). Result: chronic respiratory alkalosis (PaCO2 28-32 mmHg) with compensatory renal bicarbonate excretion (serum HCO3 18-22 mEq/L). PO2 is slightly increased (100-110 mmHg). Oxygen consumption increases 20\%. The decreased FRC means pregnant women desaturate faster during apnea (important for airway management).

\medterm{Retained Placenta}-- Failure of the placenta to deliver spontaneously within 30 minutes after delivery of the baby. Incidence: 1-3\% of deliveries. Types: placenta adherens (failure of separation due to inadequate myometrial contractions) and trapped placenta (separated but trapped behind a closed cervix). Risk factors: prior retained placenta, uterine anomaly, placenta previa, and preterm delivery. Management: manual removal of the placenta under anesthesia, oxytocin, and gentle cord traction. Retained placenta increta/percreta may require hysterectomy. Complications: hemorrhage, infection, and Asherman syndrome.

\medterm{Retained Products of Conception (RPOC)}-- Fragments of placental or fetal tissue retained in the uterus after delivery, miscarriage, or abortion. Causes: incomplete placental separation, succenturiate lobe, and placenta accreta. Presents: persistent or heavy bleeding, pain, and fever in the postpartum period. Diagnosis: ultrasound (echogenic mass in the uterine cavity, Doppler shows vascularity). Management: dilation and curettage (D\&C) or hysteroscopic resection. Antibiotics if febrile. RPOC is a common cause of secondary postpartum hemorrhage.

\medterm{Rh Isoimmunization}-- Maternal alloantibody production against fetal Rh D-positive red blood cells in an Rh D-negative mother, occurring when fetal RBCs enter the maternal circulation (fetomaternal hemorrhage). Sensitisation typically occurs at delivery of an Rh-positive infant, but may also follow miscarriage, ectopic pregnancy, amniocentesis, chorionic villus sampling, abdominal trauma, or antepartum hemorrhage. The maternal immune system produces anti-D antibodies (IgG) that cross the placenta and cause haemolysis of fetal RBCs, leading to fetal anemia, hydrops fetalis, and potentially fetal death in subsequent pregnancies. Prevention with RhoGAM (anti-D immune globulin) at 28 weeks gestation and within 72 hours of potential sensitising events has reduced the incidence from 10-15\% to less than 0.1\%. Affected pregnancies are managed with serial middle cerebral artery Doppler monitoring and intrauterine transfusions if fetal anemia develops.

\medterm{RhoGAM}-- Rh immune globulin containing anti-D antibodies, used to prevent Rh isoimmunization in Rh-negative mothers. Administered intramuscularly at 28 weeks gestation and within 72 hours after delivery of an Rh-positive infant. Also given after amniocentesis, chorionic villus sampling, miscarriage, ectopic pregnancy, or abdominal trauma. Works by coating and clearing fetal Rh-positive red blood cells before the maternal immune system can respond.

\medterm{Rubin Maneuver}-- An internal rotational maneuver for shoulder dystocia: the operator inserts a hand and applies pressure to the posterior aspect of the anterior fetal shoulder, rotating it toward the fetal chest (adducting the shoulders). This reduces the bisacromial diameter and disimpacts the anterior shoulder from the pubic symphysis. Often performed in sequence: Rubin I (external suprapubic pressure), Rubin II (internal pressure on the posterior aspect of the anterior shoulder). Performed after McRoberts and suprapubic pressure.

\medterm{Salpingectomy}-- Surgical removal of a fallopian tube, the definitive treatment for tubal ectopic pregnancy and for hydrosalpinx. Laparoscopic approach is standard (three-port technique). Indications: ruptured ectopic, hemodynamic instability, failed medical management, contralateral tube is healthy, and patient desires sterilization. Salpingectomy preserves ovarian function if the blood supply (ovarian artery anastomosis) is preserved. Methotrexate is an alternative for early, unruptured ectopic pregnancy.

\medterm{Salpingitis}-- Inflammation of the fallopian tubes, most commonly caused by ascending infection from the lower genital tract. Salpingitis is a component of pelvic inflammatory disease and is often bilateral. etiology: sexually transmitted infections—Chlamydia trachomatis (most common), Neisseria gonorrhoeae, and bacterial vaginosis-associated organisms (Gardnerella vaginalis, anaerobes—Bacteroides, Peptostreptococcus, Atopobium vaginae). Clinical: bilateral lower abdominal pain, abnormal vaginal/cervical discharge, fever, dysmenorrhoea, deep dyspareunia, intermenstrual bleeding, and cervical motion tenderness (chandelier sign) on bimanual examination. Diagnosis: clinical (CDC diagnostic criteria—minimum: cervical motion, uterine, or adnexal tenderness + one of: cervical mucopurulent discharge, increased WBC on saline microscopy of vaginal fluid, positive chlamydia/gonorrhea NAAT). Ultrasound may show thickened, fluid-filled tubes (tubo-ovarian complex or abscess). Complications: tubo-ovarian abscess, Fitz–Hugh–Curtis syndrome (perihepatitis—right upper quadrant pain), infertility (tubal factor—risk increases with number of episodes: 10\% after 1 episode, 25\% after 2, 50\% after 3), ectopic pregnancy (6--10-fold increased risk), and chronic pelvic pain. Treatment: antibiotics (regimen: ceftriaxone 500 mg IM single dose + doxycycline 100 mg PO BD for 14 days + metronidazole 500 mg PO BD for 14 days). Hospitalise if severe disease, pregnancy, tubo-ovarian abscess, or failure of oral therapy. Sexual partner(s) must be treated. Prevention: barrier contraception (condoms), chlamydia/gonorrhea screening in sexually active women <25.

\medterm{Salpingostomy (Salpingotomy)}-- A fertility-preserving surgical procedure for tubal ectopic pregnancy: a linear incision on the antimesenteric border of the fallopian tube to remove the ectopic gestation, leaving the tube in situ. Indications: unruptured ectopic, desire for future fertility, contralateral tubal damage or absence. Risk: persistent trophoblast (up to 15\%, requires hCG monitoring and possibly methotrexate). The tube may heal with scarring, reducing patency (50-60\% patency rate). Fertility rates after salpingostomy are similar to salpingectomy if the contralateral tube is normal.

\medterm{Second Stage of Labor}-- From full cervical dilatation (10 cm) to delivery of the baby. The mother feels the urge to push as the presenting part descends. Pushing: closed-glottis (Valsalva) or open-glottis. Duration: up to 3 hours in nulliparas with regional anesthesia (2 hours without), and 2 hours in multiparas (1 hour without). Prolonged second stage: longer than these limits. Management: continued maternal effort, operative vaginal delivery (vacuum or forceps), or cesarean delivery if vaginal delivery is not imminent.

\medterm{Second Trimester}-- The period from 13 to 27 weeks gestational age. Fetal growth accelerates, mother begins to feel quickening (fetal movement) between 16-20 weeks. Anatomy ultrasound at 18-22 weeks screens for structural anomalies. Maternal serum screening offered. Nausea typically resolves. Uterus becomes palpable above the symphysis pubis. Amniocentesis may be performed between 15-20 weeks. Risk of miscarriage decreases substantially after the first trimester. Cervical length screening by transvaginal ultrasound is performed in high-risk women between 18-24 weeks to assess preterm birth risk.

\medterm{Sexual Dysfunction}-- Persistent difficulty in sexual activity causing distress. In women: low sexual desire (hypoactive), arousal disorder, orgasmic disorder (anorgasmia), and genito-pelvic pain/penetration disorder (dyspareunia, vaginismus). Causes: psychological (stress, relationship, depression, anxiety), endocrine (menopause, hypogonadism, thyroid, prolactin), vascular/neurologic (diabetes, MS), medication (SSRIs, antihypertensives), gynecologic (endometriosis, vaginismus, pelvic floor, infection, atrophy), and surgery (pelvic, breast cancer). Evaluation: history (medical, psychosexual, medication), exam, and labs (hormones). Management: treat underlying cause, counseling/sex therapy, lubricants, estrogen (atrophy), flibanserin/ testosterone (selected), and address relationship factors.

\medterm{Sexual Intercourse (Coitus)}-- The insertion of a penis into a partner's body for sexual pleasure, reproduction, or both. Most commonly vaginal intercourse (penile-vaginal penetration) but may also include anal intercourse. Sexual intercourse is a key component of human pair bonding and reproduction. From a medical perspective, it is associated with risks including sexually transmitted infections (STIs), unintended pregnancy, and physical injury (e.g., vaginal or penile trauma). Safe sex practices (condoms, contraception, PrEP, vaccination) reduce these risks. Sexual dysfunction can affect any stage of the sexual response cycle (desire, arousal, orgasm).

\medterm{Sheehan Syndrome (Postpartum Pituitary Necrosis)}-- Ischemic necrosis of the anterior pituitary gland due to severe postpartum hemorrhage (PPH) and hypotension. The pituitary enlarges during pregnancy and is vulnerable to ischemia from hypotensive episodes. Clinical features: failure of lactation (first sign, most common), amenorrhea, loss of axillary and pubic hair, breast atrophy, fatigue, hypoglycemia, and hypotension. Diagnosis: low serum cortisol, low TSH/T4, low LH/FSH, low prolactin, and MRI (empty sella). Treatment: lifelong hormone replacement (glucocorticoids first, thyroid hormone, estrogen, and growth hormone as needed).

\medterm{Shirodkar Cerclage}-- A transvaginal cervical cerclage technique that involves dissection of the bladder and rectum to place the suture (Mersilene tape) as high as possible at the level of the internal os. Technically more difficult than McDonald cerclage with higher risk of hemorrhage and bladder injury. Indications: cervical insufficiency when the cervix is too short for McDonald placement, or failed McDonald cerclage. The suture is buried below the vaginal mucosa. Removed at 36-37 weeks.

\medterm{Shoulder Dystocia}-- An obstetric emergency where the fetal head delivers but the anterior shoulder is impacted behind the maternal symphysis pubis. Defined as head-to-body delivery time exceeding 60 seconds or need for ancillary maneuvers. Risk factors include macrosomia, diabetes, prior shoulder dystocia, and operative vaginal delivery. Management follows HELPERR mnemonic: call for help, episiotomy, legs (McRoberts), suprapubic pressure, internal rotation (Woods screw), delivery of posterior arm, and clavicle fracture.

\medterm{Sinusoidal Fetal Heart Rate Pattern}-- A rare, ominous FHR pattern characterized by a smooth, sine-wave-like undulation with a fixed frequency (2-5 cycles/minute) and absent variability. Strongly associated with severe fetal anemia (Rh isoimmunization, fetomaternal hemorrhage, parvovirus B19), fetal hypoxia, and fetal acidosis. Differential: pseudo-sinusoidal pattern (from maternal Demerol administration, usually resolves). A true sinusoidal pattern is a Category III tracing requiring immediate evaluation, intrauterine resuscitation, and urgent delivery.

\medterm{Soft Markers for Aneuploidy}-- Sonographic findings in the second trimester that are common in normal fetuses but have a statistically higher prevalence in aneuploid fetuses (especially trisomy 21). Common soft markers: echogenic intracardiac focus (EIF), echogenic bowel, renal pyelectasis (4-7 mm), short femur/humerus, nuchal fold \(\ge\)6 mm, choroid plexus cyst, and sandal gap. Each marker has a likelihood ratio (LR) for trisomy 21 (1.5-5). In isolation and with low-risk serum screening, no further workup is needed. Multiple markers increase risk and warrant genetic counseling.

\medterm{Sonographer}-- A healthcare professional trained to perform ultrasound scans (sonography). In obstetrics, sonographers perform: dating scans (first trimester), nuchal translucency measurement (11-14 weeks), anomaly scans (18-21 weeks), growth scans (serial assessment of fetal size and amniotic fluid), and Doppler studies (umbilical artery, middle cerebral artery). Sonographers also perform gynecological scans (pelvic ultrasound for ovarian cysts, fibroids, endometriosis). They work closely with obstetricians and radiologists to interpret findings.

\medterm{Spinal Anesthesia for Cesarean Delivery}-- A neuraxial anesthetic technique used for elective and many emergency cesarean deliveries. A small dose of local anesthetic (hyperbaric bupivacaine 12-15 mg) with opioid (fentanyl 15-25 mcg, morphine 0.15-0.25 mg) is injected into the subarachnoid space. Produces rapid surgical anesthesia (T4 sensory level for cesarean). Advantages: minimal fetal drug exposure (compared to general anesthesia), the mother is awake for delivery, and excellent postoperative analgesia with intrathecal morphine. Disadvantages: maternal hypotension (from sympathectomy — prevent with left uterine displacement, IV fluids, vasopressors), pruritus, nausea, and PDPH (post-dural puncture headache, 1-2\% with pencil-point needles).

\medterm{Stages of Labor}-- Four stages describing the progression from uterine contractions to delivery of the placenta. First stage: onset of regular contractions to full cervical dilatation (10 cm). Second stage: full dilatation to delivery of the baby. Third stage: delivery of the baby to delivery of the placenta. Fourth stage: the first 1-2 hours after placental delivery (immediate postpartum, highest risk of hemorrhage). Normal labor requires coordinated uterine contractions, progressive cervical change, and descent of the fetal presenting part.

\medterm{Station}-- The relationship of the fetal presenting part to the level of the ischial spines, measured in centimeters.

\medterm{Sterilisation}-- A permanent method of contraception. Female sterilization (tubal occlusion): the fallopian tubes are blocked or cut to prevent the egg from reaching the uterus, performed by laparoscopy (clip or diathermy) or at the time of cesarean delivery. Male sterilization (vasectomy): the vas deferens is cut or sealed to prevent sperm from entering the ejaculate. Both are highly effective (>99\%) and intended to be permanent, though reversal is sometimes possible through microsurgery.

\medterm{Stillbirth}-- Intrauterine fetal death (IUFD) at or after 20-24 weeks (definition varies by country); fetal death before viability (<20 weeks) is miscarriage. Causes: placental abruption, placental insufficiency, preeclampsia, fetal anomalies, infection (Listeria, parvovirus, CMV, syphilis), umbilical cord accidents, diabetes, IUGR, maternal disease (antiphospholipid, lupus), and unexplained (~25 percent). Diagnosis: absent fetal heart tones/ultrasound. Management: confirm death, discuss delivery options (induction), psychological support, investigations (autopsy, placental pathology, karyotype, TORCH, thrombophilia), and prevention of recurrence counseling. Bereavement support and follow-up are essential.

\medterm{Striae Gravidarum}-- Stretch marks developing on the abdomen, breasts, thighs, and buttocks during pregnancy due to mechanical stretching of the dermis combined with hormonal effects on collagen. Initially appear as red or purple striae (striae rubrae) and eventually fade to silver-white (striae albae). Occur in 50-90\% of pregnancies. Prevention with oils and creams has limited evidence. Do not affect fetal development.

\medterm{Stripping of Membranes (Membrane Sweep)}-- A procedure performed during a cervical exam where the examiner inserts a finger through the internal os and sweeps the membranes circumferentially, separating them from the lower uterine segment. Releases endogenous prostaglandins, promoting cervical ripening and potentially initiating labor. Performed at 38-41 weeks to reduce the need for formal induction. Discomfort: moderate cramping and spotting. Contraindications: low-lying placenta, placenta previa, and active bleeding. Studies show reduced postterm pregnancy and induction rates with serial sweeping. Nipple stimulation may also augment endogenous oxytocin.

\medterm{Substance Use in Pregnancy}-- Use of alcohol, tobacco, marijuana, opioids, cocaine, and other drugs during pregnancy. Tobacco: increased risk of IUGR, preterm birth, placental abruption, SIDS, and low birth weight. Opioids: neonatal abstinence syndrome (NAS, withdrawal in 60-80\% of exposed newborns), treated with morphine or methadone. Alcohol: FASD. Marijuana: limited evidence but possible impact on fetal neurodevelopment. Cocaine: abruption, IUGR, preterm birth. Screening: urine drug screen with patient consent (universal screening recommended, but risk of legal/social consequences).

\medterm{Succenturial Lobe}-- An accessory lobe of the placenta, connected to the main placenta by blood vessels that traverse the membranes. If undetected after delivery, the succenturial lobe may be retained in the uterus, causing postpartum hemorrhage or infection. Diagnosis: careful inspection of the placenta and membranes after delivery. A succenturial lobe can also be associated with vasa praevia (fetal vessels running across the internal cervical os), which can cause life-threatening fetal hemorrhage when membranes rupture.

\medterm{Superimposed Preeclampsia}-- Preeclampsia developing in a woman with preexisting chronic hypertension. Diagnosis: sudden worsening of hypertension, new-onset proteinuria or a sudden increase in proteinuria, or new-onset thrombocytopenia, elevated liver enzymes, or symptoms (headache, visual changes). Occurs in 20-50\% of women with chronic hypertension. Higher risk of adverse maternal and fetal outcomes compared to preeclampsia alone. Requires close monitoring and early delivery if severe features develop.

\medterm{Supine Hypotensive Syndrome}-- Compression of the inferior vena cava by the gravid uterus when the pregnant woman lies flat on her back after 20 weeks gestation. Results in decreased venous return, reduced cardiac output, and maternal hypotension, dizziness, pallor, and syncope. Uteroplacental perfusion may be compromised, potentially causing fetal bradycardia and distress. The syndrome typically occurs after 20 weeks gestation when the uterus is large enough to compress the great vessels. Prevention involves avoiding the supine position and maintaining left lateral tilt during labour, surgery, and transport. Treatment includes repositioning to the left lateral decubitus position and IV fluid bolus. The syndrome is especially relevant during regional anesthesia (spinal or epidural) where loss of sympathetic tone compounds the hemodynamic effect.

\medterm{Sutures}-- Surgical threads or stitches used to close wounds or surgical incisions. In obstetrics, sutures are used for: episiotomy repair, perineal tear repair (1st-4th degree), cesarean section closure (uterine, fascial, and skin layers), and cervical cerclage. Absorbable sutures (Vicryl, Monocryl, chromic catgut) dissolve over time and do not require removal. Non-absorbable sutures (nylon, Prolene) require removal. Continuous suturing is faster; interrupted sutures may provide better wound edge apposition.

\medterm{Tachysystole (Uterine Hyperstimulation)}-- More than 5 contractions in 10 minutes averaged over 30 minutes, or contractions lasting \(\ge\)2 minutes. Tachysystole with fetal heart rate changes (late decelerations, prolonged deceleration, loss of variability) is more concerning. Causes: oxytocin (most common), prostaglandins (misoprostol > dinoprostone), and cocaine. Management: stop oxytocin (if running), administer a tocolytic (terbutaline 0.25 mg SQ for persistent tachysystole with nonreassuring FHR), maternal position change, IV fluids, and oxygen. Uterine hypertonus: increased resting tone between contractions.

\medterm{Tdap Vaccination in Pregnancy}-- Administration of tetanus, diphtheria, and acellular pertussis (Tdap) vaccine during each pregnancy, optimally at 27-36 weeks. The goal: boost maternal antibodies to protect the newborn from pertussis (whooping cough) before the infant can be vaccinated at 2 months. Transplacental transfer of anti-pertussis antibodies provides passive immunity to the infant. Tdap is safe in pregnancy. The cocooning strategy (vaccinating family members) is no longer recommended as a primary strategy; maternal vaccination is more effective.

\medterm{Third Stage of Labor}-- From delivery of the baby to delivery of the placenta. The placenta separates from the uterine wall and is expelled spontaneously or with gentle cord traction. Normal duration: 5-30 minutes. Active management of the third stage: administration of oxytocin immediately after delivery of the baby, controlled cord traction, and uterine massage. This reduces the risk of postpartum hemorrhage by 60\%. Retained placenta: failure to deliver the placenta within 30 minutes; management includes manual removal or curettage.

\medterm{Third Trimester}-- The period from 28 weeks to delivery. Rapid fetal weight gain occurs, lungs mature with surfactant production. Fetal movements become regular and countable. Braxton Hicks contractions begin. Group B Streptococcus screening at 36 weeks. Fetal positioning assessed for delivery. Viability generally considered at 23-24 weeks. Cesarean section rate highest in this trimester.

\medterm{Thrush}-- A common vaginal infection caused by the yeast Candida albicans (also called candidiasis, vaginal thrush, or candida vulvovaginitis). Symptoms: itching, soreness, white discharge (like cottage cheese), dyspareunia, and external dysuria. Risk factors: pregnancy, antibiotic use, diabetes, immunosuppression, and high estrogen states. Diagnosis: clinical and confirmed by vaginal swab (microscopy and culture showing budding yeast and pseudohyphae). Treatment: topical clotrimazole or oral fluconazole (single dose). Recurrent thrush (>4 episodes/year) may require maintenance therapy.

\medterm{Tocolysis}-- The use of medications to suppress uterine contractions and delay preterm labor, providing time for corticosteroid administration (fetal lung maturity) and in utero transfer to a unit with appropriate neonatal facilities. First-line: nifedipine (calcium channel blocker) or atosiban (oxytocin receptor antagonist). Second-line: indomethacin (NSAID, short-term). Beta-agonists (terbutaline, ritodrine) are no longer first-line due to maternal side effects. Maintenance tocolysis beyond 48 hours is not recommended.

\medterm{Tocolytic Therapy}-- Medications used to inhibit uterine contractions temporarily to prolong pregnancy, allowing time for corticosteroid administration and transport to a tertiary care center. Indications: preterm labor with intact membranes, 24-34 weeks. First-line: nifedipine (calcium channel blocker, 20-30 mg PO q3-4h) or indomethacin (NSAID, 50-100 mg PR/PO then 25-50 mg PO q6h, limited to 48 hours due to oligohydramnios and premature ductus arteriosus closure). Second-line: terbutaline (beta-agonist, SQ, not for prolonged use). Atosiban (oxytocin receptor antagonist) is used in Europe. Tocolysis should not delay delivery in the setting of chorioamnionitis, severe preeclampsia, or fetal compromise.

\medterm{Total Breech Extraction}-- A maneuver where the operator grasps the fetal feet and delivers the entire baby in breech presentation. The most common indication is delivery of the second twin (malpresentation). Contraindicated for term singleton breech due to high risk of head entrapment, nuchal arms, and birth trauma. Technique: identify the feet, grasp them at the ankles, and apply gentle traction. Deliver the legs, trunk, shoulders, and head sequentially. Anesthesia: adequate relaxation (regional or general) is needed.

\medterm{Toxemia}-- An older term for pre-eclampsia (also called toxemia of pregnancy), reflecting the now-outdated theory that the condition was caused by toxins in the blood. Pre-eclampsia is a multisystem disorder of pregnancy characterized by hypertension and proteinuria developing after 20 weeks gestation, affecting 2-8\% of pregnancies. Features: elevated blood pressure (\(\ge\)140/90 mmHg), proteinuria (\(\ge\)300 mg/24h), and potentially headache, visual disturbances, epigastric pain, and seizures (eclampsia). The term toxemia has been largely replaced by pre-eclampsia in modern practice.

\medterm{Transabdominal Scan}-- An ultrasound scan performed with the transducer placed on the maternal abdomen. A full bladder improves visualization of pelvic structures. Used in early pregnancy to confirm viability and gestational age, and throughout pregnancy to assess fetal growth, anatomy, amniotic fluid volume, and placental location. Transabdominal scans provide a wider field of view than transvaginal scans but have lower resolution for early pregnancy (before 10-12 weeks) and detailed assessment.

\medterm{Transcervical Foley Catheter for Cervical Ripening}-- A mechanical method for cervical ripening: a Foley catheter (16-30 Fr) is inserted through the cervix, the balloon is inflated with 30-60 mL of saline, and traction is applied (or the catheter is taped to the thigh with gentle tension). The balloon applies direct pressure to the internal os, stimulating endogenous prostaglandin release. Advantages: no increased risk of uterine tachysystole (unlike pharmacologic agents), safe in patients with prior cesarean, and minimal systemic side effects. Often combined with oxytocin after the balloon falls out. Typically <12 hours from insertion to expulsion.

\medterm{Transverse Lie}-- A fetal lie where the long axis of the fetus is perpendicular to the long axis of the mother. The presenting part is a shoulder or arm. Occurs in approximately 1\% of term pregnancies. Risk of cord prolapse and obstructed labor. Cesarean delivery required as vaginal delivery is not possible. Version may be attempted to convert to longitudinal lie but often recurs. More common in multiparous women with lax abdominal walls.

\medterm{TRAP Sequence}-- Twin Reversed Arterial Perfusion sequence, a complication of monochorionic twin gestations where one twin (the pump twin) perfuses the other (the acardiac twin) via arterioarterial anastomoses. The acardiac twin receives retrograde, deoxygenated blood through its umbilical artery, resulting in severe anomalies (acardia, acephaly, rudimentary limbs) and no functional heart. The pump twin is structurally normal but at risk for high-output cardiac failure, polyhydramnios, and hydrops from the hemodynamic burden of perfusing the acardiac twin. Diagnosis by ultrasound showing a structurally abnormal twin with absent cardiac activity lacking the visible heartbeat. Doppler demonstrates reverse umbilical artery flow to the acardiac twin. Risk to the pump twin: mortality up to 50\% without intervention. Treatment: fetoscopic laser coagulation or radiofrequency ablation of the acardiac twin's umbilical cord to stop perfusion. Outcomes for the pump twin are excellent with early intervention.

\medterm{Trial of Labor After Cesarean (TOLAC)}-- The attempt at vaginal birth in a woman with a prior cesarean, which may result in VBAC or a failed TOLAC (requiring repeat cesarean). Candidates: one prior low-transverse cesarean, clinically adequate pelvis, no contraindications. Counseling: success rate 60-80\%, uterine rupture risk 0.5-0.9\%, perinatal mortality from rupture <0.1\%. Induction of labor increases rupture risk (especially prostaglandins, avoid misoprostol). Oxytocin may be used with caution. Patients must deliver at hospitals capable of emergency cesarean within 30 minutes.

\medterm{Triplets}-- A multiple pregnancy with three fetuses (trichorionic, dichorionic, or monochorionic combinations); commonly results from assisted reproductive technology. High-risk: preterm delivery (average ~32-33 weeks), low birth weight, preeclampsia, gestational diabetes, growth restriction, and complications of monochorionicity (TTTS). Antenatal: frequent monitoring, growth scans, corticosteroid prophylaxis, and planned cesarean delivery (usually). Neonatal: high rates of NICU admission, respiratory distress, and long-term neurodevelopmental concerns. Antenatal counseling, maternal nutrition, and multidisciplinary management are essential.

\medterm{Tubal Pregnancy}-- An ectopic pregnancy implanted in the fallopian tube, accounting for >95\% of ectopic pregnancies. Most common site: ampulla (70\%). Risk factors: prior ectopic, PID (chlamydia), tubal surgery, smoking, IUD, and ART. Presents with lower abdominal pain, amenorrhea, and vaginal bleeding. Diagnosis: transvaginal ultrasound (no intrauterine pregnancy with hCG above discriminatory zone, typically 1500-3000 IU/L), adnexal mass, or tubal ring. Unruptured: treated with methotrexate (single-dose regimen: 50 mg/m$^2$ IM) or salpingectomy/salpingostomy. Ruptured: surgical emergency with hemorrhagic shock requiring salpingectomy.

\medterm{Twin-to-Twin Transfusion Syndrome}-- See complete definition in the Multiple Gestation section above.

\medterm{Twin-Twin Transfusion Syndrome}-- A complication of monochorionic diamniotic twin pregnancies where placental vascular anastomoses (arteriovenous connections) cause unbalanced blood flow from one twin (the donor) to the other (the recipient). The donor twin becomes hypovolemic, anemic, and growth-restricted with oligohydramnios ("stuck twin"). The recipient twin becomes hypervolemic, polycythemic, and develops polyhydramnios, cardiac dysfunction, and hydrops. Quintero staging (I-V): I — oligohydramnios/polyhydramnios sequence without Doppler abnormalities; II — bladder not visible in donor; III — abnormal Doppler studies (absent/reversed end-diastolic flow in donor or reverse flow in ductus venosus in recipient); IV — hydrops in recipient; V — intrauterine fetal demise of one or both twins. Diagnosis by ultrasound in monochorionic twins. Treatment: fetoscopic laser photocoagulation of the placental anastomoses (first-line for Quintero II-IV). Severe TTTS without treatment has a mortality rate exceeding 80\%.

\medterm{Umbilical Artery Doppler}-- Ultrasound assessment of blood flow in the umbilical arteries to evaluate placental resistance. Normal flow shows a pulsatility index decreasing throughout pregnancy with forward diastolic flow. Absent or reversed end-diastolic flow indicates severe placental dysfunction and is associated with increased perinatal morbidity and mortality. Used in management of intrauterine growth restriction to guide timing of delivery.

\medterm{Umbilical Cord}-- The connection between the fetus and the placenta, containing two umbilical arteries (carrying deoxygenated blood from the fetus to the placenta) and one umbilical vein (carrying oxygenated blood from the placenta to the fetus). The vessels are embedded in Wharton jelly, a mucopolysaccharide matrix that prevents compression. Normal length: 55-60 cm. The cord inserts centrally into the placenta in most cases. Abnormalities: velamentous insertion (vessels insert into the membranes), vasa previa (vessels cross the internal os), true knots, nuchal cord (around the fetal neck), and single umbilical artery (two-vessel cord).

\medterm{Umbilical Cord Prolapse}-- Descent of the umbilical cord through the cervix alongside or past the presenting fetal part, occurring when the membranes rupture. Risk factors: malpresentation (transverse lie, breech), preterm/low birth weight, polyhydramnios, multiparity, and artificial ROM with high presenting part. Presents as cord palpable on vaginal exam or visible at the introitus after ROM, often with sudden fetal bradycardia. Management: immediate cesarean delivery. While preparing for delivery, relieve cord compression by elevating the presenting part (knee-chest position, Trendelenburg, or manual elevation). Keep the cord inside the vagina to prevent vasospasm. Tocolysis may be used.

\medterm{Upright Labor Positions}-- Maternal positions during labor that use gravity to promote fetal descent: walking, standing, squatting, kneeling, and sitting upright on a birth stool. Benefits of upright vs. recumbent positions: shorter first stage (~1 hour), less epidural use, fewer operative deliveries, less perineal trauma, and less severe pain. The supine position is discouraged because it decreases uterine blood flow and cardiac output. Upright positions increase pelvic diameters and contraction efficiency. Water immersion (hydrotherapy) in the first stage is also beneficial.

\medterm{Urinary Tract Infection in Pregnancy}-- Asymptomatic bacteriuria (2-10\% of pregnancies) and symptomatic UTI (cystitis, pyelonephritis). Pregnancy increases the risk of pyelonephritis due to hormonally mediated ureteral dilatation and stasis. Asymptomatic bacteriuria: screening and treatment (nitrofurantoin, cephalexin, amoxicillin) reduce the risk of pyelonephritis and preterm birth. Cystitis: dysuria, frequency, urgency. Pyelonephritis: fever, flank pain, nausea, higher risk of preterm labor/preterm delivery. Treatment: IV antibiotics (ceftriaxone, gentamicin, ampicillin+gentamicin) for pyelonephritis.

\medterm{Uterine Atony}-- Failure of the uterus to contract adequately after delivery, resulting in postpartum hemorrhage. The most common cause of PPH. Risk factors include prolonged labor, multiple gestation, macrosomia, oxytocin augmentation, and uterine fibroids. Diagnosis is clinical: boggy, soft uterus with excessive bleeding. Management: uterine massage, oxytocin, methylergonovine, carboprost, misoprostol, balloon tamponade, and surgical interventions (B-Lynch suture, uterine artery ligation) for refractory cases.

\medterm{Uterine Balloon Tamponade}-- A minimally invasive technique for managing postpartum hemorrhage due to uterine atony when uterotonics fail. A balloon (Bakri Postpartum Balloon, 500-1000 mL capacity) is inserted transvaginally into the uterine cavity and inflated with saline until bleeding stops. The balloon provides mechanical compression against the uterine wall. Success rate: 70-90\%. The balloon is left in place for 12-24 hours, then gradually deflated. If bleeding continues despite balloon tamponade, surgical intervention (uterine artery embolization, B-Lynch suture, hysterectomy) is indicated. Uterine rupture must be excluded before balloon placement.

\medterm{Uterine Inversion}-- A rare but life-threatening obstetric emergency where the uterus turns inside out, protruding through the cervix or vaginal introitus. Occurs in 1:2000-1:20,000 deliveries. Causes: fundal pressure, excessive cord traction (especially with fundal placenta), uterine atony, and fundal fibroids. Presents: sudden massive hemorrhage, shock out of proportion to blood loss, and a mass at the introitus or vagina. Management: Johnson maneuver (push the fundus back through the cervix with an open palm), tocolytics (terbutaline, magnesium) to relax the uterus, and if manual replacement fails, a surgical approach (Huntington or Hautlain procedure). IV fluids, blood products, and uterotonics after replacement.

\medterm{Uterine Prolapse}-- Descent of the uterus into the vaginal canal due to weakening of the uterosacral/cardinal ligaments and pelvic floor, most commonly after childbirth and menopause. Grades 1-3 based on descent. Risk: vaginal delivery, obesity, aging, chronic straining, heavy lifting. Features: vaginal bulge/pressure, backache, urinary symptoms (incontinence, retention), sexual dysfunction; worse with standing. Diagnosis: clinical (POP-Q). Management: conservative—pelvic floor physiotherapy, pessary, weight loss, treat constipation; surgical—hysterectomy with apical suspension, sacrocolpopexy, or Manchester repair for symptomatic prolapse; counseling for future pregnancy.

\medterm{Uterine Rupture}-- A full-thickness tear through the uterine wall, potentially life-threatening to both mother and fetus. Risk factors include prior caesarean section (especially classical incision), oxytocin augmentation, and operative vaginal delivery. Presents with sudden abdominal pain, loss of fetal heart rate, and maternal hemorrhage. Requires emergent caesarean delivery and surgical repair. Risk is approximately 0.5\% in women with a prior low-transverse caesarean attempting trial of labour after caesarean (TOLAC). The risk is higher with classical or T-shaped uterine incisions (5-10\%), multiple prior caesareans, short interpregnancy interval, labour induction with prostaglandins (especially misoprostol), and traumatic delivery. Complete rupture involves full-thickness separation through serosa; incomplete rupture (dehiscence) involves separation of the myometrium with intact serosa. Management: emergency caesarean delivery and surgical repair or hysterectomy. Maternal mortality is less than 1\%, but perinatal mortality ranges from 5-20\%.

\medterm{Vacuum-Assisted Delivery}-- Operative vaginal delivery using a vacuum extractor applied to the fetal scalp. Types: soft cup (Kiwi OmniCup, Mityvac, for outlet and low deliveries) and rigid cup (metal Kiwi, for posterior positions and midpelvic deliveries). Steps: confirm position, cup placed over the flexion point (3 cm anterior to the posterior fontanelle), vacuum is created to 600-650 mm Hg, traction is applied during contractions with maternal pushing. Descent should be progressive; if no descent after three pulls, abandon. Cup pops-off: reposition after inspection. Contraindications: face presentation, <34 weeks (fragile skull), unengaged head. Neonatal risks: cephalohematoma, subgaleal hemorrhage (life-threatening, monitor), retinal hemorrhage.

\medterm{Vacuum Delivery}-- Assisted vaginal delivery using a vacuum extractor (cup applied to the fetal scalp with negative pressure) to expedite delivery for prolonged second stage, fetal distress, or maternal indications. Requires full dilation, engaged presenting part, ruptured membranes, and informed consent; contraindicated for face/brow presentation, <34 weeks, suspected fetal coagulopathy or macrosomia with cephalopelvic disproportion. Procedure: cup placement (over the flexion point), traction with contractions/pushing, and careful popping/descent. Complications: fetal scalp laceration, cephalohematoma, subgaleal hemorrhage (serious), retinal hemorrhage, and maternal perineal trauma; failed vacuum converts to forceps/cesarean.

\medterm{Vaginal Birth After Cesarean (VBAC)}-- Vaginal delivery in a woman who has had a prior cesarean delivery. Success rate: 60-80\% in appropriately selected candidates. Contraindications: prior classical or T-incision, prior uterine rupture, vertical uterine incision, more than 2 prior cesareans, and other contraindications to vaginal delivery. Factors favoring success: prior vaginal delivery (especially VBAC), spontaneous labor, non-recurrent indication for prior cesarean. Risk: uterine rupture (0.5-0.9\%, higher with labor induction). Resources: continuous EFM, in-hospital capability for emergency cesarean.

\medterm{Vaginal Delivery (Spontaneous)}-- The spontaneous passage of the fetus through the birth canal without the use of operative instruments (vacuum or forceps). Typically occurs in the second stage of labor with maternal pushing. Includes crowning of the fetal head, delivery of the head (with restitution and external rotation), delivery of the anterior shoulder, delivery of the posterior shoulder, and delivery of the body. The umbilical cord is clamped and cut. The baby is placed on the mother's abdomen or chest for skin-to-skin contact. APGAR scores are assigned at 1 and 5 minutes.

\medterm{Vaginal Intercourse}-- The insertion of the penis into the vagina for sexual pleasure or reproduction. The most common form of sexual intercourse. Vaginal intercourse can lead to pregnancy if sperm are deposited in the vagina. From a medical perspective, it carries risks of STIs (HIV, gonorrhea, chlamydia, HPV, herpes, syphilis), unintended pregnancy, and rarely trauma (vaginal lacerations, especially in postmenopausal women or with inadequate lubrication). Contraception (condoms, hormonal methods, IUDs) and safe sex practices reduce these risks. Normal physiologic changes during vaginal intercourse include vaginal lubrication (Bartholin glands), vasocongestion, and orgasm.

\medterm{Variable Decelerations}-- A common fetal heart rate pattern characterized by an abrupt decrease in FHR (onset to nadir <30 seconds) with a variable shape, typically caused by umbilical cord compression. Variable decelerations are classified as mild (<30 sec duration, >80 bpm), moderate (30-60 sec, 70-80 bpm), or severe (>60 sec, <70 bpm). Recurrent variable decelerations (>50\% of contractions) with moderate variability are reassuring. Management for repetitive/severe variable decels: maternal position change, amnioinfusion, and reduction of oxytocin.

\medterm{VEAL CHOP Mnemonic}-- A memory aid for interpreting fetal heart rate patterns. VEAL: Variable decelerations $\rightarrow$ Cord compression, Early decelerations $\rightarrow$ Head compression, Accelerations $\rightarrow$ OK (fetal well-being), Late decelerations $\rightarrow$ Placental insufficiency. CHOP: Cord compression $\rightarrow$ Change position, Head compression $\rightarrow$ (no intervention needed), OK $\rightarrow$ (reassuring), Placental insufficiency $\rightarrow$ Oxygen. A quick reference for nurses and residents to correlate FHR patterns with their causes and guide initial interventions.

\medterm{Venous Thromboembolism in Pregnancy}-- Pregnancy increases the risk of VTE (DVT and PE) by 5-10 fold due to a hypercoagulable state (increased fibrinogen, factors VII-XII, vWF, decreased protein S, and venous stasis from the gravid uterus). Risk factors: prior VTE, thrombophilia (factor V Leiden, prothrombin G20210A, antithrombin deficiency, protein C/S deficiency), obesity, age >35, and cesarean section. DVT: most commonly left leg (by 3:1), proximal (iliofemoral rather than calf). Diagnosis: compression ultrasonography (DVT), CT pulmonary angiography (PE). Treatment: LMWH (enoxaparin 1 mg/kg BID for full-dose, 40 mg daily for prophylaxis), anticoagulate for entire pregnancy and at least 6 weeks postpartum.

\medterm{Vitamin K Prophylaxis}-- Intramuscular administration of vitamin K (1 mg, 0.5 mg for preterm) to all newborns at birth to prevent vitamin K deficiency bleeding (VKDB, formerly hemorrhagic disease of the newborn). Newborns have low stores of vitamin K (limited placental transfer, sterile gut, low in breast milk). VKDB presents with intracranial, GI, or umbilical bleeding. Oral vitamin K is less effective, especially for late VKDB (2-12 weeks). The IM route is superior and standard in most countries.

\medterm{Vulvar Cancer}-- Malignancy of the vulva, most commonly squamous cell carcinoma (90 percent), related to HPV (younger women, with VIN) or lichen sclerosus (older women); also melanoma, basal cell, Paget disease, Bartholin gland adenocarcinoma. Risk: age, HPV infection, smoking, immunosuppression, chronic vulvar dermatoses. Features: persistent vulvar lump, ulcer, itching, pain, bleeding, and enlarged inguinal nodes; often diagnosed late. Diagnosis: biopsy. Management: surgical excision (wide local excision, vulvectomy) with sentinel node staging/inguinal lymphadenectomy for early; chemoradiation for advanced or node-positive; HPV vaccination prevents many cases. Prognosis depends on stage, nodal involvement.

\medterm{Water Birth (Underwater Delivery)}-- Delivery of a baby in a tub of warm water. Proponents claim: reduced pain, fewer perineal lacerations, and a gentler transition for the baby. Water birth is used in many birth centers and some hospitals. Criteria: low-risk pregnancy, term, clear fluid, and no signs of fetal distress. Risks: maternal or neonatal infection, newborn aspiration (the baby does not breathe until exposed to air), umbilical cord avulsion, difficulty assessing blood loss, and water embolism. The AAP and ACOG state that water birth should be offered only in the context of clinical trials.

\medterm{Weight Gain in Pregnancy}-- Recommended gestational weight gain according to the Institute of Medicine (IOM) guidelines, based on prepregnancy BMI. Underweight (BMI <18.5): 28-40 lbs (12.5-18 kg). Normal weight (BMI 18.5-24.9): 25-35 lbs (11.5-16 kg). Overweight (BMI 25-29.9): 15-25 lbs (7-11.5 kg). Obese (BMI \(\ge\)30): 11-20 lbs (5-9 kg). Twin gestations: higher weight gain recommended. Inadequate weight gain is associated with IUGR and preterm birth; excessive gain is associated with macrosomia, gestational diabetes, cesarean delivery, and postpartum weight retention.

\medterm{Wharton Jelly}-- A gelatinous substance (mucopolysaccharide matrix) surrounding the umbilical cord vessels, composed primarily of hyaluronic acid and chondroitin sulfate. It protects the umbilical vessels from compression and kinking during fetal movement and uterine contractions. Its elastic properties allow the cord to maintain blood flow despite external pressure. Depletion of Wharton jelly is associated with oligohydramnios and cord accidents.

\medterm{Whites (Leucorrhoea)}-- A non-bloody, physiological or pathological vaginal discharge. Physiological leucorrhoea: normal, clear to whitish, odourless discharge that varies with the menstrual cycle (increased at ovulation, pre-menstrually, during pregnancy, sexual arousal, and with oral contraceptive use). It consists of desquamated epithelial cells, cervical mucus (from endocervical glands), vaginal flora (lactobacilli producing lactic acid, pH 3.8-4.5), and transudate. Pathological leucorrhoea: abnormal discharge in color, odour, or consistency. Causes: vulvovaginal candidiasis (white, thick, curd-like discharge, pruritus, erythema, pH normal 4-4.5, KOH-positive for pseudohyphae), bacterial vaginosis (thin, gray-white, fishy odour, elevated pH >4.5, clue cells on microscopy, positive whiff test with KOH), Trichomonas vaginalis (yellow-green, frothy, malodorous, pruritus, dysuria, strawberry cervix, pH >5, motile trichomonads on wet mount), Chlamydia trachomatis (mucopurulent cervicitis, yellow discharge, cervical friability), Neisseria gonorrhoeae (purulent discharge, PID risk), and atrophic vaginitis (thin, watery, sometimes blood-tinged, pH >5, in menopause). Diagnosis: history, speculum examination, pH testing, wet mount microscopy, KOH preparation, NAAT for Chlamydia/gonorrhea/Trichomonas. Management: treat underlying infection (antifungals for Candida, metronidazole for BV and Trichomonas, antibiotics for STIs). Recurrent or persistent leucorrhoea requires evaluation for cervical pathology and retained foreign body (tampon).

\medterm{Womb}-- The common term for the uterus, the muscular organ in the female pelvis where a fetus develops during pregnancy. The womb is pear-shaped, approximately 7.5 cm long, 5 cm wide, and 2.5 cm thick in the non-pregnant state. It consists of three layers: perimetrium (outer serosa), myometrium (thick muscular layer, responsible for uterine contractions during labor), and endometrium (inner lining, shed during menstruation and the site of implantation). The womb expands dramatically during pregnancy, increasing in capacity from 10 mL to 5 L.

\medterm{Woods Corkscrew Maneuver}-- An internal rotation maneuver for shoulder dystocia where the operator inserts a hand and rotates the fetal shoulders 180 degrees by applying pressure to the posterior aspect of the anterior shoulder. This disimpacts the shoulder from behind the symphysis pubis. Can be repeated if the shoulder re-impacts in the new position. Used when McRoberts maneuver and suprapubic pressure fail to resolve the dystocia.

\medterm{Yeast Infection}-- Vulvovaginal candidiasis: infection by Candida albicans (most common), causing vulvar and vaginal itching, burning, thick white 'cottage cheese' discharge, dyspareunia, and erythema; more common in pregnancy, diabetes, antibiotic use, and immunosuppression. Diagnosis: clinical; wet mount (budding yeast/hyphae), KOH, or culture. Management: topical azoles (clotrimazole, miconazole) or oral fluconazole (single/3-day dose) for uncomplicated; complicated (recurrent, severe, non-albicans, immunocompromised) needs longer therapy, culture-directed treatment, and maintenance/suppression; treat predisposing factors. Distinguish from bacterial vaginosis and trichomoniasis.

\medterm{Yolk Sac}-- The first extraembryonic structure visible on ultrasound at 5-6 weeks gestation. It provides early nutrition to the embryo before the placental circulation is established, produces blood cells until the fetal liver takes over, and forms the primitive gut. A normal yolk sac is 3-6 mm in diameter. Abnormal size (large >6 mm or small <2 mm) or shape is associated with early pregnancy failure. The yolk sac degenerates by the end of the first trimester.

\medterm{Zavanelli Maneuver}-- The final, rarely performed maneuver for an irreducible shoulder dystocia: the fetal head is replaced into the vagina (cephalic replacement) and the patient undergoes emergency cesarean delivery. Technique: the head is flexed and pushed back into the vagina, then into the birth canal under tocolysis (terbutaline, magnesium, or general anesthesia for uterine relaxation). High rates of fetal morbidity and mortality. The maneuver is used only when all other maneuvers have failed.

\medterm{Zygote}-- The single diploid cell formed by the fusion of male and female gametes (spermatozoon and secondary oocyte) at fertilization. The zygote contains 46 chromosomes (23 from each parent) and represents the beginning of a new, genetically unique individual. After fertilization, the zygote undergoes a series of mitotic cell divisions (cleavage) as it travels through the fallopian tube toward the uterine cavity. The first cleavage produces two blastomeres, then four, eight, and sixteen cells (morula), which then develops into a blastocyst (with inner cell mass $\rightarrow$ embryo and outer trophoblast $\rightarrow$ placenta) and implants in the endometrium approximately 5-6 days after fertilization. The zygote stage lasts approximately 24-30 hours, after which it becomes a 2-cell embryo. Abnormal zygote formation can result from polyspermy (fertilization by multiple sperm, resulting in triploidy), parthenogenesis (activation of the oocyte without sperm, not viable in humans), or chromosomal nondisjunction (leading to aneuploidy). In assisted reproductive technology (IVF/ICSI), the zygote is assessed for pronuclear morphology (2 pronuclei = normal fertilization) before embryo culture and transfer. Physiological leucorrhoea: normal, clear to whitish, odourless discharge that varies with the menstrual cycle (increased at ovulation, pre-menstrually, during pregnancy, sexual arousal, and with oral contraceptive use). It consists of desquamated epithelial cells, cervical mucus (from endocervical glands), vaginal flora (lactobacilli producing lactic acid, pH 3.8-4.5), and transudate. Pathological leucorrhoea: abnormal discharge in color, odour, or consistency. Causes: vulvovaginal candidiasis (white, thick, curd-like discharge, pruritus, erythema, pH normal 4-4.5, KOH-positive for pseudohyphae), bacterial vaginosis (thin, gray-white, fishy odour, elevated pH >4.5, clue cells on microscopy, positive whiff test with KOH), Trichomonas vaginalis (yellow-green, frothy, malodorous, pruritus, dysuria, strawberry cervix, pH >5, motile trichomonads on wet mount), Chlamydia trachomatis (mucopurulent cervicitis, yellow discharge, cervical friability), Neisseria gonorrhoeae (purulent discharge, PID risk), and atrophic vaginitis (thin, watery, sometimes blood-tinged, pH >5, in menopause). Diagnosis: history, speculum examination, pH testing, wet mount microscopy, KOH preparation, NAAT for Chlamydia/gonorrhea/Trichomonas. Management: treat underlying infection (antifungals for Candida, metronidazole for BV and Trichomonas, antibiotics for STIs). Recurrent or persistent leucorrhoea requires evaluation for cervical pathology and retained foreign body (tampon).


\section{Gynecologic Oncology \& Disorders}
\medterm{Adenomyosis}-- A condition in which endometrial glands and stroma are present within the myometrium (uterine muscle). Most common in multiparous women in their 40s and 50s. Presents with heavy menstrual bleeding (menorrhagia), severe dysmenorrhea, chronic pelvic pain, and uterine enlargement. Diagnosis: transvaginal ultrasound shows diffuse or focal myometrial thickening with heterogeneous echogenicity; MRI is more specific (junctional zone thickening >12 mm). Treatment: NSAIDs, hormonal contraceptives, levonorgestrel IUD, GnRH agonists, endometrial ablation, and hysterectomy for definitive management.

\medterm{Anal Intercourse}-- The insertion of the penis into the anus and rectum for sexual pleasure. The anal canal is rich in nerve endings, making it a sensitive erogenous zone. The rectum lacks natural lubrication, so artificial lubricants are essential to prevent mucosal trauma. Risks: STI transmission is higher than with vaginal intercourse (especially HIV, HPV, hepatitis A/B, gonorrhea, chlamydia, and syphilis) due to the fragile rectal mucosa and increased HIV target cells (CD4+ T-cells) in the rectal tissue. Condom use significantly reduces STI risk. Other risks: anal fissures, fecal incontinence (rare with healthy sphincter), and foreign body retention. Anal intercourse does not cause hemorrhoids.

\medterm{Breast Health}-- Breast health involves routine self-awareness, clinical examinations, and mammographic screening to detect abnormalities such as lumps, pain, or skin changes that may indicate benign conditions or breast cancer. Risk factors include genetics, age, hormone exposure, and lifestyle factors. Early detection through regular screening significantly improves treatment outcomes and survival.

\medterm{Cervical Dysplasia}-- Abnormal growth of cervical epithelial cells, representing a premalignant condition that may progress to cervical cancer if untreated. Also known as cervical intraepithelial neoplasia (CIN). CIN is caused by persistent infection with high-risk human papillomavirus (HPV) types, most commonly HPV 16 and 18. Grading: CIN 1 (mild dysplasia, involves the lower one-third of the epithelium, high rate of spontaneous regression), CIN 2 (moderate dysplasia, lower two-thirds), and CIN 3 (severe dysplasia/carcinoma in situ, full-thickness involvement). Most CIN 1 lesions regress spontaneously; CIN 2 and 3 require treatment. Diagnosis: abnormal Pap smear followed by colposcopy with directed biopsy and endocervical curettage. Treatment: CIN 1 may be observed; CIN 2/3 is treated with excisional procedures (loop electrosurgical excision procedure, cold knife conization) or ablative therapy (cryotherapy, laser ablation) depending on lesion size, location, and patient factors. Follow-up is essential as recurrence can occur. Prevention: HPV vaccination.

\medterm{Clear Margins}-- A histological finding indicating that no abnormal cells (e.g., cervical intraepithelial neoplasia or cancer cells) are present at the edge of a tissue specimen removed during a procedure such as LLETZ, cone biopsy, or surgical excision. Clear margins suggest that the entire abnormal area has been removed, reducing the risk of persistent or recurrent disease. Positive margins (abnormal cells at the edge) indicate incomplete excision and require closer follow-up or repeat treatment. Clear margins are the goal of excisional treatment for cervical cell changes and other gynecologic conditions.

\medterm{Diathermy}-- The use of high-frequency electrical current to generate heat for surgical purposes, including cutting tissue, coagulating blood vessels to stop bleeding (hemostasis), and destroying abnormal tissue (electrocautery, electrocoagulation). In gynecologic surgery, diathermy is used in laparoscopic procedures (ovarian drilling for PCOS, endometriosis ablation, myomectomy), hysteroscopic surgery (endometrial ablation, polypectomy, fibroid resection), and colposcopy (treatment of cervical intraepithelial neoplasia). Monopolar diathermy: current flows from electrode through patient to grounding pad. Bipolar diathermy: current flows between two forceps tips, safer for use near delicate structures.

\medterm{Dysmenorrhea}-- Painful menstruation with cramping pelvic pain. Primary dysmenorrhea: no underlying pathology, associated with prostaglandin release, starts within 3 years of menarche. Secondary dysmenorrhea: due to underlying pathology (endometriosis, fibroids, adenomyosis, PID). Treatment: NSAIDs (ibuprofen, naproxen), hormonal contraceptives (pills, IUD, implant), and treatment of underlying cause. Severe cases may require GnRH agonists or surgical intervention.

\medterm{Dysmenorrhoea}-- Painful menstruation, classified as primary (no identifiable pelvic pathology) or secondary (underlying structural or inflammatory condition). Primary dysmenorrhoea is caused by prostaglandin-induced uterine contractions and occurs in the first 2-3 years after menarche, typically starting 6-12 hours before menstruation and lasting 1-3 days. Secondary dysmenorrhoea is associated with endometriosis, adenomyosis, fibroids, or pelvic inflammatory disease. Treatment: NSAIDs (inhibit prostaglandin synthesis), combined oral contraceptives (reduce menstrual volume and prostaglandin production), Mirena IUS, and treatment of underlying cause.

\medterm{Electrocoagulation}-- A surgical technique using high-frequency electrical current to destroy tissue or stop bleeding by coagulating proteins, also known as diathermy or cauterization. In gynecology, electrocoagulation is used in laparoscopic treatment of endometriosis (ablation of implants), polycystic ovary syndrome (ovarian drilling), and hysteroscopic endometrial ablation. It may be monopolar or bipolar. Complications are rare but include thermal injury to adjacent organs (bowel, ureter, bladder) and infection.

\medterm{Endometrial Cancer}-- The most common gynecologic malignancy in developed countries, arising from the lining of the uterine cavity, typically in postmenopausal women. Type I (endometrioid) carcinomas account for 80 percent of cases and are estrogen-driven, associated with endometrial hyperplasia, obesity, diabetes, polycystic ovary syndrome, unopposed estrogen therapy, and tamoxifen use; they are generally low-grade with a favorable prognosis. Type II (non-endometrioid) carcinomas (serous, clear cell, carcinosarcoma) are not estrogen-driven, arise in atrophic endometrium, and are aggressive with a poor prognosis. The lifetime risk is approximately 3 percent in the general population. Clinical presentation: postmenopausal bleeding is the cardinal symptom (occurring in 90 percent of cases), along with abnormal uterine bleeding in premenopausal women, watery or bloody vaginal discharge, and pelvic pain in advanced disease. Diagnosis is by endometrial biopsy or dilation and curettage, with transvaginal ultrasound showing endometrial thickness greater than 4-5 mm in postmenopausal women requiring further evaluation. Histologic grading (FIGO grade 1-3 based on glandular architecture), subtype determination, and molecular classification (POLE ultramutated, microsatellite instability-high, copy number low, copy number high) guide prognosis and treatment. Staging is surgical according to the FIGO system, based on depth of myometrial invasion, cervical stromal involvement, adnexal spread, lymph node metastasis, and peritoneal cytology. Treatment: total hysterectomy with bilateral salpingo-oophorectomy is the mainstay; lymph node assessment is performed for high-risk histologies. Adjuvant radiotherapy and/or chemotherapy are indicated for high-stage or high-grade disease. Prognosis is generally favourable due to early presentation with postmenopausal bleeding; five-year survival exceeds 95\% for stage I disease.

\medterm{Endometriosis}-- A chronic, estrogen-dependent gynecologic disorder characterized by the presence of endometrial-like tissue (glands and stroma) outside the uterine cavity, most commonly implanted on the pelvic peritoneum, ovaries, uterosacral ligaments, and rectovaginal septum. The ectopic endometrial tissue responds to hormonal stimulation, proliferating and shedding cyclically, causing local inflammation, adhesion formation, and fibrosis. The prevalence is 6-10 percent of women of reproductive age, rising to 35-50 percent in women with infertility and chronic pelvic pain. The exact pathogenesis remains unclear, but leading theories include retrograde menstruation (Sampson theory, with menstrual debris transported through the fallopian tubes into the peritoneal cavity), coelomic metaplasia (transformation of peritoneal mesothelium into endometrial tissue), lymphatic or hematogenous dissemination, and immune dysregulation (impaired clearance of endometrial cells from the peritoneal cavity). Genetic factors contribute, with a 6-9 fold increased risk in first-degree relatives. Clinical presentation includes chronic pelvic pain (often cyclic but may become continuous), severe dysmenorrhea, dyspareunia (painful intercourse), infertility, and dyschezia (painful defecation). Physical examination may reveal tender nodularity in the posterior vaginal fornix, fixed retroverted uterus, and adnexal masses (endometriomas). Diagnosis is definitively made by direct visualization at laparoscopy, with histological confirmation of endometrial glands and stroma in biopsy specimens. Treatment includes analgesics (NSAIDs), hormonal therapy (combined oral contraceptives, progestins, GnRH agonists, levonorgestrel IUS), and surgical excision or ablation of implants. Fertility management ranges from expectant management to assisted reproductive technologies.

\medterm{Fibrocystic Breast Changes}-- A benign condition characterized by lumpy, painful breast tissue, most common in women aged 20-50, driven by hormonal fluctuations during the menstrual cycle (estrogen and progesterone). Microscopic features include stromal fibrosis, cyst formation (macrocysts and microcysts), and adenosis (increased glandular tissue). Palpable breast lumps may be tender and fluctuate in size with the menstrual cycle. Diagnosis relies on clinical breast exam, mammography (showing dense breast tissue and cysts), and ultrasound (anechoic simple cysts), with biopsy (core needle or fine-needle aspiration) reserved for suspicious or complex cysts or solid areas to exclude malignancy. Fibrocystic changes do not increase breast cancer risk except when associated with proliferative changes with atypia (columnar cell change with atypia or atypical ductal hyperplasia), which warrants increased surveillance. Management includes supportive bras, NSAIDs, oral contraceptives (suppressing hormonal fluctuations), and avoiding dietary methylxanthines (caffeine) if symptomatic, though evidence for dietary modification is weak.

\medterm{Fibroids, Cysts, and Polyps}-- Benign growths that can develop in the reproductive organs. Fibroids are smooth muscle tumors of the uterus; cysts are fluid-filled sacs that commonly form on the ovaries; and polyps are tissue overgrowths in the endometrium or cervix. Although typically noncancerous, they can cause symptoms such as pain, abnormal bleeding, and fertility issues, often requiring monitoring or surgical removal.

\medterm{Fimbriae}-- The finger-like projections at the distal end of each fallopian tube that sweep over the ovary during ovulation to capture the released egg and guide it into the fallopian tube. The fimbriae are lined with ciliated epithelium that creates a current to draw the egg into the tube. At the center of the fimbriae is the abdominal ostium, the opening into the fallopian tube. Damage to the fimbriae from infection (pelvic inflammatory disease, chlamydia), endometriosis, or surgery can impair oocyte pickup and cause infertility.

\medterm{Galactocele}-- A milk-filled retention cyst of the breast, occurring during lactation or after weaning. Results from obstruction of a milk duct. Presents as a painless, fluctuant breast mass. Diagnosis by ultrasound showing a well-defined, hypoechoic or anechoic mass. Treatment includes aspiration. May become infected, requiring antibiotics. Recurrence is uncommon after complete drainage.

\medterm{Genetic Sonogram}-- A targeted ultrasound performed at 18-22 weeks that combines standard anatomic survey with specific soft markers for aneuploidy. Findings assessed: nuchal fold thickened (\(\ge\)6 mm), echogenic intracardiac focus (EIF), echogenic bowel, renal pyelectasis (\(\ge\)4 mm anterior-posterior), short femur/humerus, choroid plexus cyst, and sandal gap. The presence of multiple soft markers increases the risk of aneuploidy. The genetic sonogram is used for risk adjustment when combined with maternal age and serum screening (likelihood ratio for each marker). Markers are nonspecific and many are normal variants in euploid fetuses.

\medterm{Genitals}-- The external and internal reproductive organs. Female genitals: external (vulva — mons pubis, labia majora, labia minora, clitoris, vaginal opening, Bartholin glands) and internal (vagina, cervix, uterus, fallopian tubes, ovaries). Male genitals: external (penis, scrotum) and internal (testes, epididymis, vas deferens, seminal vesicles, prostate, bulbourethral glands). The genitals develop from common embryological structures (genital tubercle, urogenital sinus, Wolffian and Müllerian ducts) under the influence of sex-determining genes and hormones.

\medterm{Gonadotrophins}-- See endocrinology.

\medterm{HRT}-- Hormone Replacement Therapy, the use of estrogen (with or without progesterone) to treat menopausal symptoms caused by declining ovarian hormone production. HRT is the most effective treatment for vasomotor symptoms (hot flushes, night sweats), prevents bone loss (osteoporosis), and improves urogenital atrophy. Routes: oral tablets, transdermal patches/gels, vaginal preparations. Estrogen alone is used in women after hysterectomy; combined estrogen-progestogen is used in women with an intact uterus to prevent endometrial hyperplasia/cancer. The timing hypothesis: HRT initiated within 10 years of menopause has cardiovascular benefits; late initiation increases cardiovascular and thromboembolic risks.

\medterm{Hydrocele}-- An accumulation of fluid within the tunica vaginalis surrounding the testicle, causing scrotal swelling.

\medterm{Hysterosalpingo-Contrast-Sonography}-- An ultrasound-based technique for evaluating the uterine cavity and fallopian tube patency. Sterile saline or contrast medium is injected through the cervix while transvaginal ultrasound is performed. Advantages over HSG: no ionizing radiation, can be performed in clinic, assesses the uterine cavity (polyps, fibroids, adhesions) and tubal patency simultaneously. Disadvantages: less detailed tubal assessment than X-ray HSG, operator-dependent. Used in infertility workup to screen for tubal occlusion before proceeding to IVF.

\medterm{Hysterosalpingogram (HSG)}-- An X-ray procedure used to evaluate the uterine cavity and fallopian tube patency. A radiopaque contrast medium is injected through the cervix, and serial X-ray images are taken as the contrast fills the uterine cavity and spills from the fallopian tubes. Indications: infertility workup (tubal factor), recurrent miscarriage (uterine cavity assessment), and evaluation after tubal sterilization or tubal surgery. Findings: normal (smooth uterine cavity, free spill from both tubes), tubal occlusion (proximal or distal), hydrosalpinx (dilated, blocked tube), uterine filling defects (polyps, fibroids, adhesions, septate uterus).

\medterm{Infertility}-- Inability to conceive after 12 months of regular unprotected intercourse (6 months if woman >35). Affects 15\% of couples. Female causes: ovulation disorders (PCOS, hypothalamic amenorrhea), tubal factor (blockage, PID, endometriosis), uterine (fibroids, polyps, adhesions), and diminished ovarian reserve. Male causes: low sperm count/motility, varicocele, and hormonal dysfunction. Evaluation: ovarian reserve testing, semen analysis, hysterosalpingogram, and hormonal profiling. Treatment: ovulation induction, IUI, IVF, and donor gametes/embryos.

\medterm{Leiomyoma (Fibroid)}-- A benign smooth muscle tumor of the uterus, arising from the myometrium. See Fibroids (Uterine Leiomyomas).

\medterm{LLETZ}-- Large Loop Excision of the Transformation Zone, also called LEEP (Loop Electrosurgical Excision Procedure). A common outpatient procedure used to diagnose and treat cervical intraepithelial neoplasia (CIN). A thin wire loop carrying an electrical current is used to remove the transformation zone of the cervix (where cell changes typically occur). The procedure takes approximately 5-10 minutes under local anesthesia. The removed tissue is sent for histology. Complications: bleeding, infection, and cervical stenosis. LLETZ does not significantly affect fertility but may slightly increase the risk of preterm birth in future pregnancies.

\medterm{Menorrhagia}-- Excessive or prolonged menstrual bleeding (>80 mL blood loss or >7 days duration). Causes: uterine fibroids, adenomyosis, endometrial polyps, coagulation disorders, thyroid dysfunction, and anovulation. Presents with heavy bleeding, flooding, passage of clots, and iron deficiency anemia. Diagnosis: CBC, ferritin, TSH, transvaginal ultrasound, saline infusion sonography, and endometrial biopsy if indicated. Treatment: NSAIDs, tranexamic acid, hormonal contraceptives, levonorgestrel IUD, endometrial ablation, and hysterectomy in severe cases.

\medterm{Mesh}-- A synthetic or biological implant used to provide structural support in pelvic organ prolapse and stress urinary incontinence surgery. Types: polypropylene (most common, mid-urethral slings for SUI, vaginal mesh for prolapse), biological (porcine, bovine, or cadaveric dermis). Controversy: transvaginal mesh for prolapse has been associated with complications including mesh erosion (exposure into vagina), chronic pain, dyspareunia, infection, and organ perforation. The NHS stopped routine use of transvaginal mesh for prolapse in 2018. Mid-urethral slings for SUI (retropubic and transobturator tapes) remain in use with careful patient selection and counseling.

\medterm{Metrorrhagia}-- Uterine bleeding at irregular intervals, particularly between expected menstrual periods. Causes: anovulation, endometrial polyps, submucosal fibroids, endometrial hyperplasia/cancer, cervical pathology, and breakthrough bleeding with hormonal contraceptives. Diagnosis: transvaginal ultrasound, endometrial biopsy, and hysteroscopy. Treatment depends on the underlying cause and patient age.

\medterm{Newborn Metabolic Screen}-- A dried blood spot test (heel stick) collected 24-48 hours after birth to screen for congenital disorders before symptoms develop. The panel varies by country/state but typically includes: phenylketonuria (PKU), congenital hypothyroidism, cystic fibrosis, sickle cell disease, medium-chain acyl-CoA dehydrogenase deficiency (MCADD), maple syrup urine disease (MSUD), galactosemia, severe combined immunodeficiency (SCID), and many others (50-80 conditions in some states). Positive results require confirmatory testing.

\medterm{Oophorectomy}-- Surgical removal of one (unilateral) or both (bilateral) ovaries. Indications: ovarian masses (benign or malignant), ovarian torsion, endometriosis, pelvic inflammatory disease (abscess), risk-reducing surgery (BRCA1/BRCA2 mutation carriers—risk-reducing salpingo-oophorectomy reduces ovarian cancer risk by 80--96\% and breast cancer risk by 50\%), and as part of treatment for hormone-sensitive cancers (breast cancer, endometrial cancer). Bilateral oophorectomy causes surgical menopause (immediate loss of oestrogen, progesterone, and testosterone production), with symptoms (hot flushes, vaginal atrophy, mood changes, decreased libido) and long-term health effects (increased cardiovascular disease, osteoporosis, cognitive decline). Hormone replacement therapy mitigates these risks in women aged <50 without contraindications. Oophorectomy is often performed with salpingectomy (salpingo-oophorectomy); bilateral salpingo-oophorectomy plus hysterectomy is common for benign indications and ovarian cancer staging.

\medterm{Oral Sex (Fellatio, Cunnilingus)}-- Oral stimulation of the genitals for sexual pleasure. Fellatio: oral stimulation of the penis. Cunnilingus: oral stimulation of the vulva and clitoris. Analingus (rimming): oral stimulation of the anus. Oral sex carries lower STI risk than unprotected vaginal or anal intercourse, but STIs can still be transmitted (HSV-1/2, gonorrhea (pharyngeal), chlamydia (pharyngeal), syphilis, HPV (oropharyngeal cancer risk with HPV-16), and HIV (lower risk than anal/vaginal). Pharyngeal gonorrhea and chlamydia are often asymptomatic but can be transmitted to partners. Dental dams and condoms reduce risk. Oral sex is a common cause of primary HSV-1 infection (cold sores transmitted to the genitals or vice versa).

\medterm{Ovum}-- The female reproductive cell (egg) after ovulation and before fertilization. The ovum is a large, non-motile cell (~100 \(\mu\)m diameter) surrounded by the zona pellucida (glycoprotein layer that mediates sperm binding and prevents polyspermy after fertilization) and the corona radiata (follicular cells). The ovum is arrested in metaphase II of meiosis at ovulation. Upon sperm penetration, the ovum completes meiosis II (producing a mature oocyte and a second polar body), and the male and female pronuclei fuse (syngamy) to form a diploid zygote. The ovum is viable for approximately 12--24 hours after ovulation; fertilization must occur within this window. Ovum donation is used in assisted reproduction for women with poor oocyte quality, premature ovarian insufficiency, or genetic disorders.

\medterm{Pessaries}-- Devices inserted into the vagina to support pelvic organs in cases of pelvic organ prolapse (cystocele, rectocele, uterine prolapse). Types: ring pessary (most common, with or without supportive membrane), shelf pessary (for severe prolapse), cube pessary (suction-based), and Gellhorn pessary (for advanced prolapse). Pessaries are fitted by a trained clinician and can be self-managed for removal, cleaning, and reinsertion. Indications: women who decline or are unsuitable for surgery, women awaiting surgery, or those who wish to conceive (pessary avoids surgical risks).

\medterm{Pessary}-- A removable device placed in the vagina to provide structural support for pelvic organ prolapse (uterine prolapse, cystocele, rectocele) or stress urinary incontinence. Pessaries are available in various shapes and sizes: ring pessary (most common—used for mild to moderate prolapse), Gellhorn pessary (for severe prolapse—has a knob that sits above the levator plate), cube pessary (suction-cup shaped, for advanced prolapse), and incontinence pessary (with a knob that supports the urethrovesical junction). Pessaries are fitted by a trained clinician; the correct size prevents discomfort or expulsion. Patients are taught to remove, clean, and reinsert the pessary (if they can) or return every 3--6 months for cleaning and inspection. Complications: vaginal discharge, ulceration, bleeding, and rarely fistula formation. Oestrogen cream may reduce vaginal erosion in postmenopausal women.

\medterm{Prolapse of the Uterus (Uterine Prolapse)}-- Descent of the uterus into the vaginal canal due to weakening of pelvic floor supports, most commonly the uterosacral and cardinal ligaments. Graded by the POP-Q system (Pelvic Organ Prolapse Quantification): Stage 0 (no prolapse), Stage I (descent >1 cm above hymen), Stage II (descent to within 1 cm of hymen), Stage III (descent >1 cm below hymen), Stage IV (complete eversion of vagina). Risk factors: vaginal childbirth (especially multiple, prolonged, or traumatic deliveries), aging, menopause (oestrogen deficiency), obesity, chronic cough, constipation, and heavy lifting. Symptoms: vaginal bulge or pressure, sensation of something falling out, urinary incontinence or retention, incomplete bladder emptying, and difficulty with defecation. Diagnosis: clinical, confirmed on pelvic examination with Valsalva manoeuvre. Treatment: observation for asymptomatic mild prolapse; pelvic floor muscle training (Kegel exercises); vaginal pessary (ring or shelf pessary for symptomatic relief); surgery (vaginal hysterectomy with uterosacral ligament suspension, sacrospinous fixation, or sacrocolpopexy for apical prolapse). Concomitant anti-incontinence procedures may be performed.

\medterm{Rectocele}-- A weakening of the vaginal wall that allows the rectum to bulge into the vagina.

\medterm{Sporadic}-- See genetics.

\medterm{Teratoma}-- A type of germ cell tumor composed of tissues derived from two or three germ layers (ectoderm, endoderm, mesoderm), reflecting pluripotent differentiation. Teratomas are classified as mature (benign—well-differentiated tissues, most common in children and young adults) or immature (malignant—primitive/embryonal tissues, higher risk of recurrence and metastasis). Sites: ovaries (most common—dermoid cyst/mature cystic teratoma, usually benign), testes (prepubertal boys—benign; postpubertal men—may contain malignant germ cell elements), sacrococcygeal region (most common congenital tumor in neonates, 1 in 35,000 births), mediastinum (anterior mediastinal germ cell tumor), and intracranial (pineal region, suprasellar). Ovarian dermoid cyst: contains hair, teeth, sebum, bone, and thyroid tissue, often discovered incidentally on ultrasound. Treatment: surgical resection (ovarian cystectomy or oophorectomy for ovarian teratoma). Mature teratoma has an excellent prognosis; immature teratoma requires staging, chemotherapy (BEP: bleomycin, etoposide, cisplatin), and surveillance. Sacrococcygeal teratoma in neonates requires early resection (risk of malignant transformation and bleeding).

\medterm{Urethracele}-- A herniation or prolapse of the urethra into the anterior vaginal wall, often occurring together with a cystocele (bladder prolapse). Urethracele is caused by weakening of the pubocervical fascia, often from childbirth, aging, or increased intra-abdominal pressure. Symptoms: stress urinary incontinence (leakage with cough, sneeze, exercise), urinary urgency, and a vaginal bulge. Diagnosis: clinical (pelvic examination). Treatment: pelvic floor physiotherapy, pessary, or surgical repair (anterior colporrhaphy).

\medterm{Urodynamics}-- A set of tests that assess the function of the lower urinary tract (bladder and urethra). Urodynamic studies include: uroflowmetry (flow rate), filling cystometry (bladder pressure during filling, assesses detrusor overactivity and compliance), pressure-flow study (voiding phase, assesses obstruction), and leak point pressure (stress incontinence). Indicated when surgery is planned for stress incontinence, when the diagnosis is unclear, or when mixed incontinence is suspected. Urodynamics helps distinguish genuine stress incontinence from detrusor overactivity.

\medterm{Uterine Fibroid (Leiomyoma)}-- Benign smooth muscle tumors of the uterus, very common (70-80\% of women by age 50). Submucosal: most symptomatic, cause heavy bleeding. Intramural: most common. Subserosal: cause pressure symptoms. Symptoms: menorrhagia, pelvic pain/pressure, urinary frequency, and reproductive dysfunction (infertility, miscarriage). Diagnosis: ultrasound, MRI. Treatment: symptomatic management (NSAIDs, tranexamic acid), hormonal contraceptives, levonorgestrel IUD, GnRH agonists, myomectomy, uterine artery embolization, and hysterectomy.

\medterm{Uterine Fibroids}-- See Uterine Fibroid (Leiomyoma).

\medterm{Vaginal Dryness and Atrophy (Genitourinary Syndrome of Menopause)}-- Vaginal dryness and atrophy, collectively termed genitourinary syndrome of menopause, result from declining estrogen levels that cause thinning, drying, and loss of elasticity of the vaginal tissues. Symptoms include vaginal dryness, irritation, burning, dyspareunia, and urinary symptoms such as frequency and recurrent infections. Treatment options include over-the-counter vaginal lubricants and moisturizers, prescription topical estrogen therapy in various formulations, and nonhormonal options such as ospemifene or vaginal laser therapy.

\medterm{Vaginismus}-- Involuntary spasms of the vaginal muscles.

\medterm{Vaginitis}-- Inflammation of the vaginal mucosa, one of the most common gynecologic complaints. Types: bacterial vaginosis (BV: thin, fishy-odor discharge, clue cells, pH >4.5), vulvovaginal candidiasis (thick, white, cottage-cheese discharge, pruritus, pH normal), and trichomoniasis (frothy, green-yellow discharge, strawberry cervix, pH >4.5, motile trichomonads). Atrophic vaginitis: postmenopausal, due to estrogen deficiency, thin vaginal mucosa, and dyspareunia. Diagnosis: vaginal pH, wet mount microscopy, KOH preparation, and NAAT for trichomonas. Treatment: metronidazole or clindamycin for BV; topical azoles or oral fluconazole for candidiasis; metronidazole or tinidazole for trichomoniasis; topical estrogen for atrophic vaginitis.

\medterm{Vulvar vestibulitis}-- Pain and inflammation of the vulvar vestibule.

\medterm{Vulvitis}-- Inflammation of the vulva (external female genitalia), often presenting with pruritus, burning, irritation, erythema, edema, and discharge. Vulvitis is rarely isolated and is most commonly part of vulvovaginitis (inflammation of both vulva and vagina). Causes: infections (Candida albicans — most common, Trichomonas vaginalis, bacterial vaginosis, herpes simplex, human papillomavirus), irritants (soaps, detergents, douches, perfumed hygiene products, sanitary pads, synthetic underwear, chlorinated water), allergens (latex condoms, spermicides, topical medications), dermatologic conditions (lichen sclerosus, lichen planus, psoriasis, atopic dermatitis, contact dermatitis), and systemic conditions (diabetes mellitus, immunosuppression, oestrogen deficiency in menopause causing atrophic vulvovaginitis). Diagnosis: history (symptom onset, triggers, sexual history), examination (vulvar erythema, fissures, discharge, atrophy, excoriations, lichenification), swabs for culture, pH testing, and KOH preparation. Management: treat underlying cause (antifungals for Candida, antibiotics for bacterial infection, topical corticosteroids for inflammatory dermatoses, oestrogen cream for atrophic vulvitis). General measures: avoidance of irritants, sitz baths, loose cotton underwear, and good vulvar hygiene.


\section{Reproductive Endocrinology \& Infertility}
\medterm{Abnormal Uterine Bleeding}-- Uterine bleeding that is abnormal in frequency, regularity, duration, or volume, occurring outside of normal menstruation or in women who are not pregnant. The PALM-COEIN classification system categorizes causes into structural (PALM: polyps, adenomyosis, leiomyomas, malignancy/hyperplasia) and non-structural (COEIN: coagulopathy, ovulatory dysfunction, endometrial, iatrogenic, not otherwise classified). Ovulatory dysfunction is the most common cause in reproductive-age women, resulting from anovulation due to polycystic ovary syndrome, thyroid disorders, hyperprolactinemia, hypothalamic dysfunction, or perimenopause. Structural causes include endometrial polyps (benign overgrowths of endometrial glands and stroma), uterine leiomyomas (fibroids, benign smooth muscle tumors), adenomyosis (endometrial tissue within the myometrium), and endometrial hyperplasia or carcinoma (particularly in postmenopausal women).

\medterm{Bladder Exstrophy}-- A rare congenital abdominal wall defect where the bladder is exposed on the outside of the body due to failure of the lower abdominal wall and bladder to close during embryonic development. Part of the exstrophy-epispadias complex. Associated with pubic bone diastasis, external genital anomalies, and potential bladder dysfunction. Treatment involves staged surgical repair beginning in the neonatal period, followed by long-term urologic and orthopedic follow-up.

\medterm{Breast Abscess}-- A localized collection of pus within the breast tissue, most commonly a complication of lactational mastitis. Presents as a fluctuant, painful mass with overlying erythema and systemic symptoms. Diagnosis by ultrasound. Treatment requires incision and drainage plus antibiotics. Non-lactational breast abscesses may be associated with smoking or duct ectasia and may require biopsy to rule out underlying malignancy.

\medterm{Childbirth}-- See Parturition (Childbirth).

\medterm{Chromosomes}-- See genetics.

\medterm{Clitoris}-- A small, sensitive erectile organ located at the anterior junction of the labia minora, beneath the clitoral hood (prepuce). The clitoris is the primary source of female sexual pleasure and is homologous to the male penis. It consists of the external glans (visible, about the size of a pea), the body (corpus), and the internal crura (two elongated erectile structures that extend along the pubic rami). The clitoris contains approximately 8,000 nerve endings and becomes engorged with blood during sexual arousal. The only known function of the clitoris is sexual pleasure.

\medterm{Dyspareunia}-- Painful sexual intercourse.

\medterm{Endometrium}-- The lining of the uterus.

\medterm{Erythromycin Ophthalmic Prophylaxis}-- Erythromycin 0.5\% ophthalmic ointment administered to all newborns within 1 hour of birth to prevent ophthalmia neonatorum (neonatal conjunctivitis) from Neisseria gonorrhoeae and Chlamydia trachomatis. Instilled into both conjunctival sacs. Silver nitrate (Credé's method) was used historically but is less effective for chlamydia. The prophylaxis is required by law in most US states. Does not prevent nasopharyngeal colonization.

\medterm{Fertilisation}-- The union of a sperm and an egg (oocyte) to form a zygote, the first cell of a new individual. Fertilization occurs in the ampulla of the fallopian tube and involves: capacitation of sperm, binding of sperm to the zona pellucida, acrosome reaction, fusion of sperm and oocyte plasma membranes, cortical reaction (prevents polyspermy), completion of the second meiotic division of the oocyte, formation of the male and female pronuclei, and syngamy (fusion of pronuclei). The resulting zygote then begins cell division (cleavage) as it travels down the fallopian tube to implant in the uterus.

\medterm{Galactogogue}-- A substance that promotes breast milk production.

\medterm{Gaskin Maneuver (All-Fours Position)}-- A maneuver for shoulder dystocia resolution: the mother is rolled onto her hands and knees (all-fours position). Gravity and the change in pelvic dimensions (increased anteroposterior diameter of the pelvic outlet) may dislodge the impacted shoulder. The posterior arm may also be delivered more easily. Success rate: 25-50\% after failed McRoberts. If the posterior shoulder delivers, the anterior shoulder usually follows.

\medterm{gynecology}-- The medical specialty concerned with the health of the female reproductive system (uterus, ovaries, fallopian tubes, cervix, vagina, and vulva).

\medterm{Haemolytic Disease of the Newborn}-- See blood-related entries in hematology.

\medterm{Hepatitis B Vaccine (Birth Dose)}-- The first dose of the hepatitis B vaccine is recommended within 24 hours of birth for all medically stable infants. Administered IM in the vastus lateralis. Prevents perinatal transmission of HBV. For infants born to HBsAg-positive mothers: HBIG (0.5 mL IM) is given within 12 hours of birth along with the vaccine. The three-dose series (0, 1-2, and 6-18 months) provides >95\% protection. Universal vaccination has dramatically reduced childhood HBV infection.

\medterm{Hysterectomy and Surgical Menopause}-- Hysterectomy is the surgical removal of the uterus, and when the ovaries are also removed (oophorectomy), it induces surgical menopause regardless of age. Surgical menopause causes an abrupt drop in estrogen levels, leading to more severe menopausal symptoms than natural menopause, including hot flashes, bone loss, and increased cardiovascular risk. Hormone replacement therapy is often recommended for younger women to mitigate these effects.

\medterm{Infant Sleep Position}-- The supine sleep position (back to sleep) is recommended by the AAP to reduce the risk of sudden infant death syndrome (SIDS). The Back to Sleep campaign (1994) reduced SIDS rates by >50\%. Always place the baby on the back for sleep (naps and nighttime). Use a firm sleep surface with no soft bedding, pillows, bumpers, or stuffed animals. Room-sharing (same room, separate bed) is recommended for the first 6-12 months. Prone and side positions increase SIDS risk. Tummy time when awake promotes motor development and prevents positional plagiocephaly.

\medterm{Kegel}-- An exercise that helps prevent and treat incontinence by strengthening pelvic floor muscles.

\medterm{Lactation}-- The physiological process of milk production and secretion from the mammary glands, occurring after childbirth (puerperium). See Lactation (Breastfeeding).

\medterm{Lanugo}-- Fine, soft hair that grows all over the body of a fetus and is typically shed before birth.

\medterm{Levator Ani}-- The primary pelvic floor muscle complex supporting the pelvic organs, forming a hammock from the pubic symphysis to the coccyx. Composed of the pubococcygeus, puborectalis, and iliococcygeus muscles. Innervated by the pudendal nerve (S2-S4).

\medterm{Male Infertility}-- Male infertility—the inability to conceive after 12 months of unprotected intercourse—contributes to approximately 50\% of couple infertility cases. Causes include varicocele (most common correctable cause, ~40\%), obstructive azoospermia (congenital bilateral absence of the vas deferens from CFTR mutations, vasectomy, infection), hormonal disorders (hypogonadotropic hypogonadism, hyperprolactinemia), spermatogenic failure (Y-chromosome microdeletions, Klinefelter syndrome, chemotherapy, cryptorchidism), and idiopathic oligoasthenoteratozoospermia. Evaluation: semen analysis per WHO criteria (volume $\ge$1.4 mL, concentration $\ge$16 million/mL, motility $\ge$30\%, morphology $\ge$4\%), hormonal profile (FSH, LH, testosterone), genetic testing (karyotype, Y-deletion, CFTR), and scrotal ultrasound. Treatment: varicocelectomy, hormonal therapy (hCG, clomiphene), surgical sperm extraction (TESE, micro-TESE) for IVF/ICSI, and lifestyle optimization (weight loss, smoking cessation, reducing scrotal heat).

\medterm{McRoberts Maneuver}-- A positioning maneuver for managing shoulder dystocia where the mother’s legs are hyperflexed onto her abdomen, with thighs against the abdomen. This flattens the sacrum, increases the anteroposterior diameter of the pelvis, and rotates the symphysis pubis superiorly. Often combined with suprapubic pressure. First-line maneuver for shoulder dystocia resolution, with a reported success rate of approximately 40\%.

\medterm{Menopause}-- The point marking the end of menstruation, officially designated as one year after a woman’s final period.

\medterm{Newborn Hearing Screen}-- See medical_tests.

\medterm{Oocyte}-- The female gamete (egg cell) produced in the ovary through oogenesis. Oocytes develop from primordial germ cells, which undergo mitosis to form oogonia, then enter meiosis I before birth and arrest in prophase I until puberty. Each menstrual cycle, a cohort of follicles is recruited; one dominant follicle matures and releases its oocyte at ovulation (meiosis I completed, arrested in metaphase II). The oocyte is arrested in metaphase II until fertilization, which triggers completion of meiosis II. The oocyte is the largest cell in the human body (~100--120 \(\mu\)m diameter) and contains maternal mRNA, proteins, organelles, and nutrients (yolk) necessary for early embryonic development until the embryonic genome activates (8-cell stage). Oocyte quality declines with maternal age (increased aneuploidy risk, reduced mitochondrial function). Oocyte cryopreservation (egg freezing) and in vitro fertilization (IVF) are key assisted reproductive technologies.

\medterm{Ovariectomy}-- Surgical removal of one or both ovaries.

\medterm{Ovulation}-- The release of a mature egg from the ovary, at which time it is available to be fertilized by sperm.

\medterm{Perimenopause}-- The transition time in a woman’s life that begins when ovaries produce less estrogen and menstruation becomes less frequent, and ends when the ovaries no longer produce eggs and menstruation stops.

\medterm{Physiologic Jaundice of Newborn}-- A common, benign condition characterized by unconjugated hyperbilirubinemia appearing after 24 hours of age (peak at 3-5 days) in healthy term newborns. Caused by increased bilirubin production (high RBC turnover, short RBC lifespan) and immature hepatic conjugation (low UDP-glucuronyl transferase activity). Total bilirubin typically <12-15 mg/dL. Resolves spontaneously within 1-2 weeks. Management: monitor bilirubin, assess risk factors, and treat with phototherapy if bilirubin exceeds thresholds on the Bhutani nomogram.

\medterm{Postmenopausal osteoporosis}-- Bone loss caused by lower estrogen levels associated with menopause. Sometimes called type I osteoporosis.

\medterm{Postmenopause}-- The period in a woman’s life lasting from the end of perimenopause until the end of life.

\medterm{Premenstrual Syndrome}-- A recurrent pattern of physical, emotional, and behavioral symptoms during the luteal phase of the menstrual cycle that resolve with menses. Affects up to 50\% of reproductive-age women. Symptoms: irritability, mood swings, depression, anxiety, bloating, breast tenderness, fatigue, and appetite changes. Premenstrual dysphoric disorder (PMDD) is a severe form with marked mood symptoms and functional impairment. Treatment: lifestyle modification, SSRIs (continuous or luteal-phase), oral contraceptives, spironolactone, and CBT.

\medterm{Rhesus (Rh) Factor}-- See hematology.

\medterm{Semen}-- The fluid ejaculated from the male reproductive tract, containing sperm (produced in the testes) suspended in seminal plasma (produced by the seminal vesicles, prostate, and bulbourethral glands). Semen provides nutrients (fructose, amino acids), protective agents (zinc, antioxidants), and a medium for sperm transport. Normal semen parameters (WHO 2021): volume \(\ge\)1.4 mL, sperm concentration \(\ge\)16 million/mL, total sperm count \(\ge\)39 million, progressive motility \(\ge\)30\%, normal morphology \(\ge\)4\%.

\medterm{Special Care Baby Unit}-- A ward within a hospital that provides specialist care for newborn babies who need additional support but not intensive care. SCBU provides: monitoring of premature babies (usually born at 32-36 weeks), phototherapy for jaundice, treatment of hypoglycemia, intravenous fluids or nasogastric tube feeding, and treatment of mild respiratory distress. SCBU differs from a Neonatal Intensive Care Unit (NICU), which provides full intensive care including mechanical ventilation for the sickest or most premature infants.

\medterm{Sperm}-- The male gamete (reproductive cell), produced in the seminiferous tubules of the testis through spermatogenesis. Sperm consists of: head (contains the haploid nucleus and acrosome containing enzymes for egg penetration), midpiece (contains mitochondria for energy), and tail (flagellum for motility). Approximately 1,500 sperm are produced per second. Sperm travel from the testis through the epididymis (maturation), vas deferens, seminal vesicles, and urethra during ejaculation.

\medterm{Vulvodynia}-- Pain in the vulva that may or may not be brought on by touch or pressure.

\medterm{Woods Screw Maneuver}-- See Woods Corkscrew Maneuver in the main section above.
